\chapter{使徒行传总结}

\section{圣灵}
\subsection{聖父聖子聖靈关系}

\subsubsection{只有一位神}
\begin{table}[H]
\centering
\begin{tabular}{p{6em}p{41em}}  \toprule
经节&内容 \\  \midrule
申6：4-5&以色列啊，你要聽！耶和華我們神是獨一的主。 5 你要盡心、盡性、盡力愛耶和華你的神。 \\  \midrule
林前8：4&論到吃祭偶像之物，我們知道偶像在世上算不得什麼，也知道神只有一位，再沒有別的神。\\  \midrule
加3：20&但中保本不是為一面做的，神卻是一位。 \\  \midrule
提前2：5&因為只有一位神，在神和人中間只有一位中保，乃是降世為人的基督耶穌，\\  \midrule
\bottomrule
\end{tabular}%
%  \label{tab:addlabel}%
% \caption{法利賽人求神蹟}  
\end{table}

\subsubsection{三位一體包含三個位格——聖父聖子聖靈同时出现}
\begin{table}[H]
\centering
\begin{tabular}{p{5.6em}p{43em}}  \toprule
经节&内容 \\  \midrule

創1:1-3&起初，\textcolor{red}{神}創造天地。地是空虛混沌，淵面黑暗，\textcolor{red}{神的靈}運行在水面上。神說：\textcolor{red}{「要有光。」 }\\  \midrule
創1:26&神說：「我們要照著我們的形象，按著我們的樣式造人，使他們....」 \\  \midrule
創3:22&耶和華神說：「那人已經與我們相似，能知道善惡。現在恐怕...」\\  \midrule
創11:7&我們下去，在那裡變亂他們的口音，使他們的言語彼此不通。\\  \midrule
賽6:8&我又聽見主的聲音說：「我可以差遣誰呢？誰肯為我們去呢？」我說：「我在這裡，請差遣我！」\\  \midrule
賽11:1-5&從耶西的本 a必發一條，從他根生的枝子必結果實。 2 耶和華的靈必住在他身上，就是使他有智慧和聰明的靈，謀略和能力的靈，知識和敬畏耶和華的靈。 3 他必以敬畏耶和華為樂，行審判不憑眼見，斷是非也不憑耳聞； 4 卻要以公義審判貧窮人，以正直判斷世上的謙卑人，以口中的杖擊打世界，以嘴裡的氣殺戮惡人。 5 公義必當他的腰帶，信實必當他脅下的帶子。\\  \midrule

賽42:1-4&看哪，我的僕人，我所扶持、所揀選，心裡所喜悅的，我已將我的靈賜給他，他必將公理傳給外邦。 2 他不喧嚷，不揚聲，也不使街上聽見他的聲音。 3 壓傷的蘆葦他不折斷，將殘的燈火他不吹滅，他憑真實將公理傳開。 4 他不灰心，也不喪膽，直到他在地上設立公理，海島都等候他的訓誨。\\  \midrule
賽48:16&你們要就近我來聽這話，我從起頭並未曾在隱密處說話，自從有這事，我就在那裡。」現在，\textcolor{red}{主耶和華}差遣\textcolor{red}{我}和\textcolor{red}{他的靈}來\\  \midrule
賽61:1-3&\textcolor{red}{主耶和華的靈}在\textcolor{red}{我身}上，因為\textcolor{red}{耶和華}用膏膏我，叫我傳好信息給謙卑的人，差遣我醫好傷心的人，報告被擄的得釋放、被囚的出監牢，報告耶和華的恩年和我們神報仇的日子，安慰...\\  \midrule
太3:16-17&\textcolor{red}{耶穌}受了洗，隨即從水裡上來。天忽然為他開了，他就看見\textcolor{red}{神的靈}彷彿鴿子降下，落在他身上； \textcolor{red}{從天上有聲音}說：「這是我的愛子，我所喜悅的。」\\  \midrule
太28:19-20&所以你們要去，使萬民做我的門徒，奉\textcolor{red}{父、子、聖靈}的名給他們施洗，凡我所吩咐你們的，都教訓他們遵守。我就常與你們同在，直到世界的末了。\\  \midrule
路4:18-19&\textcolor{red}{主的靈}在\textcolor{red}{我身}上，因為\textcolor{red}{他}用膏膏我，叫我傳福音給貧窮的人，差遣我報告被擄的得釋放、瞎眼的得看見，叫那受壓制的得自由，報告神悅納人的禧年。\\  \midrule
約14:16&\textcolor{red}{我}要求\textcolor{red}{父}，父就另外賜給你們一位保惠師，叫他永遠與你們同在，就是\textcolor{red}{真理的聖靈}。\\  \midrule
林後13:14&願\textcolor{red}{主耶穌基督}的恩惠、\textcolor{red}{神}的慈愛、\textcolor{red}{聖靈}的感動，常與你們眾人同在！\\   \midrule
彼前1:2&就是照\textcolor{red}{父神}的先見被揀選，藉著\textcolor{red}{聖靈}得成聖潔，以致順服\textcolor{red}{耶穌基督}，又蒙他血所灑的人。\\ \bottomrule

\end{tabular}%
\end{table}


\subsubsection{聖父聖子聖靈各有區分}
\begin{table}[H]
\centering
\begin{tabular}{p{5.6em}p{43em}}  \toprule
经节&内容 \\  \midrule


&舊約裏，大寫的“主”不同於小寫的“主”\\  \midrule
創19:24&當時，耶和華LORD將硫磺與火從天上耶和華那裡降於所多瑪和蛾摩拉，\\  \midrule
何1:4&耶和華LORD對何西阿說：「給他起名叫耶斯列，因為再過片時，我必討耶戶家在耶斯列殺人流血的罪，也必使以色列家的國滅絕。\\  \midrule

&大寫的“主”有一個“兒子”\\  \midrule
詩2:7-9&受膏者說：「我要傳聖旨。耶和華曾對我說：『你是我的兒子，我今日生你。你求我，我就將列國賜你為基業，將地極賜你為田產。你必用鐵杖打破他們，你必將他們如同窯匠的瓦器摔碎。』」\\  \midrule
詩2:11-12&當存畏懼侍奉耶和華，又當存戰兢而快樂。當以嘴親子，恐怕他發怒，你們便在道中滅亡，因為他的怒氣快要發作。凡投靠他的，都是有福的！\\  \midrule
箴言30:2-4&這人對以鐵和烏甲說：我比眾人更蠢笨，也沒有人的聰明。 3 我沒有學好智慧，也不認識至聖者。 4 誰升天又降下來？誰聚風在掌握中？誰包水在衣服裡？誰立定地的四極？他名叫什麼？他兒子名叫什麼？你知道嗎？\\  \midrule

&聖靈跟大寫的“主”和“神”被區分開來。\\  \midrule
民27:18&耶和華LORD對摩西說：「嫩的兒子約書亞是心中有聖靈的，你將他領來，按手在他頭上，\\  \midrule
詩51:10-12&神God啊，求你為我造清潔的心，使我裡面重新有正直的靈。不要丟棄我，使我離開你的面，不要從我收回你的聖靈。求你使我仍得救恩之樂，賜我樂意的靈扶持我。\\  \midrule

&神子同父神區分開來\\  \midrule
詩45:6-7&神God啊，你的寶座是永永遠遠的，你的國權是正直的。你喜愛公義，恨惡罪惡，所以神，就是你的神，用喜樂油膏你，勝過膏你的同伴。\\  \midrule
來1:8-9&論到子卻說：「神啊，你的寶座是永永遠遠的，你的國權是正直的。你喜愛公義，恨惡罪惡，所以神，就是你的神，用喜樂油膏你，勝過膏你的同伴。」\\  \midrule

%&耶穌要求天父賜下一位保惠師，聖靈。這表明耶穌沒有把自己當作天父或聖靈\\  \midrule
約14:16-20&我要求父，父就另外賜給你們一位保惠師，叫他永遠與你們同在，就是真理的聖靈，乃世人不能接受的，因為不見他，也不認識他。你們卻認識他，因他常與你們同在，也要在你們裡面。我不撇下你們為孤兒，我必到你們這裡來。 還有不多的時候，世人不再看見我，你們卻看見我，因為我活著，你們也要活著。到那日，你們就知道我在父裡面，你們在我裡面，我也在你們裡面。\\  \midrule
\bottomrule
\end{tabular}%
\end{table}

\subsubsection{三位一體中的每一位格都是神}
\begin{table}[H]
\centering
\begin{tabular}{p{5.6em}p{43em}}  \toprule
经节&内容 \\  \midrule

&\textcolor{blue}{天父是神}\\  \midrule
約6:27&不要為那必壞的食物勞力，要為那存到永生的食物勞力，就是人子要賜給你們的，因為人子是父神所印證的。\\  \midrule
羅1:7&我寫信給你們在羅馬、為神所愛、奉召做聖徒的眾人。願恩惠、平安從我們的父神並主耶穌基督歸於你們！\\  \midrule
彼前1:2&就是照父神的先見被揀選，藉著聖靈得成聖潔，以致順服耶穌基督，又蒙他血所灑的人。\\  \midrule

&\textcolor{blue}{聖子是神}\\  \midrule
約1:1-2&太初有道，道與神同在，道就是神。這道太初與神同在。\\  \midrule
約1:14&道成了肉身,住在我們中間,充充滿滿地有恩典有真理.我們也見過他的榮光,正是父獨生子的榮光\\  \midrule
羅9:5&列祖就是他們的祖宗；按肉體說，基督也是從他們出來的——他是在萬有之上永遠可稱頌的神！\\  \midrule
西2:9&因為神本性一切的豐盛，都有形有體地居住在基督裡面； \\  \midrule
來1:8&論到子卻說：「神啊，你的寶座是永永遠遠的，你的國權是正直的。你喜愛公義，恨惡罪惡，所以神，就是你的神，用喜樂油膏你，勝過膏你的同伴。」\\  \midrule
約一5:20&我們也知道神的兒子已經來到，且將智慧賜給我們，使我們認識那位真實的；我們也在那位真實的裡面，就是在他兒子耶穌基督裡面。這是真神，也是永生。\\  \midrule

&\textcolor{blue}{聖靈是神}\\  \midrule
徒5:3-4&彼得說：「亞拿尼亞，為什麼撒旦充滿了你的心，叫你欺哄聖靈，把田地的價銀私自留下幾份呢？田地還沒有賣，不是你自己的嗎？既賣了，價銀不是你做主嗎？你怎麼心裡起這意念呢？你不是欺哄人，是欺哄神了！」\\  \midrule
林前3:16&豈不知你們是神的殿，神的靈住在你們裡頭嗎？ \\  \midrule
羅8:5-11&因為隨從肉體的人，體貼肉體的事；隨從聖靈的人，體貼聖靈的事。體貼肉體的就是死，體貼聖靈的乃是生命、平安。原來體貼肉體的就是與神為仇，因為不服神的律法，也是不能服；而且屬肉體的人不能得神的喜歡。 如果神的靈住在你們心裡，你們就不屬肉體，乃屬聖靈了。人若沒有基督的靈，就不是屬基督的。 基督若在你們心裡，身體就因罪而死，心靈卻因義而活。 然而叫耶穌從死裡復活者的靈若住在你們心裡，那叫基督耶穌從死裡復活的，也必藉著住在你們心裡的聖靈，使你們必死的身體又活過來\\  \midrule
約14:16-17&我要求父，父就另外賜給你們一位保惠師，叫他永遠與你們同在，就是真理的聖靈，乃世人不能接受的，因為不見他，也不認識他。你們卻認識他，因他常與你們同在，也要在你們裡面。\\  \midrule
徒2:1-4&五旬節到了，門徒都聚集在一處。忽然，從天上有響聲下來，好像一陣大風吹過，充滿了他們所坐的屋子； 又有舌頭如火焰顯現出來，分開落在他們各人頭上。他們就都被聖靈充滿，按著聖靈所賜的口才說起別國的話來。\\  \midrule
\bottomrule
\end{tabular}%
\end{table}

\subsubsection{三位一體中的從屬關係}

\begin{table}[H]
\centering
\begin{tabular}{p{5.7em}p{43em}}  \toprule
经节&内容 \\  \midrule

&聖子聽從聖父\\  \midrule
路22:42&說：「父啊！你若願意，就把這杯撤去！然而，不要成就我的意思，只要成就你的意思。」\\  \midrule
約5:36-47&但我有比約翰更大的見證，因為父交給我要我成就的事，就是我所做的事，這便見證我是父所差來的。差我來的父也為我作過見證。你們從來沒有聽見他的聲音，也沒有看見他的形象。你們並沒有他的道存在心裡，因為他所差來的，你們不信。你們查考聖經，因你們以為內中有永生。給我作見證的就是這經，然而你們不肯到我這裡來得生命。  我不受從人來的榮耀。但我知道，你們心裡沒有神的愛。我奉我父的名來，你們並不接待我；若有別人奉自己的名來，你們倒要接待他。你們互相受榮耀，卻不求從獨一之神來的榮耀，怎能信我呢？不要想我在父面前要告你們，有一位告你們的，就是你們所仰賴的摩西。你們如果信摩西，也必信我，因為他書上有指著我寫的話。你們若不信他的書，怎能信我的話呢？\\  \midrule
約20:21-23&耶穌又對他們說：「願你們平安！父怎樣差遣了我，我也照樣差遣你們。」 說了這話，就向他們吹一口氣，說：「你們受聖靈！你們赦免誰的罪，誰的罪就赦免了；你們留下誰的罪，誰的罪就留下了。」\\  \midrule
約一4：14-15&父差子做世人的救主，這是我們所看見且作見證的。凡認耶穌為神兒子的，神就住在他裡面，他也住在神裡面。\\  \midrule

&聖靈聽從聖父和聖子\\  \midrule
約14:1-206&我要求父，父就另外賜給你們一位保惠師，叫他永遠與你們同在，就是真理的聖靈，乃世人不能接受的，因為不見他，也不認識他。你們卻認識他，因他常與你們同在，也要在你們裡面。我不撇下你們為孤兒，我必到你們這裡來。還有不多的時候，世人不再看見我，你們卻看見我，因為我活著，你們也要活著。到那日，你們就知道我在父裡面，你們在我裡面，我也在你們裡面。\\  \midrule
約14:26&但保惠師，就是父因我的名所要差來的聖靈，他要將一切的事指教你們，並且要叫你們想起我對你們所說的一切話。我留下平安給你們，我將我的平安賜給你們。我所賜的不像世人所賜的，你們心裡不要憂愁，也不要膽怯。你們聽見我對你們說了，我去，還要到你們這裡來。你們若愛我，因我到父那裡去，就必喜樂，因為父是比我大的。\\  \midrule
約15:26-27&但我要從父那裡差保惠師來，就是從父出來真理的聖靈；他來了，就要為我作見證。你們也要作見證，因為你從起頭就與我同在。\\  \midrule
約16:7-16&然而我將真情告訴你們，我去是於你們有益的。我若不去，保惠師就不到你們這裡來；我若去，就差他來。他既來了，就要叫世人為罪、為義、為審判，自己責備自己。為罪，是因他們不信我；為義，是因我往父那裡去，你們就不再見我；為審判，是因這世界的王受了審判。我還有好些事要告訴你們，但你們現在擔當不了。 只等真理的聖靈來了，他要引導你們明白一切的真理，因為他不是憑自己說的，乃是把他所聽見的都說出來，並要把將來的事告訴你們。他要榮耀我，因為他要將受於我的告訴你們。凡父所有的，都是我的，所以我說，他要將受於我的告訴你們。等不多時，你們就不得見我；再等不多時，你們還要見我。\\ \bottomrule
\end{tabular}%
\end{table}

\subsubsection{三位一體中聖父聖子的工作}
聖父是以下的源頭和起因
\begin{table}[H]
\centering
\begin{tabular}{p{5.7em}p{6em}p{33em}}  \toprule
经节&内容&经文 \\  \midrule
萬有&林書8：6&然而我們只有一位神，就是父，萬物都本於他，我們也歸於他。並有一位主，就是耶穌基督，萬物都是藉著他有的，我們也是藉著他有的。\\  \midrule
&啟4：11&我們的主，我們的神，你是配得榮耀、尊貴、權柄的，因為你創造了萬物，並且萬物是因你的旨意被創造而有的。\\  \midrule
神性的啟示&啟1：1&耶穌基督的啟示，就是神賜給他，叫他將必要快成的事指示他的眾僕人。他就差遣使者曉諭他的僕人約翰。\\  \midrule
救贖&約3：16-17&神愛世人，甚至將他的獨生子賜給他們，叫一切信他的不致滅亡，反得永生。 17 因為神差他的兒子降世，不是要定世人的罪，乃是要叫世人因他得救。\\  \midrule
耶穌做的事&約5：17&耶穌就對他們說：「我父做事直到如今，我也做事。」\\  \midrule
&約14：10&我在父裡面，父在我裡面，你不信嗎？我對你們所說的話，不是憑著自己說的，乃是住在我裡面的父做他自己的事。\\  \midrule
\bottomrule
\end{tabular}\caption{聖父是以下的源頭和起因} 
\end{table}



聖子是聖父的代表做了以下的工作
\begin{table}[H]
\centering
\begin{tabular}{p{5.7em}p{6em}p{33em}}  \toprule
经节&内容&经文 \\  \midrule
創造和維持萬有&林前8：6&然而我們只有一位神，就是父，萬物都本於他，我們也歸於他。並有一位主，就是耶穌基督，萬物都是藉著他有的，我們也是藉著他有的。\\  \midrule
&約1：3&萬物是藉著他造的；凡被造的，沒有一樣不是藉著他造的。\\  \midrule
&西1：16-17&因為萬有都是靠他造的，無論是天上的、地上的，能看見的、不能看見的，或是有位的、主治的、執政的、掌權的，一概都是藉著他造的，又是為他造的。他在萬有之先，萬有也靠他而立。\\  \midrule
神性的啟示&約1：1&太初有道，道與神同在，道就是神。\\  \midrule
&太11：27&一切所有的，都是我父交付我的。除了父，沒有人知道子；除了子和子所願意指示的，沒有人知道父。\\  \midrule
&約16：12-15&我還有好些事要告訴你們，但你們現在擔當不了。 只等真理的聖靈來了，他要引導你們明白一切的真理，因為他不是憑自己說的，乃是把他所聽見的都說出來，並要把將來的事告訴你們。他要榮耀我，因為他要將受於我的告訴你們。凡父所有的，都是我的，所以我說，他要將受於我的告訴你們。\\  \midrule
&啟1：1&耶穌基督的啟示，就是神賜給他，叫他將必要快成的事指示他的眾僕人。他就差遣使者曉諭他的僕人約翰。\\  \midrule
救贖&林後5：19&這就是神在基督裡叫世人與自己和好，不將他們的過犯歸到他們身上，並且將這和好的道理託付了我們。\\  \midrule
&太1：21&她將要生一個兒子，你要給他起名叫耶穌，因他要將自己的百姓從罪惡裡救出來。\\  \midrule
&約4：42&便對婦人說：「現在我們信，不是因為你的話，是我們親自聽見了，知道這真是救世主。」\\  \midrule
\bottomrule
\end{tabular}\caption{聖子是聖父的代表做了以下的工作} 
\end{table}










\subsection{圣灵名词解释}


\subsubsection{圣灵施洗}




\subsubsubsection{圣灵施洗是使教会成为基督的“身体”}
\begin{table}[H]
  \centering
  \begin{tabular}{p{36em}|p{13em}}  \toprule
成为基督的“身体” - 林前12:4-20 & “身体” 的定义-  約 1:3/3:5-8 \\  \midrule
 \multirow{2}{36.2em}{恩賜原有分別,聖靈卻是一位；職事也有分別,主卻是一位； 功用也有分別,神卻是一位,在眾人裡面運行一切的事。聖靈顯在各人身上,是叫人得益處。這人蒙聖靈賜他智慧的言語,那人也蒙這位聖靈賜他知識的言語,又有一人蒙這位聖靈賜他信心,還有一人蒙這位聖靈賜他醫病的恩賜,又叫一人能行異能,又叫一人能做先知,又叫一人能辨別諸靈,又叫一人能說方言,又叫一人能翻方言。這一切都是這位聖靈所運行,隨己意分給各人的。就如身子是一個,卻有許多肢體,而且肢體雖多,仍是一個身子；基督也是這樣。我們不拘是猶太人，是希臘人，是為奴的，是自主的，都從一位聖靈受洗，成了一個身體,飲於一位聖靈。身子原不是一個肢體,乃是許多肢體 設若腳說：「我不是手，所以不屬乎身子」,它不能因此就不屬乎身子。設若耳說：「我不是眼，所以不屬乎身子」，它也不能因此就不屬乎身子。若全身是眼,從哪裡聽聲呢？若全身是耳,從哪裡聞味呢?但如今，神隨自己的意思把肢體俱各安排在身上了。若都是一個肢體,身子在哪裡呢?但如今肢體是多的，身子卻是一個。}& 這等人不是從血氣生的，不是從情慾生的,也不是從人意生的，乃是從神生的。\\\cline{2-2}
&耶穌說：「我實實在在地告訴你：人若不是從水和聖靈生的，就不能進神的國。從肉身生的就是肉身，從靈生的就是靈。我說你們必須重生，你不要以為稀奇。風隨著意思吹，你聽見風的響聲，卻不曉得從哪裡來，往哪裡去。凡從聖靈生的，也是如此。」\\ \bottomrule
\end{tabular}
\end{table}


\subsubsubsection{只有死里复活的耶稣是用\textcolor{red}{圣灵施洗}的那一位}
\begin{table}[H]
\centering
\begin{tabular}{p{4em}|p{44em}}  \toprule
经文&内容 \\  \midrule

太3:11-12&我是用水給你們施洗,叫你們悔改;但那在我以後來的,能力比我更大,我就是給他提鞋也不配,他要用聖靈與火給你們施洗。他手裡拿著簸箕,要揚淨他的場,把麥子收在倉裡,把糠用不滅的火燒盡了。\\  \midrule
可1:8&我是用水給你們施洗，他卻要用聖靈給你們施洗。 \\  \midrule
路3:16&約翰說：「我是用水給你們施洗。但有一位能力比我更大的要來，我就是給他解鞋帶也不配，他要用聖靈與火給你們施洗。 \\  \midrule
約1:26-27,32-33&約翰回答說：「我是用水施洗，但有一位站在你們中間，是你們不認識的，就是那在我以後來的，我給他解鞋帶也不配。」
約翰又作見證說：「我曾看見聖靈彷彿鴿子從天降下，住在他的身上。我先前不認識他，只是那差我來用水施洗的對我說：『你看見聖靈降下來，住在誰的身上，誰就是用聖靈施洗的。』 \\  \midrule

%賽11:1-5&從耶西的本必發一條，從他根生的枝子必結果實。 2 耶和華的靈必住在他身上，就是使他有智慧和聰明的靈，謀略和能力的靈，知識和敬畏耶和華的靈。 3 他必以敬畏耶和華為樂，行審判不憑眼見，斷是非也不憑耳聞； 4 卻要以公義審判貧窮人，以正直判斷世上的謙卑人，以口中的杖擊打世界，以嘴裡的氣殺戮惡人。 5 公義必當他的腰帶，信實必當他脅下的帶子。 \\  \midrule


%賽42:1-4&看哪，我的僕人，我所扶持、所揀選，心裡所喜悅的，我已將我的靈賜給他，他必將公理傳給外邦。 2 他不喧嚷，不揚聲，也不使街上聽見他的聲音。 3 壓傷的蘆葦他不折斷，將殘的燈火他不吹滅，他憑真實將公理傳開。 4 他不灰心，也不喪膽，直到他在地上設立公理，海島都等候他的訓誨。 \\  \midrule

%賽61:1-3&主耶和華的靈在我身上，因為耶和華用膏膏我，叫我傳好信息給謙卑的人 a，差遣我醫好傷心的人，報告被擄的得釋放、被囚的出監牢， 2 報告耶和華的恩年和我們神報仇的日子，安慰一切悲哀的人， 3 賜華冠於錫安悲哀的人代替灰塵，喜樂油代替悲哀，讚美衣代替憂傷之靈，使他們稱為公義樹，是耶和華所栽的，叫他得榮耀。\\  \midrule

路13:13&你們雖然不好，尚且知道拿好東西給兒女，何況天父，豈不更將聖靈給求他的人嗎？ \\  \midrule

約4:10-14&耶穌回答說：「你若知道神的恩賜和對你說『給我水喝』的是誰，你必早求他，他也必早給了你活水。」 11 婦人說：「先生，沒有打水的器具，井又深，你從哪裡得活水呢？ 12 我們的祖宗雅各將這井留給我們，他自己和兒子並牲畜也都喝這井裡的水，難道你比他還大嗎？」 耶穌回答說：「凡喝這水的，還要再渴， 14 人若喝我所賜的水，就永遠不渴。我所賜的水要在他裡頭成為泉源，直湧到永生。」 \\  \midrule

約7:37-39&節期的末日，就是最大之日，耶穌站著高聲說：「人若渴了，可以到我這裡來喝！ 38 信我的人就如經上所說，從他腹中要流出活水的江河來。」 39 耶穌這話是指著信他之人要受聖靈說的。那時還沒有賜下聖靈來，因為耶穌尚未得著榮耀。\\  \midrule

徒1:4-5&耶穌和他們聚集的時候，囑咐他們說：「不要離開耶路撒冷，要等候父所應許的，就是你們聽見我說過的。约翰是用水施洗。但不多几日，你们要受圣灵的洗。\\ \bottomrule
%亚12:10&我必将那施恩叫人\textcolor{red}{恳求的灵}, 浇灌大卫家和耶路撒冷的居民, 他们必仰望我, 就是他们所扎的\\  \midrule
\end{tabular}%
\end{table}







\subsubsection{聖靈充满（动词,希腊文被圣灵管制的意思）和满有聖靈（形容词）}
\begin{table}[H]
  \centering
  \begin{tabular}{p{5em}|p{5em}|p{38em}}  \toprule
经文 & 人 & \\  \midrule
路1:15&施洗约翰&他在主面前將要為大，淡酒濃酒都不喝，從母腹裡\textcolor{red}{就被聖靈充滿了 shall be filled with the Holy}。\\  \midrule
路1:41-42&伊利莎白&伊利莎白一聽馬利亞問安，所懷的胎就在腹裡跳動，伊利莎白且\textcolor{red}{被聖靈充滿 was filled with the Holy}，高聲喊著說：\\  \midrule
路1:67&撒迦利亞&他父親撒迦利亞\textcolor{red}{被聖靈充滿了was filled with the Holy}，就預言說：\\  \midrule
徒2:4,14-42&門徒&他们就都\textcolor{red}{被圣灵充满 filled with the Holy}，按着圣灵所赐的口才，说起别国的话来。...彼得和十一個使徒站起，高聲說：...於是，領受他話的人就受了洗，那一天，門徒約添了三千人； 都恆心遵守使徒的教訓，彼此交接、掰餅、祈禱。\\  \midrule
%本节所说的圣灵充满，也就是一章五节和八节所说的“圣灵的洗”和“圣灵降临”。重生得救的信徒在“受圣灵”之后，可能一再被圣灵充满，
徒4:5-22&彼得&殘疾人得了痊癒,官府、長老和文士在耶路撒冷聚會，又有大祭司亞那和該亞法、約翰、亞歷山大，並大祭司的親族都在那裡，叫使徒站在當中，就問他們說：「你們用什麼能力，奉誰的名做這事呢？」那時彼得\textcolor{red}{被聖靈充滿 filled with the Holy}，對他們說：.....\\  \midrule
徒4:31&會友&禱告完了，聚會的地方震動，他們就都\textcolor{red}{被聖靈充滿 filled with the Holy}，放膽講論神的道。\\  \midrule
徒9:17&保罗&亞拿尼亞就去了，進入那家，把手按在掃羅身上，說：「兄弟掃羅，在你來的路上向你顯現的主，就是耶穌，打發我來叫你能看見，又\textcolor{red}{被聖靈充滿 be filled with the Holy}。」\\  \midrule
徒13:9&保羅&掃羅又名保羅，\textcolor{red}{被聖靈充滿 filled with the Holy}，定睛看他（行邪术之巴耶稣）說\\  \midrule
徒13:52&門徒&門徒滿心喜樂，又\textcolor{red}{被聖靈充滿 were filled with the Spirit Holy}。\\  \midrule
弗5:18&所有信徒&不要醉酒，酒能使人放蕩；乃要\textcolor{red}{被聖靈充滿 be filled with the Spirit}。\\  \midrule

路4:1&耶穌&耶穌被\textcolor{blue}{聖靈充滿 being full of the Holy Spirit}，從約旦河回來，聖靈將他引到曠野，四十天受魔鬼的試探。\\  \midrule
徒6:3-4&管理飯食&所以弟兄們，當從你們中間選出七個有好名聲、\textcolor{blue}{被聖靈充滿 full of the Spirit}、智慧充足的人，我們就派他們管理這事。但我們要專心以祈禱、傳道為事。\\  \midrule
徒6:3-5&司提反&大眾都喜悅這話，就揀選了司提反，乃是大有信心、\textcolor{blue}{被聖靈充滿 full of the Holy}的人；又揀選腓利、伯羅哥羅、尼迦挪、提門、巴米拿，並進猶太教的安提阿人尼哥拉。\\  \midrule
徒7:55-56&司提反&但司提反被\textcolor{blue}{聖靈充滿being full of the Holy Spirit}，定睛望天，看見神的榮耀，又看見耶穌站在神的右邊，就說：「我看見天開了，人子站在神的右邊！」\\  \midrule
徒11:24&巴拿巴&這巴拿巴原是個好人，被\textcolor{blue}{聖靈充滿 full of the Holy Spirit}，大有信心。於是，有許多人歸服了主。\\  \midrule
加5:22-26&所有信徒,滿有聖靈的定义&聖靈所結的果子，就是仁愛、喜樂、和平、忍耐、恩慈、良善、信實、聖靈所結的果子，就是仁愛、喜樂、和平、忍耐、恩慈、良善、信實、溫柔、節制；這樣的事沒有律法禁止。凡屬基督耶穌的人，是已經把肉體連肉體的邪情私慾同釘在十字架上了。我們若是靠聖靈得生，就當靠聖靈行事。不要貪圖虛名，彼此惹氣，互相嫉妒。\\  \midrule

\bottomrule
\end{tabular}
\end{table}




\subsubsection{圣灵浇灌(倾倒,用膏油“倒”在蒙拣选的人身上)}
\begin{table}[H]
\centering
\begin{tabular}{p{4em}|p{45em}}  \toprule
经文&内容\\  \midrule
利8:6-30&摩西帶了亞倫和他兒子來，用\textcolor{red}{水}洗了他們，給亞倫穿上內袍，....。摩西用\textcolor{red}{膏油}抹帳幕和其中所有的，使它成聖，又用\textcolor{red}{膏油}在壇上彈了七次，又抹了壇和壇的一切器皿，並洗濯盆和盆座，使它成聖，又把\textcolor{red}{膏油}倒在亞倫的頭上膏他，使他成聖。摩西帶了亞倫的兒子來，給他們穿上內袍，束上腰帶，包上裹頭巾，都是照耶和華所吩咐摩西的。

他牽了贖罪祭的公牛來，亞倫和他兒子按手在贖罪祭公牛的頭上，就宰了公牛。摩西用指頭蘸\textcolor{red}{血}，抹在壇上四角的周圍，使壇潔淨，把\textcolor{red}{血}倒在壇的腳那裡，使壇成聖，壇就潔淨了，又取臟上所有的脂油和肝上的網子，並兩個腰子與腰子上的脂油，都\textcolor{red}{燒}在壇上，唯有公牛，連皮帶肉並糞，用火\textcolor{red}{燒}在營外，都是照耶和華所吩咐摩西的。

他奉上燔祭的公綿羊，亞倫和他兒子按手在羊的頭上，就宰了公羊。摩西把\textcolor{red}{血}灑在壇的周圍，把羊切成塊子，把頭和肉塊並脂油都燒了，用水洗了臟腑和腿，就把全羊燒在壇上為馨香的燔祭，是獻給耶和華的火祭，都是照耶和華所吩咐摩西的。

他又奉上第二隻公綿羊，就是承接聖職之禮的羊，亞倫和他兒子按手在羊的頭上，就宰了羊。摩西把些\textcolor{red}{血}抹在亞倫的右耳垂上和右手的大拇指上，並右腳的大拇指上；又帶了亞倫的兒子來，把些\textcolor{red}{血}抹在他們的右耳垂上和右手的大拇指上，並右腳的大拇指上；又把血灑在壇的周圍。取脂油和肥尾巴，並臟上一切的脂油與肝上的網子，兩個腰子和腰子上的脂油，並右腿；

再從耶和華面前盛無酵餅的筐子裡取出一個無酵餅、一個油餅、一個薄餅，都放在脂油和右腿上；把這一切放在亞倫的手上和他兒子的手上做搖祭，在耶和華面前搖一搖。摩西從他們的手上拿下來，燒在壇上的燔祭上，都是為承接聖職獻給耶和華馨香的火祭。

摩西拿羊的胸作為搖祭，在耶和華面前搖一搖，是承接聖職之禮，歸摩西的份，都是照耶和華所吩咐摩西的。摩西取點\textcolor{red}{膏油}和壇上的\textcolor{red}{血}，彈在亞倫和他的衣服上，並他兒子和他兒子的衣服上，使他和他們的衣服一同成聖。\\  \midrule

撒上10:1&撒母耳拿瓶膏油倒在掃羅的頭上，與他親嘴，說：「這不是耶和華膏你做他產業的君嗎？\\  \midrule

王上19:15-16&耶和華對他說：「你回去，從曠野往大馬士革去。到了那裡，就要膏哈薛做亞蘭王， 16 又膏寧示的孫子耶戶做以色列王，並膏亞伯米何拉人沙法的兒子以利沙做先知接續你。\\  \midrule

箴1:23&你們當因我的責備回轉，我要將我的靈\textcolor{red}{澆灌}你們，將我的話指示你們。\\  \midrule
賽32:15&等到聖靈從上\textcolor{red}{澆灌}我們，曠野就變為肥田，肥田看如樹林。\\  \midrule
賽44:3&因為我要將水\textcolor{red}{澆灌}口渴的人，將河澆灌乾旱之地，我要將我的靈澆灌你的後裔，將我的福澆灌你的子孫。\\  \midrule
珥2:28&以後我要將我的靈\textcolor{red}{澆灌}凡有血氣的,你們的兒女要說預言,你們的老年人要做異夢,少年人要見異象\\  \midrule
徒2:17-21&這正是先知約珥所說的： 17 『神說：在末後的日子，我要將我的靈澆灌凡有血氣的，你們的兒女要說預言，你們的少年人要見異象，老年人要做異夢。 18 在那些日子，我要將我的靈\textcolor{red}{澆灌}我的僕人和使女，他們就要說預言。 19 在天上我要顯出奇事，在地下我要顯出神蹟，有血，有火，有煙霧。 20 日頭要變為黑暗，月亮要變為血，這都在主大而明顯的日子未到以前。 21 到那時候，凡求告主名的，就必得救。』\\  \midrule
徒2:33&他既被神的右手高举，又从父受了所应许的聖靈，就把你们所看见所听见的，\textcolor{red}{澆灌}下来。\\  \midrule
徒10:45&那些奉割禮、和彼得同來的信徒，見聖靈的恩賜也\textcolor{red}{澆}在外邦人身上，就都稀奇，\\  \midrule
羅5:1-4&我們既因信稱義，就藉著我們的主耶穌基督得與神相和。我們又藉著他，因信得進入現在所站的這恩典中，並且歡歡喜喜盼望神的榮耀。不但如此，就是在患難中也是歡歡喜喜的。因為知道患難生忍耐忍耐生老練，老練生盼望，盼望不至於羞恥，因為所賜給我們的聖靈將神的愛\textcolor{red}{澆灌}在我們心裡。\\  \midrule
提3:4-7&但到了神我們救主的恩慈和他向人所施的慈愛顯明的時候,他便救了我們,並不是因我們自己所行的義，乃是照他的憐憫，藉著重生的洗和聖靈的更新。聖靈就是神藉著耶穌基督——我們救主厚厚\textcolor{red}{澆灌}在我們身上的，好叫我們因他的恩得稱為義，可以憑著永生的盼望成為後嗣(可以憑著盼望承受永生)。\\  \midrule
\end{tabular}
\end{table}

%\subsubsection{圣灵的洗，圣灵降临，圣灵充满，圣灵浇灌}
\subsubsection{领受圣灵，圣灵的洗，圣灵降临-旧约圣徒进入新约（国度）时期}
“受圣灵”常用于信主得救时领受圣灵的经历。这经历又被看为是圣灵降在人身上或圣灵的洗。如：彼得把“圣灵降在…人身上”跟“受圣灵”看作同一回事。稍后，彼得回到耶路撒冷，向使徒与众弟兄报告到哥尼流家传道经过时，却把他们受圣灵的经历看作受圣灵的洗。他在徒11:15节所说的“圣灵降在他们身上”，也就是徒11:16末句的“圣灵的洗”；这圣灵的洗就是徒11:18节所说那些外邦人悔改得生命时“受圣灵”的经历（参徒10：47）。
\begin{itemize}
\item 领受所赐的圣灵在我們心裡做憑據(質)-徒2:38,林后1:21-22
\item 身份上是兒子，靠著神為後嗣-加4:6-7
\end{itemize}

\begin{table}[H]
\centering
 \begin{tabular}{p{6em}|p{40em}}  \toprule
经文&内容\\  \midrule
結36:25-27&我必用清水灑在你們身上，你們就潔淨了。我要潔淨你們，使你們脫離一切的汙穢，棄掉一切的偶像。…\\  \midrule
徒1:5&約翰是用水施洗，但不多幾日，你們要\textcolor{red}{受聖靈的洗}。\\  \midrule
徒2:38&彼得说：你们各人要悔改，奉耶稣基督的名受洗，叫你们的罪得赦，就必\textcolor{red}{领受所赐的圣灵}。\\  \midrule

徒19:4-6&保羅說：「約翰所行的是悔改的洗，告訴百姓當信那在他以後要來的，就是耶穌。」…\\  \midrule
林前12:13&我們不拘是猶太人，是希臘人，是為奴的，是自主的，都\textcolor{red}{從一位聖靈受洗}，成了一個身體，飲於一位聖靈。\\  \midrule
林后1:21-22&那在基督裡堅固我們和你們，並且膏我們的就是神； 22 他又用印印了我們，並賜聖靈在我們心裡做憑據(質)。\\  \midrule
多3:5,6&他便救了我們，並不是因我們自己所行的義，乃是照他的憐憫，藉著重生的洗和聖靈的更新。…\\  \midrule
弗1:13-14&你們既聽見真理的道，就是那叫你們得救的福音，也信了基督，既然信他，就受了所應許的聖靈為印記。這聖靈是我們得基業的憑據，直等到神之民被贖，使他的榮耀得著稱讚。\\  \midrule
加4:6-7&你們既為兒子，神就差他兒子的靈進入你們的心，呼叫：「阿爸！父！」可見，從此以後你不是奴僕，乃是兒子了；既是兒子，就靠著神為後嗣。\\  \midrule

徒2:1-4，耶路撒冷&五旬節到了，門徒都聚集在一處。 2 忽然，從天上有響聲下來，好像一陣大風吹過，充滿了他們所坐的屋子； 3 又有舌頭如火焰顯現出來，分開落在他們各人頭上。 4 他們就都被聖靈充滿，按著聖靈所賜的口才說起別國的話來。\\  \midrule
徒8:14-17，撒玛利亚&使徒在耶路撒冷聽見撒馬利亞人領受了神的道，就打發彼得、約翰往他們那裡去。 15 兩個人到了，就為他們禱告，要叫他們\textcolor{red}{受聖靈}， 16 因為聖靈還沒有降在他們一個人身上，他們只奉主耶穌的名受了洗。 17 於是使徒按手在他們頭上，他們就\textcolor{red}{受了聖靈}。\\  \midrule
徒10:44-48，该撒利亚，哥尼流家&彼得還說這話的時候，\textcolor{red}{聖靈降在}一切聽道的人身上。 45 那些奉割禮、和彼得同來的信徒，見\textcolor{red}{聖靈的恩賜也澆}在外邦人身上，就都稀奇， 46 因聽見他們說方言，稱讚神為大。 47 於是彼得說：「這些人既\textcolor{red}{受了聖靈}，與我們一樣，誰能禁止用水給他們施洗呢？」 48 就吩咐奉耶穌基督的名給他們施洗。他們又請彼得住了幾天。\\  \midrule
徒11:15-18，该撒利亚，哥尼流家&我一開講，聖靈便\textcolor{red}{降}在他們身上，正像當初降在我們身上一樣。 16 我就想起主的話說：『約翰是用水施洗，但你們要\textcolor{red}{受聖靈的洗}。』 17 神既然給他們恩賜，像在我們信主耶穌基督的時候給了我們一樣，我是誰，能攔阻神呢？」 18 眾人聽見這話，就不言語了，只歸榮耀於神，說：「這樣看來，神也賜恩給外邦人，叫他們悔改得生命了。」\\  \midrule
徒19:1-7，以弗所&亞波羅在哥林多的時候，保羅經過了上邊一帶地方，就來到以弗所。在那裡遇見幾個門徒， 2 問他們說：「你們信的時候受了聖靈沒有？」他們回答說：「沒有，也未曾聽見有聖靈賜下來。」 3 保羅說：「這樣，你們受的是什麼洗呢？」他們說：「是約翰的洗。」 4 保羅說：「約翰所行的是悔改的洗，告訴百姓當信那在他以後要來的，就是耶穌。」 5 他們聽見這話，就奉主耶穌的名受洗。 6 保羅按手在他們頭上，聖靈便降在他們身上，他們就說方言，又說預言。 7 一共約有十二個人。\\  \midrule
\end{tabular}
\end{table}

\subsubsection{圣灵的别名-兒子的靈，}
\begin{table}[H]
  \centering
  \begin{tabular}{p{6em}|p{42em}}  \toprule
经文&内容\\  \midrule
%荣耀的灵(The Spirit of Glory and of God, 彼前4:14)、

%弗4:4&身体只有一个, 圣灵只有\textcolor{red}{一个(KJV: one Spirit)}, 正如你们蒙召, 同有一个指望.\\  \midrule
启1:4&约翰写信给亚西亚的七个教会. 但愿从那昔在今在以後永在的神, 和他宝座前的\textcolor{red}{七靈}.\\  \midrule
启3:1&你要寫信給撒狄教會的使者說：『那有神的\textcolor{red}{七靈}和七星的說：我知道你的行為，按名你是活的，其實是死的。\\  \midrule
启4:5&有閃電、聲音、雷轟從寶座中發出。又有七盞火燈在寶座前點著，這七燈就是神的\textcolor{red}{七靈}。\\  \midrule
启5:6&我又看見寶座與四活物並長老之中有羔羊站立，像是被殺過的，有七角七眼，就是神的\textcolor{red}{七靈}，奉差遣往普天下去的。\\  \midrule

路4:18&\textcolor{red}{主的靈}在我身上, 因为他用膏膏我, 叫我传福音给贫穷的人….”\\  \midrule
 
徒5:9&彼得說：「你們為什麼同心試探\textcolor{red}{主的靈}呢？埋葬你丈夫之人的腳已到門口\\  \midrule
徒8:39&從水裡上來，\textcolor{red}{主的靈}把腓利提了去。太監也不再見他了，\\  \midrule
林后3:17&主就是那靈，\textcolor{red}{主的靈}在哪裡，哪裡就得以自由。 18 我們眾人既然敞著臉得以看見主的榮光，好像從鏡子裡返照，就變成主的形狀，榮上加榮，如同從\textcolor{red}{主的靈}變成的。\\  \midrule


约16:13-15&只等\textcolor{red}{真理的聖靈}來了，他要引導你們明白一切的真理，因為他不是憑自己說的，乃是把他所聽見的都說出來，並要把將來的事告訴你們。 14 他要榮耀我，因為他要將受於我的告訴你們。 15 凡父所有的，都是我的，所以我說，他要將受於我的告訴你們。\\  \midrule
约14:16-17&我要求父, 父就另外赐给你们一位\textcolor{red}{保惠师(或作“训慰师”)}, 叫他永远与你们同在,就是\textcolor{red}{真理的圣灵}, 乃世人不能接受的. 因为不见他, 也不认识他. 你们却认识他, 因他常与你们同在, 也要在你们里面.\\  \midrule


约14:26&但\textcolor{red}{保惠師}，就是父因我的名所要差來的聖靈，他要將一切的事指教你們，並且要叫你們想起我對你們所說的一切話。\\  \midrule
约15:26&但我要從父那裡差\textcolor{red}{保惠師}來，就是從父出來真理的聖靈；他來了，就要為我作見證。\\  \midrule
约16:7-11&我去是於你們有益的。我若不去，\textcolor{red}{保惠師}就不到你們這裡來；我若去，就差他來。 8 他既來了，就要叫世人為罪、為義、為審判，自己責備自己。 9 為罪，是因他們不信我； 10 為義，是因我往父那裡去，你們就不再見我； 11 為審判，是因這世界的王受了審判。\\  \midrule

林后3:3&你们明显是基督的信(KJV: epistle), 借着我们修成的. 不是用墨写的, 乃是用\textcolor{red}{永生神的灵}写的. 不是写在石版上, 乃是写在心版上.\\  \midrule
林前6:11&你们中间也有人从前是这样. 但如今你们奉主耶稣基督的名, 并借着我们\textcolor{red}{神的灵}, 已经洗净, 成圣称义了.\\  \midrule

创1:2&地是空虚混沌, 渊面黑暗. \textcolor{red}{神的灵}运行在水面上.\\  \midrule
太3:16&耶穌受了洗，隨即從水裡上來。天忽然為他開了，他就看見\textcolor{red}{神的灵}彷彿鴿子降下，落在他身上；從天上有聲音說：「這是我的愛子，我所喜悅的。」\\  \midrule
林前6:11&你们中间也有人从前是这样. 但如今你们奉主耶稣基督的名, 并借着我们\textcolor{red}{神的灵}, 已经洗净, 成圣称义了.\\  \midrule
林前2:11-12&除了在人裡頭的靈，誰知道人的事？像這樣，除了\textcolor{red}{神的靈}，也沒有人知道神的事。 12 我們所領受的，並不是世上的靈，乃是從神來的靈，叫我們能知道神開恩賜給我們的事。\\  \midrule


罗8:9&如果\textcolor{red}{神的灵}住在你们心里,你们就不属肉体,乃属圣灵了. 人若没有\textcolor{red}{基督的灵},就不是属基督的.\\  \midrule 
彼前1:8-12&如今雖不得看見，卻因信他就有說不出來、滿有榮光的大喜樂，並且得著你們信心的果效，就是靈魂的救恩。論到這救恩，那預先說你們要得恩典的眾先知早已詳細地尋求考察，就是考察在他們心裡\textcolor{red}{基督的灵}，預先證明基督受苦難、後來得榮耀是指著什麼時候，並怎樣的時候。\\  \midrule
腓1:19&因为我知道这事借着你们的祈祷, 和\textcolor{red}{耶稣基督的灵}的帮助, 终必叫我得救.\\  \midrule
徒16:7&到了每西亚的边界, 他们想要往庇推尼去, \textcolor{red}{耶稣的灵}却不许.\\  \midrule

\end{tabular}
\end{table}


\begin{table}[H]
\centering
\begin{tabular}{p{5em}|p{44em}}  \toprule
经文&内容\\  \midrule

民11:29&摩西对他说:你为我的缘故嫉妒人吗?惟愿耶和华的百姓都受感说话,愿耶和华把\textcolor{red}{他的靈}降在他们身上\\  \midrule

弗3:15-19&天上地上的各家，都是從他得名—— 求他按著他豐盛的榮耀，藉著\textcolor{red}{他的靈}，叫你們心裡的力量剛強起來， 使基督因你們的信住在你們心裡，叫你們的愛心有根有基，能以和眾聖徒一同明白基督的愛是何等長闊高深， 並知道這愛是過於人所能測度的，便叫神一切所充滿的充滿了你們。\\  \midrule

罗8:15-16&你们所受的不是奴仆的心, 仍旧害怕, 所受的乃是\textcolor{red}{儿子的心(原文是“养子/义子的灵”)}, 因此我们呼叫阿爸、父.聖靈與我們的心同證我們是神的兒女。 \\  \midrule

加4:6-7&你們既為兒子，神就差他\textcolor{red}{兒子的靈}進入你們的心，呼叫：「阿爸！父！」可見，從此以後你不是奴僕，乃是兒子了；既是兒子，就靠著神為後嗣。\\  \midrule

士3:10&\textcolor{red}{耶和华的靈}降在他身上, 他就作了以色列的士师, 出去争战….\\  \midrule
赛4:3&主以\textcolor{red}{公義的靈}和\textcolor{red}{焚燒的靈}，將錫安女子的汙穢洗去，又將耶路撒冷中殺人的血除淨。那時，剩在錫安留在耶路撒冷的，就是一切住耶路撒冷在生命冊上記名的，必稱為聖。\\  \midrule
赛11:2&\textcolor{red}{耶和华的灵}必住在他身上,就是使他有\textcolor{red}{智慧和聪明的灵},\textcolor{red}{谋略和能力的灵},\textcolor{red}{知识和敬畏耶和华的灵}.\\  \midrule

太10:20&因为不是你们自己说的, 乃是你们\textcolor{red}{父的灵}在你们里头说的.\\  \midrule

创6:3&耶和华说: 人既属乎血气, \textcolor{red}{我的灵}就不永远住在他里面, 然而他的日子还可到一百二十年.\\  \midrule
诗139:7&我往那里去躲避\textcolor{red}{你的灵}, 我往那里逃躲避你的面.\\  \midrule

啟11:11&過了這三天半，有生氣(生命圣灵the Spirit of life from God)從神那裡進入他們裡面，他們就站起來，看見他們的人甚是害怕。\\  \midrule

罗8:2&因为\textcolor{red}{赐生命圣灵(The Spirit of Life)}的律, 在基督耶稣里释放了我, 使我脱离罪和死的律了.\\  \midrule

罗1:4&按\textcolor{red}{圣善(洁)的灵 (The Spirit of Holiness)}说, 因从死里复活, 以大能显明是神的儿子.\\  \midrule
来9:14&何况基督借着\textcolor{red}{永远的灵}, 将自己无瑕无疵献给神, 他的血岂不更能洗净你们的心, 除去你们的死行, 使你们事奉那永生神吗!\\  \midrule
彼前4:14&你们若为基督的名受辱骂, 便是有福的. 因为\textcolor{red}{神荣耀的灵}, 常住在你们身上.\\  \midrule

罗8:11&然而叫耶稣从\textcolor{red}{死里复活者的灵(原文: 他的灵; KJV: the Spirit of Him that raised up Jesus from the dead…)}, 若住在你们心里, 那叫基督耶稣从死里复活的, 也必借着住在你们心里的圣灵, 使你们必死的身体又活过来.\\  \midrule

亚12:10&我必将那施恩叫人\textcolor{red}{恳求的灵}, 浇灌大卫家和耶路撒冷的居民, 他们必仰望我, 就是他们所扎的\\  \midrule
诗51:12&求你使我仍得救恩之乐, 赐我\textcolor{red}{乐意的灵}扶持我.\\  \midrule
\end{tabular}
\end{table}














\begin{table}[H]
\centering
\begin{tabular}{p{10em}|p{18em}|p{13em}}  \toprule
%旧约&新约&末后\\  \midrule
乐意的灵(诗)&你的灵&我的灵\\  \midrule
耶和华的靈(赛士)&谋略和能力的灵(赛)&知识和敬畏耶和华的灵(赛)\\  \midrule
智慧和聪明的灵(赛)&公義的靈(赛)&焚燒的靈(赛)\\  \midrule
他的靈(民弗)&神的灵(创)&恳求的灵(亚)\\  \midrule

死里复活者的灵&永远的灵&圣善(洁)的灵\\  \midrule
神荣耀的灵&生命圣灵the Spirit of life from God&赐生命圣灵\\  \midrule
父的灵(太)&养子义子的灵(罗)&兒子的靈(加)\\  \midrule
神的灵(太罗林前)&基督的灵(罗彼前)&耶稣基督的灵(腓)\\  \midrule
耶稣的灵(徒)&永生神的灵(林后)&主的靈(路徒林后)\\  \midrule
保惠師/训慰师(约)&真理的圣灵(约)&七靈(启)\\  \midrule
\end{tabular}
\end{table}





\subsubsection{圣灵图像}
\subsubsubsection{风(Wind) (“灵,Pneuma" {G:4151}可译作“风”)}
\begin{table}[H]
\centering
\begin{tabular}{p{6em}|p{40em}}  \toprule
经文&内容\\  \midrule
约3:8&我實實在在地告訴你：人若不是從水和聖靈生的，就不能進神的國。從肉身生的就是肉身，從靈生的就是靈。我說你們必須重生，你不要以為稀奇。風隨著意思吹，你聽見風的響聲，卻不曉得從哪裡來，往哪裡去\textcolor{red}{(有主权 His Sovereignty)}。凡從聖靈生的，也是如此。\\  \midrule

林前12:4-11&恩賜原有分別，聖靈卻是一位；職事也有分別，主卻是一位；功用也有分別，神卻是一位，在眾人裡面運行一切的事。聖靈顯在各人身上，是叫人得益處。這人蒙聖靈賜他智慧的言語，...又叫一人能說方言，又叫一人能翻方言。這一切都是這位聖靈所運行，\textcolor{red}{隨己意}分給各人的。\\  \midrule

徒2:2-4&忽然從天上有響聲下來,好像一陣大風吹過,充滿了他們所坐的屋子;又有舌頭如火焰顯現出來,分開落在他們各人頭上。他們就都被聖靈充滿,按著聖靈所賜的口才說起別國的話來。 \\  \midrule
徒1:8&但聖靈降臨在你們身上，你們就必得著\textcolor{red}{能力}，並要在耶路撒冷、猶太全地和撒馬利亞，直到地極，作我的見證。\\  \midrule
\end{tabular}
\end{table}


\subsubsubsection{鸽子(Dove)}
\begin{table}[H]
\centering
\begin{tabular}{p{4em}|p{44em}}  \toprule
经文&内容\\  \midrule
太3:16&耶稣受了洗, 随即从水里上来. 天忽然为他开了, 他就看见神的灵, 彷佛鸽子降下, 落在他身上\\  \midrule
路3:22&圣灵降临在他身上, 形状彷佛鸽子…\\  \midrule
可1:10&他從水裡一上來，就看見天裂開了，聖靈彷彿鴿子降在他身上； 11 又有聲音從天上來說：「你是我的愛子，我喜悅你。\\  \midrule
约1:32&约翰又作见证说: 我曾看见圣灵彷佛鸽子, 从天降下, 住在他的身上”(\\  \midrule
%&鸽子是洁净的动物,故可用作献祭\\  \midrule
利5:7&他的力量若不夠獻一隻羊羔，就要因所犯的罪，把兩隻斑鳩或是\textcolor{red}{兩隻雛鴿}帶到耶和華面前為贖愆祭，一隻做贖罪祭，一隻做燔祭。\\  \midrule
利12:6&滿了潔淨的日子，無論是為男孩是為女孩，她要把一歲的羊羔為燔祭，\textcolor{red}{一隻雛鴿}或是一隻斑鳩為贖罪祭，帶到會幕門口交給祭司。\\  \midrule
歌6:9&我的\textcolor{red}{鴿子}，我的完全人，只有這一個是她母親獨生的，是生養她者所寶愛的。眾女子見了就稱她有福，王后妃嬪見了也讚美她。\\  \midrule
%&鸽子有温和的性情(The Dove is Gentle in Manner)\\  \midrule
太10:16&我差你們去，如同羊進入狼群，所以你們要靈巧像蛇，\textcolor{red}{馴良像鴿子}。\\  \midrule
路4:18-19&主的靈在我身上，因為他用膏膏我，叫我傳福音給貧窮的人，差遣我報告被擄的得釋放、瞎眼的得看見，叫那受壓制的得自由，報告神悅納人的禧年。\\  \midrule
赛42:3&壓傷的蘆葦他不折斷，將殘的燈火他不吹滅，他憑真實將公理傳開。\\  \midrule
%&鸽子有哀鸣的表现(The Dove is Mournful in Expression): \\  \midrule
赛38:14&我像燕子呢喃，像白鶴鳴叫，又像\textcolor{red}{鴿子哀鳴}，我因仰觀，眼睛困倦。耶和華啊，我受欺壓，求你為我作保！\\  \midrule
约11:35&耶穌哭了。\\  \midrule
太23:37-39&耶路撒冷啊，耶路撒冷啊！你常殺害先知，又用石頭打死那奉差遣到你這裡來的人。我多次願意聚集你的兒女，好像母雞把小雞聚集在翅膀底下，只是你們不願意。看哪，你們的家成為荒場留給你們！我告訴你們：從今以後，你們不得再見我，直等到你們說：『奉主名來的是應當稱頌的！』\\  \midrule
\end{tabular}
\end{table}

\subsubsubsection{油(Oil)-膏抹先知、祭司和君王}
\begin{table}[H]
\centering
\begin{tabular}{p{4em}|p{44em}}  \toprule
经文&内容\\  \midrule
%油也是圣灵的表记. 在旧约中, 人用油膏抹先知、祭司和君王, 以说明圣灵的事奉. 
亚4:1-14,建殿&那與我說話的天使又來叫醒我,好像人睡覺被喚醒一樣。他問我說：「你看見了什麼?」我說：「我看見了一個純金的燈臺，頂上有燈盞，燈臺上有七盞燈,每盞有七個管子。旁邊有兩棵橄欖樹,一棵在燈盞的右邊,一棵在燈盞的左邊。」我問與我說話的天使說：「主啊，這是什麼意思?」與我說話的天使回答我說：「你不知道這是什麼意思嗎?」我說：「主啊，我不知道。」他對我說：「這是耶和華指示所羅巴伯的。萬軍之耶和華說：\textcolor{red}{不是倚靠勢力，不是倚靠才能，乃是倚靠我的靈方能成事}。大山哪,你算什麼呢?在所羅巴伯面前,你必成為平地。他必搬出一塊石頭,安在殿頂上,人且大聲歡呼說：『願恩惠，恩惠歸於這殿!』」耶和華的話又臨到我說： 「所羅巴伯的手立了這殿的根基，他的手也必完成這工，你就知道萬軍之耶和華差遣我到你們這裡來了。誰藐視這日的事為小呢？這七眼乃是耶和華的眼睛,遍察全地,見所羅巴伯手拿線砣就歡喜。」我又問天使說：「這燈臺左右的兩棵橄欖樹是什麼意思?」我二次問他說：「這兩根橄欖枝在兩個流出\textcolor{red}{金色油}的金嘴旁邊是什麼意思?」他對我說：「你不知道這是什麼意思嗎?」我說：「主啊,我不知道。」他說：「這是兩個受膏者,站在普天下主的旁邊。」\\  \midrule 
撒上10:1,6,9-10&撒母耳拿瓶膏油倒在掃羅的頭上，與他親嘴，說：「這不是耶和華膏你做他產業的君嗎？...耶和華的靈必大大感動你，你就與他們一同受感說話，你要變為新人。...掃羅轉身離別撒母耳，神就賜他一個新心。當日這一切兆頭都應驗了。掃羅到了那山，有一班先知遇見他，神的靈大大感動他，他就在先知中受感說話。\\  \midrule
路4:18&主的靈在我身上，因為他用膏膏我，叫我傳福音給貧窮的人，差遣我報告被擄的得釋放、瞎眼的得看見，叫那受壓制的得自由.\\  \midrule
\end{tabular}
\end{table}


\subsubsubsection{舌头}


\subsubsubsection{火(Fire)}
\begin{table}[H]
\centering
\begin{tabular}{p{6em}|p{40em}}  \toprule
经文&内容\\  \midrule
徒2:3&又有舌頭如火焰顯現出來，分開落在他們各人頭上。 4 他們就都被聖靈充滿，按著聖靈所賜的口才說起別國的話來。\\  \midrule
出3:1-8&耶和华的使者从荆棘里火焰中向摩西显现… 神说: 不要近前来, 当把你脚上的鞋脱下来, 因为你所站之地是圣地\\  \midrule

出19:18&西乃全山冒烟, 因为耶和华在火中降于山上, 山的烟气上腾, 如烧瑶一般. 遍山大大的震动\\  \midrule
申4:24&因为耶和华你的神乃是烈火, 是忌邪的神\\  \midrule
来12:29&因為我們的神乃是烈火。\\  \midrule

 
%&论到圣灵降临, 三一神学院(Trinity College)研究院的院长殷保罗(Paul P. Enns)评述道: “火象征神的显现. … 这显现(指圣灵五旬节降临一事)说明了有神的准许, 彼得跟着宣讲耶稣的复活,这火则说明神核准了彼得的信息(比较 \\  \midrule
利9:24&有火從耶和華面前出來，在壇上燒盡燔祭和脂油，眾民一見，就都歡呼，俯伏在地。\\  \midrule
王上18:38-39&於是，耶和華降下火來，燒盡燔祭、木柴、石頭、塵土，又燒乾溝裡的水。 39 眾民看見了，就俯伏在地，說：「耶和華是神！耶和華是神！」\\  \midrule
利10:1-3&亞倫的兒子拿答、亞比戶各拿自己的香爐，盛上火，加上香，在耶和華面前獻上凡火，是耶和華沒有吩咐他們的， 2 就有火從耶和華面前出來，把他們燒滅，他們就死在耶和華面前。 3 於是摩西對亞倫說：「這就是耶和華所說：『我在親近我的人中要顯為聖，在眾民面前我要得榮耀。』」\\  \midrule
\end{tabular}
\end{table}
 

\subsubsubsection{水(Water)}
\begin{table}[H]
\centering
\begin{tabular}{p{4em}|p{46em}}  \toprule
经文&内容\\  \midrule
静止的水&圣经也用水来象征圣灵. 李德尔(John M. Riddle)“静止的水(still water)象征神的道, \\  \midrule
活水&流动的水或活水(springing or moving water)却是圣灵的写照, \\  \midrule
海水&(seawater)则是列国万民的写照.”[7]\\  \midrule
约4:14&人若喝我所赐的水就永远不渴. 我所赐的水, 要在他里头成为泉源, 直涌到永生.\\  \midrule
約7:37-39&节期的末日, 就是最大之日, 耶稣站着高声说: 人若渴了, 可以到我这里来喝. 信我的人, 就如经上所说, 从他腹中要流出活水的江河来(KJV: shall flow rivers of living water). 耶稣这话是指着信他之人, 要受圣灵说的, 那时还没有赐下圣灵来, 因为耶稣尚未得着荣耀”\\  \midrule
賽44:2-4&造做你，又從你出胎造就你，並要幫助你的耶和華如此說：『我的僕人雅各，我所揀選的耶書崙哪，不要害怕！ 3 因為我要將水澆灌口渴的人，將河澆灌乾旱之地，我要將我的靈澆灌你的後裔，將我的福澆灌你的子孫。 4 他們要發生在草中，像溪水旁的柳樹。\\  \midrule
结36:25-27&我必用清水灑在你們身上，你們就潔淨了。我要潔淨你們，使你們脫離一切的汙穢，棄掉一切的偶像。我也要賜給你們一個新心，將新靈放在你們裡面，又從你們的肉體中除掉石心，賜給你們肉心。我必將我的靈放在你們裡面，使你們順從我的律例，謹守遵行我的典章。\\  \midrule
結47:1-12&他帶我回到殿門，見殿的門檻下有水往東流出（原來殿面朝東）。這水從檻下，由殿的右邊，在祭壇的南邊往下流。他帶我出北門，又領我從外邊轉到朝東的外門，見水從右邊流出。他手拿準繩往東出去的時候，量了一千肘，使我趟過水，水到踝子骨。他又量了一千肘，使我趟過水，水就到膝。再量了一千肘，使我趟過水，水便到腰。又量了一千肘，水便成了河，使我不能趟過，因為水勢漲起，成為可洑的水，不可趟的河。他對我說：「人子啊，你看見了什麼？」他就帶我回到河邊。我回到河邊的時候，見在河這邊與那邊的岸上有極多的樹木。他對我說：「這水往東方流去，必下到亞拉巴，直到海。所發出來的水必流入鹽海，使水變甜。這河水所到之處，凡滋生的動物都必生活，並且因這流來的水必有極多的魚，海水也變甜了。這河水所到之處，百物都必生活。必有漁夫站在河邊，從隱基底直到隱以革蓮，都做曬網之處。那魚各從其類，好像大海的魚，甚多。只是泥濘之地與窪濕之處不得治好，必為鹽地。在河這邊與那邊的岸上必生長各類的樹木，其果可做食物，葉子不枯乾，果子不斷絕。每月必結新果子，因為這水是從聖所流出來的。樹上的果子必做食物，葉子乃為治病。\\  \midrule
\end{tabular}
\end{table}
 
\subsubsubsection{凭据(抵押)}
\begin{table}[H]
\centering
\begin{tabular}{p{6em}|p{43em}}  \toprule
经文&内容\\  \midrule
民13：233&他们到了「以实各谷，从那里砍了葡萄树的一枝，上头有一挂葡萄。」\\  \midrule
太25:1-13&那時，天國好比十個童女拿著燈出去迎接新郎。其中有五個是愚拙的，五個是聰明的。愚拙的拿著燈，卻不預備油；聰明的拿著燈，又預備油在\textcolor{red}{器皿}裡。新郎遲延的時候，她們都打盹睡著了。半夜有人喊著說：『新郎來了，你們出來迎接他！』那些童女就都起來收拾燈。愚拙的對聰明的說：『請分點油給我們，因為我們的\textcolor{red}{燈要滅了}。』聰明的回答說：『恐怕不夠你我用的，不如你們自己到賣油的那裡去買吧。』她們去買的時候，新郎到了，那預備好了的同他進去坐席，門就關了。其餘的童女隨後也來了，說：『主啊，主啊，給我們開門！』他卻回答說：『我實在告訴你們：我\textcolor{red}{不認識}你們。』所以，你們要警醒，因為那日子、那時辰，你們不知道。\\  \midrule
林后1:22&那在基督里坚固我们和你们，并且膏我们的，就是神；他又用印印(sealed us)了我们, 并赐圣灵在我们心里作凭据”(“凭据”原文作质; KJV: the earnest of the Spirit;) \\  \midrule
林后5:5&為此培植我們的就是神，他又賜給我們聖靈做憑據 。(“凭据” arrabôn {G:728})可解为“抵押、首期的付款、订金", 是交易中先付出的价值, 以保证法律上的权利, 和条约的有效性. \\  \midrule
弗1:14&这圣灵是我们得基业的凭据(质), 直等到神之民(产业)被赎, 使他的荣耀得着称赞. \\  \midrule
\end{tabular}
\end{table}




 

\subsubsubsection{印记(Seal)}
\begin{table}[H]
\centering
\begin{tabular}{p{6em}|p{40em}}  \toprule
经文&内容\\  \midrule
耶32:10-14&我便向我叔叔的兒子哈拿篾買了亞拿突的那塊地，平了十七舍客勒銀子給他。我在契上畫押，將契封緘，又請見證人來，並用天平將銀子平給他。我便將照例按規所立的買契，就是封緘的那一張和敞著的那一張，當著我叔叔的兒子哈拿篾和畫押做見證的人，並坐在護衛兵院內的一切猶大人眼前，交給瑪西雅的孫子、尼利亞的兒子巴錄。當著他們眾人眼前，我囑咐巴錄說：「萬軍之耶和華以色列的神如此說：要將這封緘的和敞著的兩張契放在瓦器裡，可以存留多日。因為萬軍之耶和華以色列的神如此說：將來在這地必有人再買房屋、田地和葡萄園。」\\  \midrule


帖8:8&国王是用自己的戒子「盖印」在他所颁诏的信函及文件上\\  \midrule 
但&当巴比伦王将但以理「封」进狮子坑时，他就不能逃出来\\  \midrule 
太27:65-66&你们有看守的兵，去罢，尽你们所能的，把守妥当。他们就带着看守的兵同去，封了石头，将坟墓把守妥当。\\  \midrule 


&圣灵印记的重要经文：罗8：12-17；弗1：13-14\\  \midrule 

结9:4&耶和華對他說：「你去走遍耶路撒冷全城，那些因城中所行可憎之事嘆息哀哭的人，畫記號在額上。」\\  \midrule 
启7:2-4&我又看見另有一位天使，從日出之地上來，拿著永生神的印。他就向那得著權柄能傷害地和海的四位天使大聲喊著說：「地與海並樹木你們不可傷害，等我們印了我們神眾僕人的額。」我聽見以色列人各支派中受印的數目有十四萬四千：...\\  \midrule 

来10:14、15&他一次献祭，便叫那得以成圣的人永远完全。圣灵也对我们作见证。\\  \midrule 
罗马书八：1&圣灵与我们的心同证我们是神的儿女。\\  \midrule 
加拉太书四：6、7&你们既为儿子，神就差他儿子的灵，进入你们的心，呼叫阿爸、父。可见从此以后，你不是奴仆，乃是儿子了；既是儿子，就靠着神为后嗣。\\  \midrule 
约一5:10-13&信神儿子的，就有这见证在他心里。这见证，就是神赐给我们永生，这永生也是在他儿子里面。人有了神的儿子就有生命，没有神的儿子就没有生命。我将这些话写给你们信奉神儿子之名的人，要叫你们知道自己有永生。\\  \midrule 


约6：27&因为人子是父神所印证的。\\  \midrule 
約3:33&那領受他見證的，就印上印，證明神是真的。\\  \midrule
林后1:22&他又用印印了我們，並賜聖靈在我們心裡做憑據。\\  \midrule
弗1:13&你們既聽見真理的道，就是那叫你們得救的福音，也信了基督，既然信他，就受了所應許的聖靈為印記。 這聖靈是我們得基業的憑據，直等到神之民被贖，使他的榮耀得著稱讚。\\  \midrule 
弗4:30&不要叫神的聖靈擔憂，你們原是受了他的印記，等候得贖的日子來到。 \\  \midrule
%太27:66&提到, 罗马官方的印记(seal, 在此可译作“封印”), 是一块打了印的石头. 按象征性的意义, “盖上印表示认同…古时所养的所有动物身上都带着印记(标记), 以表明它的拥有者是谁, 也表明拥有者所给予的保障”.[10] 这古代的印记, 就犹如现今印在牛畜身上的烙印(比较 \\  \midrule 
%赛44:5&\\  \midrule 


%&殷保罗(Paul P. Enns)指出, 关乎圣灵的印记, 有几个重要的真理: (1)印记说明神的拥有权; 圣灵加给信徒印记, 说明信徒是属于神的; (2)印记表明稳妥: 印记是永久性的, “不要叫神的圣灵担忧. 你们原是受了他的印记, 等候得赎的日子来到”(弗4:30); (3)印记代表权柄: 正如罗马官方的权柄所及, 罗马的印记就发生功效, 照样地, 神赐给信徒圣灵为印记, 就在信徒身上拥有权柄.[11]
\end{tabular}
\end{table}
 


 
\subsubsubsection{圣灵的殿}
\begin{table}[H]
\centering
\begin{tabular}{p{6em}|p{40em}}  \toprule
经文&内容\\  \midrule
林前6:19-20&豈不知你們的身子就是聖靈的殿嗎？這聖靈是從神而來，住在你們裡頭的。並且你們不是自己的人，因為你們是重價買來的，所以要在你們的身子上榮耀神。\\  \midrule
加5:25-26&我們若是靠聖靈得生，就當靠聖靈行事。不要貪圖虛名，彼此惹氣，互相嫉妒。\\  \midrule
\end{tabular}
\end{table}




\subsection{旧约中聖靈的工作，记载}
\subsubsection{圣灵的两次浇灌(秋雨春雨)预言}
\begin{table}[H]
\centering
\begin{tabular}{p{4em}|p{44em}}  \toprule
经文&内容 \\  \midrule
申11:11-14&你們要過去得為業的那地乃是有山有谷，雨水滋潤之地，是耶和華你神所眷顧的。從歲首到年終，耶和華你神的眼目時常看顧那地。 「你們若留意聽從我今日所吩咐的誡命，愛耶和華你們的神，盡心、盡性侍奉他，他必按時降\textcolor{red}{秋雨春雨}在你們的地上，使你們可以收藏五穀、新酒和油，\\  \midrule
耶5:23-24&但這百姓有背叛忤逆的心，他們叛我而去。心內也不說：『我們應當敬畏耶和華我們的神，他按時賜雨，就是\textcolor{red}{秋雨春雨 (both the former and the latter rains)}，又為我們定收割的節令，永存不廢。』\\  \midrule

珥2:22-23&田野的走獸啊，不要懼怕！因為曠野的草發生，樹木結果，無花果樹、葡萄樹也都效力。錫安的民哪，你們要快樂，為耶和華你們的神歡喜！因他賜給你們合宜的\textcolor{red}{秋雨}，為你們降下甘霖，就是\textcolor{red}{秋雨、春雨}，和先前一樣。\\  \midrule

箴16:15&王的臉光使人有生命，王的恩典好像\textcolor{red}{春雲時雨(the latter rain)}。\\  \midrule

%賽44:3-5&因為我要將水澆灌口渴的人，將河澆灌乾旱之地，我要將我的靈澆灌你的後裔，將我的福澆灌你的子孫。 4 他們要發生在草中，像溪水旁的柳樹。 5 這個要說「我是屬耶和華的」，那個要以雅各的名自稱，又一個要親手寫「歸耶和華的」，並自稱為以色列。』\\  \midrule


耶3:2-3&你坐在道旁等候，好像阿拉伯人在曠野埋伏一樣，並且你的淫行、邪惡玷汙了全地。因此甘霖停止，\textcolor{red}{春(晚)雨不降(no latter rain)}，你還是有娼妓之臉，不顧羞恥。\\  \midrule

何6:3?&我們務要認識耶和華，竭力追求認識他。他出現確如晨光，他必臨到我們像甘雨，像滋潤田地的\textcolor{red}{春雨}。\\  \midrule

亞10:1&當\textcolor{red}{春雨}的時候，你們要向發閃電的耶和華求雨，他必為眾人降下甘霖，使田園生長菜蔬。\\  \midrule

伯29:19-23&我的根長到水邊，露水終夜沾在我的枝上。我的榮耀在身上增新，我的弓在手中日強。』人聽見我而仰望，靜默等候我的指教。我說話之後，他們就不再說，我的言語像雨露滴在他們身上。他們仰望我如仰望雨，又張開口如切慕\textcolor{red}{春雨( the latter rain)}。\\  \midrule
\end{tabular}%
\end{table}



\subsubsection{圣灵降临-主耶稣在世时圣灵还没降下 约7:37-39/16:7}
%\subsubsubsection{主耶稣在世时圣灵还没降下 约7:37-39/16:7}
\begin{table}[H]
\centering
\begin{tabular}{p{6em}|p{40em}}  \toprule
经文&内容\\  \midrule

\end{tabular}
\end{table}


\subsubsection{应许给人圣灵的新约预言}
\begin{table}[H]
\centering
\begin{tabular}{p{4em}|p{44em}}  \toprule
经文&内容 \\  \midrule
太3:11-12&我是用水給你們施洗, 叫你們悔改; 但那在我以後來的, 能力比我更大, 我就是給他提鞋也不配, 他要用聖靈與火給你們施洗。他手裡拿著簸箕, 要揚淨他的場, 把麥子收在倉裡, 把糠用不滅的火燒盡了。\\  \midrule

可1:8&我是用水給你們施洗，他卻要用聖靈給你們施洗。\\  \midrule

路3:16&約翰說：「我是用水給你們施洗。但有一位能力比我更大的要來，我就是給他解鞋帶也不配，他要用聖靈與火給你們施洗。\\  \midrule

約1:26-27,32-33&約翰回答說：「我是用水施洗，但有一位站在你們中間，是你們不認識的，就是那在我以後來的，我給他解鞋帶也不配。」約翰又作見證說：「我曾看見聖靈彷彿鴿子從天降下，住在他的身上。我先前不認識他，只是那差我來用水施洗的對我說：『你看見聖靈降下來，住在誰的身上，誰就是用聖靈施洗的。』\\  \midrule

路13:13&你們雖然不好，尚且知道拿好東西給兒女，何況天父，豈不更將聖靈給求他的人嗎？\\  \midrule

路24:49&我要將我父所應許的降在你們身上，你們要在城裡等候，直到你們領受從上頭來的能力。」\\  \midrule

約4:10-14&耶穌回答說：「你若知道神的恩賜和對你說『給我水喝』的是誰，你必早求他，他也必早給了你活水。」...耶穌回答說：「凡喝這水的，還要再渴，人若喝我所賜的水，就永遠不渴。我所賜的水要在他裡頭成為泉源，直湧到永生。」\\  \midrule

約7:37-39&節期的末日，就是最大之日，耶穌站著高聲說：「人若渴了，可以到我這裡來喝！信我的人就如經上所說，從他腹中要流出活水的江河來。」耶穌這話是指著信他之人要受聖靈說的。那時還沒有賜下聖靈來，因為耶穌尚未得著榮耀。\\  \midrule

約16:7&然而我將真情告訴你們，我去是於你們有益的。我若不去，保惠師就不到你們這裡來；我若去，就差他來。\\  \midrule

徒1:4-5&耶穌和他們聚集的時候，囑咐他們說：「不要離開耶路撒冷，要等候父所應許的，就是你們聽見我說過的。约翰是用水施洗。但不多几日，你们要受圣灵的洗。\\  \midrule

徒2:17-20&神說：在末後的日子，我要將我的靈澆灌凡有血氣的，你們的兒女...在那些日子，我要將我的靈澆灌我的僕人和使女，他們就要說預言。在天上我要顯出奇事...，這都在主大而明顯的日子未到以前。\\  \midrule
啟22:17&聖靈和新婦都說：「來！」聽見的人也該說：「來！」口渴的人也當來，願意的都可以白白取生命的水喝。\\  \midrule

\end{tabular}%
\end{table}

\subsubsection{应许给人圣灵的旧约预言}
\begin{table}[H]
\centering
\begin{tabular}{p{4em}|p{44em}}  \toprule
经文&内容 \\  \midrule

詩36:8-9&他們必因你殿裡的肥甘得以飽足，你也必叫他們喝你樂河的水。因為在你那裡有生命的源頭，在你的光中我們必得見光。\\  \midrule

箴1:23&你們當因我的責備回轉，我要將我的\textcolor{red}{靈}澆灌你們，將我的話指示你們。\\  \midrule

賽4:3&主以\textcolor{red}{公義的靈}和\textcolor{red}{焚燒的靈}，將錫安女子的汙穢洗去，又將耶路撒冷中殺人的血除淨。那時，剩在錫安留在耶路撒冷的，就是一切住耶路撒冷在生命冊上記名的，必稱為聖。\\  \midrule
賽28:6&到那日，萬軍之耶和華必做他餘剩之民的榮冠華冕， 6也做了在位上行審判者\textcolor{red}{公平之靈}，並城門口打退仇敵者的力量。\\  \midrule

賽32:15-20&等到\textcolor{red}{聖靈}從上澆灌我們，曠野就變為肥田，肥田看如樹林。那時，公平要居在曠野，公義要居在肥田。公義的果效必是平安，公義的效驗必是平穩，直到永遠。我的百姓必住在平安的居所，安穩的住處，平靜的安歇所。但要降冰雹打倒樹林，城必全然拆平。你們在各水邊撒種牧放牛驢的有福了！\\  \midrule


賽41:17-20&困苦窮乏人尋求水卻沒有，他們因口渴舌頭乾燥，我耶和華必應允他們，我以色列的神必不離棄他們。我要在淨光的高處開江河，在谷中開泉源；我要使沙漠變為水池，使乾地變為湧泉。我要在曠野種上香柏樹、皂莢樹、番石榴樹和野橄欖樹，我在沙漠要把松樹、杉樹並黃楊樹一同栽植，...。\\  \midrule


賽43:19-21&我必在曠野開道路，在沙漠開江河。 20 野地的走獸必尊重我，野狗和鴕鳥也必如此，因我使曠野有水，使沙漠有河，好賜給我的百姓、我的選民喝。 21 這百姓是我為自己所造的，好述說我的美德。\\  \midrule

賽44:2-5&我的僕人雅各，我所揀選的耶書崙哪，不要害怕！因為我要將水澆灌口渴的人，將河澆灌乾旱之地，我要將\textcolor{red}{我的靈}澆灌你的後裔，將我的福澆灌你的子孫。他們要發生在草中，像溪水旁的柳樹。這個要說「我是屬耶和華的」，那個要以雅各的名自稱，又一個要親手寫「歸耶和華的」，並自稱為以色列。\\  \midrule

賽55:1-2&你們一切乾渴的，都當就近水來！沒有銀錢的，也可以來！你們都來，買了吃。不用銀錢，不用價值，也來買酒和奶。 ...你們要留意聽我的話，就能吃那美物，得享肥甘，心中喜樂。\\  \midrule

賽59:20-21&「必有一位救贖主來到錫安，雅各族中轉離過犯的人那裡。至於我，與他們所立的約乃是這樣：\textcolor{red}{我加給你的靈，傳給你的話，必不離你的口}，也不離你後裔與你後裔之後裔的口，從今直到永遠。」\\  \midrule

結11:19-20&我要使他們有合一的心，也要將\textcolor{red}{新靈}放在他們裡面，又從他們肉體中除掉石心，賜給他們肉心，使他們順從我的律例，謹守遵行我的典章。他們要做我的子民，我要做他們的神。\\  \midrule

結18:31-32&你們要將所犯的一切罪過盡行拋棄，自做一個新心和\textcolor{red}{新靈}。以色列家啊，你們何必死亡呢？主耶和華說：我不喜悅那死人之死，所以你們當回頭而存活。」\\  \midrule
結36:26-28&我也要賜給你們一個新心，將\textcolor{red}{新靈}放在你們裡面，又從你們的肉體中除掉石心，賜給你們肉心。 我必將\textcolor{red}{我的靈}放在你們裡面，使你們順從我的律例，謹守遵行我的典章。你們必住在我所賜給...\\  \midrule
結37:14&我必將\textcolor{red}{我的靈}放在你們裡面，你們就要活了。我將你們安置在本地，你們就知道我耶和華如此說，也如此成就了。這是耶和華說的。\\  \midrule

結39:25-29&主耶和華如此說：我要使雅各被擄的人歸回，要憐憫以色列全家，又為我的聖名發熱心。他們在本地安然居住，無人驚嚇，是我將他們從萬民中領回，從仇敵之地召來。我在許多國的民眼前，在他們身上顯為聖的時候，他們要擔當自己的羞辱....我必不再留他們一人在外邦。我也不再掩面不顧他們，因我已將\textcolor{red}{我的靈}澆灌以色列家。這是主耶和華說的。\\  \midrule
珥2:28-31&以後我要將\textcolor{red}{我的靈}澆灌凡有血氣的,你們的兒女...。在那些日子,我要將\textcolor{red}{我的靈}澆灌我的僕人和使女。在天上地下,我要顯出奇事,...,這都在耶和華大而可畏的日子未到以前。到那時候凡求告耶和華名的，就必得救。...在錫安山耶路撒冷必有逃脫的人,在剩下的人中必有耶和華所召的。\\  \midrule

亞4:6&這是耶和華指示所羅巴伯的:不是倚靠勢力不是倚靠才能,乃是倚靠\textcolor{red}{我的靈}方能成事\\  \midrule
亞8:12&他們必平安撒種,葡萄樹必結果子,地土必有出產,天也必降\textcolor{red}{甘露}。我要使這餘剩的民享受這一切的福。\\  \midrule

亞12:10&我必將那\textcolor{red}{施恩叫人懇求的靈}澆灌大衛家和耶路撒冷的居民。他們必仰望我，就是他們所扎的；必為我悲哀，如喪獨生子；又為我愁苦，如喪長子。\\ \midrule 

\end{tabular}%
\end{table}



\subsubsection{創造諸天,掌管自然界}
\begin{table}[H]
\centering
\begin{tabular}{p{6em}p{40em}}  \toprule
经节&经文 \\  \midrule
創1:2&地是空虛混沌，淵面黑暗，神的靈運行在水面上。\\  \midrule
伯26:7-13&約伯回答說：...神將北極鋪在空中，將大地懸在虛空。 ...他以能力攪動大海，他藉知識打傷拉哈伯； 藉他的靈使天有裝飾，他的手刺殺快蛇。\\  \midrule

詩33:6&諸天藉耶和華的命而造，萬象藉他口中的氣而成。\\  \midrule
賽34:8-17&因耶和華有報仇之日，為錫安的爭辯有報應之年。 9 以東的河水要變為石油，塵埃要變為硫磺，地土成為燒著的石油。 10 晝夜總不熄滅，煙氣永遠上騰，必世世代代成為荒廢，永永遠遠無人經過。....你們要查考、宣讀耶和華的書，這都無一缺少，無一沒有伴偶。因為我的口已經吩咐，\textcolor{red}{他的靈}將他們聚集。 17 他也為他們拈鬮，又親手用準繩給他們分地，他們必永得為業，世世代代住在其間。\\  \midrule
\end{tabular}%\caption{創造}  
\end{table}





\subsubsection{赐人生命}
\begin{table}[H]
\centering
\begin{tabular}{p{6em}p{40em}}  \toprule
经节&经文 \\  \midrule
創2:7&耶和華神用地上的塵土造人，將生氣吹在他鼻孔裡，他就成了有靈的活人，名叫亞當。\\  \midrule
創6:3&耶和華說：「人既屬乎血氣，我的靈就不永遠住在他裡面，然而他的日子還可到一百二十年。」\\  \midrule
伯27:3&（我的生命尚在我裡面，神所賜呼吸之氣仍在我的鼻孔內）\\  \midrule
伯32:6-8&布西人巴拉迦的兒子以利戶回答說：「我年輕，你們老邁，因此我退讓，不敢向你們陳說我的意見。 7 我說，年老的當先說話，壽高的當以智慧教訓人。 8 但在人裡面有靈，全能者的氣使人有聰明。\\  \midrule

伯33:4-7&神的靈造我，全能者的氣使我得生。 5 你若回答我，就站起來，在我面前陳明。 6 我在神面前與你一樣，也是用土造成。 7 我不用威嚴驚嚇你，也不用勢力重壓你。\\  \midrule
伯34:12-15&神必不作惡，全能者也不偏離公平。 13 誰派他治理地，安定全世界呢？ 14 他若專心為己，將靈和氣收歸自己， 15 凡有血氣的就必一同死亡，世人必仍歸塵土。\\  \midrule


民27:16-18&願耶和華萬人之靈的神，立一個人治理會眾，\\  \midrule
詩104:30&你發出你的靈，牠們便受造；你使地面更換為新。\\\bottomrule
瑪2:15&雖然神有靈的餘力能造多人，他不是單造一人嗎？為何只造一人呢？乃是他願人得虔誠的後裔。所以當謹守你們的心，誰也不可以詭詐待幼年所娶的妻。\\  \midrule
賽40:6-8&有人聲說：「你喊叫吧！」有一個說：「我喊叫什麼呢？」說：「凡有血氣的盡都如草，他的美容都像野地的花。 7 草必枯乾，花必凋殘，因為耶和華的氣吹在其上——百姓誠然是草！ 8 草必枯乾，花必凋殘，唯有我們神的話必永遠立定。」\\  \midrule
傳12:7&塵土仍歸於地，靈仍歸於賜靈的神。\\  \midrule
賽42:5-9&創造諸天，鋪張穹蒼，將地和地所出的一併鋪開，賜氣息給地上的眾人，又賜靈性給行在其上之人的神耶和華，他如此說： 6 「我耶和華憑公義召你，必攙扶你的手，保守你，使你做眾民的中保，做外邦人的光....」\\  \midrule

結37:12-14&主耶和華如此說：我的民哪，我必開你們的墳墓，使你們從墳墓中出來，領你們進入以色列地。 13 我的民哪，我開你們的墳墓，使你們從墳墓中出來，你們就知道我是耶和華。 14 我必將我的靈放在你們裡面，你們就要活了。我將你們安置在本地，你們就知道我耶和華如此說，也如此成就了。這是耶和華說的。\\  \midrule
約20:22&說了這話，就向他們吹一口氣，說：「你們受聖靈！
\\\bottomrule
\end{tabular}
\end{table}





\subsubsection{赐人能力 prophetic spirit}
\begin{table}[H]
\centering
\begin{tabular}{p{6em}p{40em}}  \toprule
经节&经文 \\  \midrule
申34:9&嫩的兒子約書亞因為摩西曾按手在他頭上，就被智慧的靈充滿，以色列人便聽從他，照著耶和華吩咐摩西的行了。\\  \midrule
書27:1821&耶和華對摩西說：「嫩的兒子約書亞是心中有聖靈的，你將他領來，按手在他頭上， 19 使他站在祭司以利亞撒和全會眾面前，囑咐他。 20 又將你的尊榮給他幾分，使以色列全會眾都聽從他。 21 他要站在祭司以利亞撒面前，以利亞撒要憑烏陵的判斷，在耶和華面前為他求問。他和以色列全會眾都要遵以利亞撒的命出入。」\\  \midrule

士3:9-11&以色列人呼求耶和華的時候，耶和華就為他們興起一位拯救者救他們，就是迦勒兄弟基納斯的兒子俄陀聶。 10 耶和華的靈降在他身上，他就做了以色列的士師，出去爭戰。耶和華將美索不達米亞王古珊利薩田交在他手中，他便勝了古珊利薩田。 11 於是國中太平四十年。基納斯的兒子俄陀聶死了。\\  \midrule
士6:34-35&耶和華的靈降在基甸身上，他就吹角，亞比以謝族都聚集跟隨他。 35 他打發人走遍瑪拿西地，瑪拿西人也聚集跟隨他。又打發人去見亞設人、西布倫人、拿弗他利人，他們也都出來與他們會合。\\  \midrule
士11:29&耶和華的靈降在耶弗他身上，他就經過基列和瑪拿西，來到基列的米斯巴，又從米斯巴來到亞捫人那裡。 \\  \midrule
士13:24-25&後來婦人生了一個兒子，給他起名叫參孫。孩子長大，耶和華賜福於他。 25 在瑪哈尼但，就是瑣拉和以實陶中間，耶和華的靈才感動他。\\  \midrule

士14:6&耶和華的靈大大感動參孫，他雖然手無器械，卻將獅子撕裂，如同撕裂山羊羔一樣。他行這事並沒有告訴父母。\\  \midrule
士14:19&耶和華的靈大大感動參孫，他就下到亞實基倫，擊殺了三十個人，奪了他們的衣裳，將衣裳給了猜出謎語的人。參孫發怒，就上父家去了。\\  \midrule

士15:14-15&參孫到了利希，非利士人都迎著喧嚷。耶和華的靈大大感動參孫，他臂上的繩就像火燒的麻一樣，他的綁繩都從他手上脫落下來。 15 他見一塊未乾的驢腮骨，就伸手拾起來，用以擊殺一千人。\\  \midrule

王上18:45&耶和華的靈(手)降在以利亞身上，他就束上腰，奔在亞哈前頭，直到耶斯列的城門。\\  \midrule

王下2:9-10&過去之後，以利亞對以利沙說：「我未曾被接去離開你，你要我為你做什麼，只管求我。」以利沙說：「願感動你的靈加倍地感動我。」 10 以利亞說：「你所求的難得。雖然如此，我被接去離開你的時候，你若看見我，就必得著，不然必得不著了。」\\  \midrule
王下2:15-16&住耶利哥的先知門徒從對面看見他，就說：「感動以利亞的靈感動以利沙了。」他們就來迎接他，在他面前俯伏於地， 16 對他說：「僕人們這裡有五十個壯士，求你容他們去尋找你師傅，或者耶和華的靈將他提起來，投在某山某谷。」\\  \midrule
\end{tabular}
\end{table}


\subsubsection{給人有聰明智慧-解梦}
\begin{table}[H]
\centering
\begin{tabular}{p{6em}p{42em}}  \toprule
经节&经文 \\  \midrule
創41:37-40&法老和他一切臣僕都以這事為妙。法老對臣僕說：「像這樣的人，有神的靈在他裡頭，我們豈能找得著呢？」 法老對約瑟說：「神既將這事都指示你，可見沒有人像你這樣有聰明有智慧。 \\  \midrule% 40 你可以掌管我的家，我的民都必聽從你的話，唯獨在寶座上我比你大。」 \\  \midrule
但5:11-12&在你國中有一人，他裡頭有聖神的靈，你父在世的日子，這人心中光明，又有聰明、智慧，好像神的智慧。你父尼布甲尼撒王，就是王的父，立他為術士、用法術的和迦勒底人並觀兆的領袖。在他裡頭有美好的靈性，又有知識、聰明，能圓夢，釋謎語，解疑惑。\\  \midrule%這人名叫但以理，尼布甲尼撒王又稱他為伯提沙撒。現在可以召他來，他必解明這意思。\\  \midrule
\end{tabular}
\end{table}



\subsubsection{以西结书}
\begin{table}[H]
\centering
\begin{tabular}{p{5em}p{45em}}  \toprule
经节&经文 \\  \midrule
結1:1-3&當三十年四月初五日，以西結在迦巴魯河邊被擄的人中，天就開了，得見神的異象。正是約雅斤王被擄去第五年四月初五日，在迦勒底人之地迦巴魯河邊，耶和華的話特特臨到布西的兒子祭司以西結，耶和華的\textcolor{red}{靈}降在他身上。 \\  \midrule


結2:2&他對我說：「人子啊，你站起來，我要和你說話。」 2 他對我說話的時候，\textcolor{red}{靈}就進入我裡面，使我站起來，我便聽見那位對我說話的聲音。 3 他對我說：「人子啊，我差你往悖逆的國民以色列人那裡去。他們是悖逆我的，他們和他們的列祖違背我直到今日。 4 這眾子面無羞恥，心裡剛硬，我差你往他們那裡去，你要對他們說：『主耶和華如此說。』 5 他們或聽或不聽（他們是悖逆之家），必知道在他們中間有了先知。\\  \midrule

結3:12-15&那時,\textcolor{red}{靈}將我舉起,我就聽見在我身後有震動轟轟的聲音說：「從耶和華的所在顯出來的榮耀是該稱頌的！」 我又聽見那活物翅膀相碰與活物旁邊輪子旋轉震動轟轟的響聲。於是\textcolor{red}{靈}將我舉起，帶我而去。我心中甚苦，靈性憤激，並且耶和華的\textcolor{red}{靈}在我身上大有能力。我就來到提勒亞畢——住在迦巴魯河邊被擄的人那裡，到他們所住的地方，在他們中間憂憂悶悶地坐了七日。\\  \midrule

結3:22-27&耶和華的\textcolor{red}{靈}在那裡降在我身上，他對我說：「你起來，往平原去，我要在那裡和你說話。」於是我起來往平原去，不料耶和華的榮耀正如我在迦巴魯河邊所見的一樣，停在那裡，我就俯伏於地。\textcolor{red}{靈}就進入我裡面，使我站起來。耶和華對我說：「你進房屋去，將門關上。人子啊，人必用繩索捆綁你，你就不能出去在他們中間來往。我必使你的舌頭貼住上膛，以致你啞口，不能做責備他們的人，他們原是悖逆之家。但我對你說話的時候，必使你開口，你就要對他們說：『主耶和華如此說。』聽的可以聽，不聽的任他不聽，因為他們是悖逆之家。\\  \midrule

結8:1-4&第六年六月初五日，我坐在家中，猶大的眾長老坐在我面前，在那裡主耶和華的\textcolor{red}{靈}降在我身上。我觀看，見有形象彷彿火的形狀，從他腰以下的形狀有火，從他腰以上有光輝的形狀，彷彿光耀的精金。他伸出彷彿一隻手的樣式，抓住我的一綹頭髮，\textcolor{red}{靈}就將我舉到天地中間，在神的異象中，帶我到耶路撒冷朝北的內院門口，在那裡有觸動主怒偶像的座位，就是惹動忌邪的。誰知，在那裡有以色列神的榮耀，形狀與我在平原所見的一樣。\\  \midrule


結11:1-8&\textcolor{red}{靈}將我舉起，帶到耶和華殿向東的東門，誰知，在門口有二十五個人，我見其中有民間的首領押朔的兒子雅撒尼亞和比拿雅的兒子毗拉提。耶和華對我說：「人子啊，這就是圖謀罪孽的人，在這城中給人設惡謀。他們說：『蓋房屋的時候尚未臨近，這城是鍋，我們是肉。』人子啊，因此你當說預言，說預言攻擊他們。」\textcolor{red}{耶和華的靈}降在我身上，對我說：「你當說：『耶和華如此說：以色列家啊，你們口中所說的，心裡所想的，我都知道。你們在這城中殺人增多，使被殺的人充滿街道。所以主耶和華如此說：「你們殺在城中的人就是肉，這城就是鍋，你們卻要從其中被帶出去。你們怕刀劍，我必使刀劍臨到你們。」\\  \midrule


結11:22-25&於是基路伯展開翅膀，輪子都在他們旁邊，在他們以上有以色列神的榮耀。 23 耶和華的榮耀從城中上升，停在城東的那座山上。\textcolor{red}{靈}將我舉起，在異象中，藉著\textcolor{red}{神的靈}將我帶進迦勒底地，到被擄的人那裡。我所見的異象就離我上升去了。 25 我便將耶和華所指示我的一切事都說給被擄的人聽。\\  \midrule
結37:1-14&\textcolor{red}{耶和華的靈}降在我身上，耶和華藉\textcolor{red}{他的靈}帶我出去，將我放在平原中，這平原遍滿骸骨。...主對我說：「人子啊，這些骸骨就是以色列全家。他們說：『我們的骨頭枯乾了，我們的指望失去了，我們滅絕淨盡了。』 所以你要發預言對他們說：『主耶和華如此說：我的民哪，我必開你們的墳墓，使你們從墳墓中出來，領你們進入以色列地。我的民哪，我開你們的墳墓，使你們從墳墓中出來，你們就知道我是耶和華。我必將\textcolor{red}{我的靈}放在你們裡面，你們就要活了。我將你們安置在本地，你們就知道我耶和華如此說，也如此成就了。這是耶和華說的。』」\\  \midrule
結43:4-5&耶和華的榮光從朝東的門照入殿中。\textcolor{red}{靈}將我舉起，帶入內院，不料，耶和華的榮光充滿了殿。\\  \midrule

\end{tabular}
\end{table}


%\subsubsection{以西结书}
\begin{table}[H]
\centering
\begin{tabular}{p{5em}p{45em}}  \toprule
经节&经文 \\  \midrule

結1:4-21&我觀看,見狂風從北方颳來,隨著有一朵包括閃爍火的大雲,周圍有光輝,從其中的火內發出好像光耀的精金。又從其中顯出四個活物的形象來,他們的形狀是這樣:有人的形象,各有四個臉面、四個翅膀。他們的腿是直的,腳掌好像牛犢之蹄,都燦爛如光明的銅。在四面的翅膀以下有人的手。這四個活物的臉和翅膀乃是這樣:翅膀彼此相接,行走並不轉身,俱各直往前行。至於臉的形象,前面各有人的臉,右面各有獅子的臉,左面各有牛的臉,後面各有鷹的臉。各展開上邊的兩個翅膀相接,各以下邊的兩個翅膀遮體。他們俱各直往前行,\textcolor{red}{靈}往哪裡去,他們就往哪裡去,行走並不轉身。至於四活物的形象,就如燒著火炭的形狀,又如火把的形狀。火在四活物中間上去下來,這火有光輝,從火中發出閃電。這活物往來奔走,好像電光一閃。我正觀看活物的時候,見活物的臉旁各有一輪在地上。輪的形狀和顏色好像水蒼玉。四輪都是一個樣式,形狀和做法好像輪中套輪。輪行走的時候,向四方都能直行,並不掉轉。至於輪輞,高而可畏,四個輪輞周圍滿有眼睛。活物行走,輪也在旁邊行走;活物從地上升,輪也都上升。\textcolor{red}{靈}往哪裡去,活物就往哪裡去;活物上升,輪也在活物旁邊上升，因為活物的\textcolor{red}{靈}在輪中。那些行走,這些也行走;那些站住,這些也站住;那些從地上升,輪也在旁邊上升,因為活物的\textcolor{red}{靈}在輪中。 \\  \midrule
結10:15-17&基路伯升上去了，這是我在迦巴魯河邊所見的活物。基路伯行走，輪也在旁邊行走；基路伯展開翅膀，離地上升，輪也不轉離他們旁邊。那些站住，這些也站住，那些上升，這些也一同上升，因為活物的\textcolor{red}{靈}在輪中。\\  \midrule



結11:19-20&我要使他們有合一的心，也要將\textcolor{red}{新靈}放在他們裡面，又從他們肉體中除掉石心，賜給他們肉心， 20 使他們順從我的律例，謹守遵行我的典章。他們要做我的子民，我要做他們的神。 \\  \midrule




結18:31-32&你們要將所犯的一切罪過盡行拋棄，自做一個新心和\textcolor{red}{新靈}。以色列家啊，你們何必死亡呢？ 32 主耶和華說：我不喜悅那死人之死，所以你們當回頭而存活。\\  \midrule

結36:25-28&我必用清水灑在你們身上，你們就潔淨了。我要潔淨你們，使你們脫離一切的汙穢，棄掉一切的偶像。我也要賜給你們一個新心，將\textcolor{red}{新靈}放在你們裡面，又從你們的肉體中除掉石心，賜給你們肉心。 27 我必將\textcolor{red}{我的靈}放在你們裡面，使你們順從我的律例，謹守遵行我的典章。 28 你們必住在我所賜給你們列祖之地，你們要做我的子民，我要做你們的神。\\  \midrule

結39:25-29&主耶和華如此說：我要使雅各被擄的人歸回，...我必不再留他們一人在外邦。我也不再掩面不顧他們，因我已將\textcolor{red}{我的靈}澆灌以色列家。這是主耶和華說的。\\  \midrule

\end{tabular}
\end{table}

\subsubsection{給人有聰明智慧-做圣工}
\begin{table}[H]
\centering
\begin{tabular}{p{5em}p{45em}}  \toprule
经节&经文 \\  \midrule

出28:2-4&你要給你哥哥亞倫做聖衣為榮耀，為華美。又要吩咐一切心中有智慧的，就是我用智慧的靈所充滿的，給亞倫做衣服，使他分別為聖，可以給我供祭司的職分。所要做的就是胸牌、以弗得、外袍、雜色的內袍、冠冕、腰帶，使你哥哥亞倫和他兒子穿這聖服，可以給我供祭司的職分。\\  \midrule
出31:3-6&看哪，猶大支派中，戶珥的孫子、烏利的兒子比撒列，我已經提他的名召他。 3 我也以我的靈充滿了他，使他有智慧，有聰明，有知識，能做各樣的工， 4 能想出巧工，用金、銀、銅製造各物， 5 又能刻寶石，可以鑲嵌，能雕刻木頭，能做各樣的工。 6 我分派但支派中亞希撒抹的兒子亞何利亞伯與他同工。凡心裡有智慧的，我更使他們有智慧，能做我一切所吩咐的，\\  \midrule
出35:20-29&以色列全會眾從摩西面前退去。凡心裡受感和甘心樂意的,都拿耶和華的禮物來,用以做會幕和其中一切的使用,又用以做聖衣。凡心裡樂意獻禮物的,連男帶女，各將金器,就是胸前針、耳環 、打印的戒指和手釧,帶來獻給耶和華。凡有藍色、紫色、朱紅色線,細麻,山羊毛,染紅的公羊皮,海狗皮的,都拿了來。凡獻銀子和銅給耶和華為禮物的,都拿了來。凡有皂莢木可做什麼使用的，也拿了來。凡心中有智慧的婦女，親手紡線，把所紡的藍色、紫色、朱紅色線和細麻都拿了來。凡有智慧心裡受感的婦女，就紡山羊毛。眾官長把紅瑪瑙和別樣的寶石，可以鑲嵌在以弗得與胸牌上的，都拿了來，又拿香料做香，拿油點燈,做膏油。以色列人，無論男女，凡甘心樂意獻禮物給耶和華的，都將禮物拿來，做耶和華藉摩西所吩咐的一切工。\\  \midrule

出35:30-35&摩西對以色列人說：「猶大支派中，戶珥的孫子、烏利的兒子比撒列，耶和華已經提他的名召他， 31 又以神的靈充滿了他，使他有智慧、聰明、知識，能做各樣的工， 32 能想出巧工，用金、銀、銅製造各物， 33 又能刻寶石，可以鑲嵌，能雕刻木頭，能做各樣的巧工。 34 耶和華又使他和但支派中亞希撒抹的兒子亞何利亞伯心裡靈明，能教導人。 35 耶和華使他們的心滿有智慧，能做各樣的工，無論是雕刻的工，巧匠的工，用藍色、紫色、朱紅色線和細麻繡花的工並機匠的工，他們都能做，也能想出奇巧的工\\  \midrule

\end{tabular}
\end{table}



\subsubsection{说预言的靈,劝慰引导人的灵}
\begin{table}[H]
\centering
\begin{tabular}{p{5em}p{43em}}  \toprule
经节&经文 \\  \midrule
民11:16-17&摩西出去，將耶和華的話告訴百姓，又招聚百姓的長老中七十個人來，使他們站在會幕的四圍。 25 耶和華在雲中降臨，對摩西說話，把降於他身上的靈分賜那七十個長老。靈停在他們身上的時候，他們就受感說話，以後卻沒有再說。\\  \midrule

民11:24-30&摩西出去，將耶和華的話告訴百姓，又招聚百姓的長老中七十個人來，使他們站在會幕的四圍。 25 耶和華在雲中降臨，對摩西說話，把降於他身上的靈分賜那七十個長老。靈停在他們身上的時候，他們就受感說話，以後卻沒有再說。但有兩個人仍在營裡，一個名叫伊利達，一個名叫米達。他們本是在那些被錄的人中，卻沒有到會幕那裡去。靈停在他們身上，他們就在營裡說預言。 27 有個少年人跑來告訴摩西說：「伊利達、米達在營裡說預言。」 28 摩西的幫手，嫩的兒子約書亞，就是摩西所揀選的一個人，說：「請我主摩西禁止他們！」 29 摩西對他說：「你為我的緣故嫉妒人嗎？唯願耶和華的百姓都受感說話，願耶和華把他的靈降在他們身上！」 30 於是摩西和以色列的長老都回到營裡去。\\  \midrule

民24:1-3&巴蘭見耶和華喜歡賜福於以色列，就不像前兩次去求法術，卻面向曠野。 2 巴蘭舉目，看見以色列人照著支派居住，神的靈就臨到他身上， 3 他便題起詩歌說：\\  \midrule
代下15:1-7&神的靈感動俄德的兒子亞撒利雅， 2 他出來迎接亞撒，對他說：「亞撒和猶大、便雅憫眾人哪，要聽我說！你們若順從耶和華，耶和華必與你們同在。你們若尋求他，就必尋見；你們若離棄他，他必離棄你們。 3 以色列人不信真神，沒有訓誨的祭司，也沒有律法，已經好久了， 4 但他們在急難的時候歸向耶和華以色列的神，尋求他，他就被他們尋見。 5 那時，出入的人不得平安，列國的居民都遭大亂， 6 這國攻擊那國，這城攻擊那城，互相破壞，因為神用各樣災難擾亂他們。 7 現在你們要剛強，不要手軟，因你們所行的必得賞賜。」\\  \midrule

代下20:14-19&那時，耶和華的靈在會中臨到利未人亞薩的後裔瑪探雅的玄孫、耶利的曾孫、比拿雅的孫子、撒迦利雅的兒子雅哈悉， 15 他說：「猶大眾人、耶路撒冷的居民和約沙法王，你們請聽。耶和華對你們如此說：不要因這大軍恐懼驚惶，因為勝敗不在乎你們，乃在乎神。 16 明日你們要下去迎敵。他們是從洗斯坡上來，你們必在耶魯伊勒曠野前的谷口遇見他們。 17 猶大和耶路撒冷人哪，這次你們不要爭戰，要擺陣站著，看耶和華為你們施行拯救。不要恐懼，也不要驚惶，明日當出去迎敵，因為耶和華與你們同在。」 18 約沙法就面伏於地，猶大眾人和耶路撒冷的居民也俯伏在耶和華面前，叩拜耶和華。 19 哥轄族和可拉族的利未人都起來，用極大的聲音讚美耶和華以色列的神。\\  \midrule

代下24:20-22&那時，神的靈感動祭司耶何耶大的兒子撒迦利亞，他就站在上面對民說：「神如此說：你們為何干犯耶和華的誡命，以致不得亨通呢？因為你們離棄耶和華，所以他也離棄你們。」 21 眾民同心謀害撒迦利亞，就照王的吩咐，在耶和華殿的院內用石頭打死他。 22 這樣，約阿施王不想念撒迦利亞的父親耶何耶大向自己所施的恩，殺了他的兒子。撒迦利亞臨死的時候說：「願耶和華鑒察申冤！」\\  \midrule

撒7:11-12&他們卻不肯聽從，扭轉肩頭，塞耳不聽， 12 使心硬如金鋼石，不聽律法和萬軍之耶和華用靈藉從前的先知所說的話。故此，萬軍之耶和華大發烈怒。\\  \midrule

尼9:20-21&你也賜下你良善的靈教訓他們，未嘗不賜嗎哪使他們糊口，並賜水解他們的渴。 21 在曠野四十年，你養育他們，他們就一無所缺，衣服沒有穿破，腳也沒有腫。\\  \midrule
尼9:30-31&但你多年寬容他們,又用你的靈藉眾先知勸誡他們,他們仍不聽從,所以你將他們交在列國之民的手中。然而你大發憐憫,不全然滅絕他們,也不丟棄他們,因為你是有恩典有憐憫的神。\\  \midrule

彌3:8&至於我，我藉耶和華的靈，滿有力量、公平、才能，可以向雅各說明他的過犯，向以色列指出他的罪惡。\\  \midrule
伯32:6-8&布西人巴拉迦的兒子以利戶回答說：「我年輕,你們老邁,因此我退讓,不敢向你們陳說我的意見。我說,年老的當先說話,壽高的當以智慧教訓人。但在人裡面有靈,全能者的氣使人有聰明。\\  \midrule

詩51:9-12&求你掩面不看我的罪，塗抹我一切的罪孽。神啊，求你為我造清潔的心，使我裡面重新有正直的靈。11 不要丟棄我，使我離開你的面，不要從我收回你的聖靈。12 求你使我仍得救恩之樂，賜我樂意的靈扶持我。
\\\bottomrule
\end{tabular}
\end{table}




\subsubsection{亞瑪撒,掃羅,大衛}
\begin{table}[H]
\centering
\begin{tabular}{p{6em}p{40em}}  \toprule
经节&经文 \\  \midrule
代上12:18&那時\textcolor{red}{神的靈感動}那三十個勇士的首領亞瑪撒，他就說：「大衛啊，我們是歸於你的；耶西的兒子啊，我們是幫助你的！願你平平安安，願幫助你的也都平安，因為你的神幫助你！」大衛就收留他們，立他們做軍長。 \\  \midrule

撒上10:5-7&此後你到神的山，在那裡有非利士人的防兵。你到了城的時候，必遇見一班先知從丘壇下來，前面有鼓瑟的、擊鼓的、吹笛的、彈琴的，他們都受感說話。\textcolor{red}{耶和華的靈必大大感動你}，你就與他們一同受感說話，你要變為新人。 7 這兆頭臨到你，你就可以趁時而做，因為神與你同在。\\  \midrule
撒上10:9-13&掃羅轉身離別撒母耳，神就賜他一個新心。當日這一切兆頭都應驗了。 10 掃羅到了那山，有一班先知遇見他，\textcolor{red}{神的靈大大感動他}，他就在先知中受感說話。 11 素來認識掃羅的，看見他和先知一同受感說話，就彼此說：「基士的兒子遇見什麼了？掃羅也列在先知中嗎？」 12 那地方有一個人說：「這些人的父親是誰呢？」此後有句俗語說：「掃羅也列在先知中嗎？」 13 掃羅受感說話已畢，就上丘壇去了。\\  \midrule

撒上11:6-7&掃羅聽見這話，就\textcolor{red}{被神的靈大大感動}，甚是發怒。 7 他將一對牛切成塊子，託付使者傳送以色列的全境，說：「凡不出來跟隨掃羅和撒母耳的，也必這樣切開他的牛。」於是耶和華使百姓懼怕，他們就都出來，如同一人。\\  \midrule
撒上16:12-16&耶西就打發人去叫了他來。他面色光紅，雙目清秀，容貌俊美。耶和華說：「這就是他，你起來膏他。」 13 撒母耳就用角裡的膏油，在他諸兄中膏了他。從這日起，\textcolor{red}{耶和華的靈就大大感動大衛}。撒母耳起身回拉瑪去了。\textcolor{red}{耶和華的靈}離開掃羅，有惡魔從耶和華那裡來擾亂他。 15 掃羅的臣僕對他說：「現在有惡魔從神那裡來擾亂你。 16 我們的主可以吩咐面前的臣僕，找一個善於彈琴的來，等神那裡來的惡魔臨到你身上的時候，使他用手彈琴，你就好了。」\\  \midrule

撒上19:19-24&有人告訴掃羅，說大衛在拉瑪的拿約。 20 掃羅打發人去捉拿大衛。去的人見有一班先知都受感說話，撒母耳站在其中監管他們。打發去的人也\textcolor{red}{受神的靈感動說話}。 21 有人將這事告訴掃羅，他又打發人去，他們也受感說話。掃羅第三次打發人去，他們也受感說話。 22 然後掃羅自己往拉瑪去，到了西沽的大井，問人說：「撒母耳和大衛在哪裡呢？」有人說：「在拉瑪的拿約。」 23 他就往拉瑪的拿約去。神的靈也感動他，一面走一面說話，直到拉瑪的拿約。 24 他就脫了衣服，在撒母耳面前受感說話，一晝一夜露體躺臥。因此有句俗語說：「掃羅也列在先知中嗎？」\\  \midrule

%撒上16:12-13&耶西就打發人去叫了他來。他面色光紅，雙目清秀，容貌俊美。耶和華說：「這就是他，你起來膏他。」撒母耳就用角裡的膏油，在他諸兄中膏了他。從這日起，\textcolor{red}{耶和華的靈就大大感動大衛。撒母耳起身回拉瑪去了。\\  \midrule

撒下23:1-7&以下是大衛末了的話。耶西的兒子大衛得居高位，是雅各神所膏的，做以色列的美歌者，說： 2 「\textcolor{red}{耶和華的靈藉著我說}，他的話在我口中。 3 以色列的神，以色列的磐石曉諭我說：那以公義治理人民的，敬畏神執掌權柄， 4 他必像日出的晨光，如無雲的清晨，雨後的晴光，使地發生嫩草。』 5 我家在神面前並非如此，神卻與我立永遠的約。這約凡事堅穩，關乎我的一切救恩和我一切所想望的，他豈不為我成就嗎？ 6 但匪類都必像荊棘被丟棄，人不敢用手拿它。 7 拿它的人必帶鐵器和槍桿，終久它必被火焚燒。」\\  \midrule

詩51:9-12&(大衛)求你掩面不看我的罪，塗抹我一切的罪孽。神啊，求你為我造清潔的心，使我裡面\textcolor{red}{重新有正直的靈}。11 不要丟棄我，使我離開你的面，不要從我收回\textcolor{red}{你的聖靈}。12 求你使我仍得救恩之樂，\textcolor{red}{賜我樂意的靈}扶持我。\\  \midrule
代上28:12&大衛將殿的遊廊、旁屋、府庫、樓房、內殿和施恩所的樣式指示他兒子所羅門。 12 又將\textcolor{red}{被靈感動}所得的樣式，就是耶和華神殿的院子、周圍的房屋、殿的府庫和聖物府庫的一切樣式，都指示他。 13 又指示他祭司和利未人的班次，與耶和華殿裡各樣的工作，並耶和華殿裡一切器皿的樣式， 14 以及各樣應用金器的分量和各樣應用銀器的分量， 15 金燈臺和金燈的分量，銀燈臺和銀燈的分量（輕重各都合宜）， 16 陳設餅金桌子的分量，銀桌子的分量， 17 精金的肉叉子、盤子和爵的分量，各金碗與各銀碗的分量， 18 精金香壇的分量，並用金子做基路伯，基路伯張開翅膀，遮掩耶和華的約櫃。 19 大衛說：「這一切工作的樣式，都是耶和華用手劃出來使我明白的。」
\\\bottomrule
\end{tabular}  
\end{table}





\subsubsection{住在百姓中間/引导/擔憂/使人得安息}
\begin{table}[H]
\centering
\begin{tabular}{p{5em}|p{43em}}  \toprule
经文&内容\\  \midrule
出33:14&耶和華說：「我必親自和你同去，使你得安息。」\\  \midrule

該2:4-5&耶和華說：所羅巴伯啊，雖然如此，你當剛強！約薩答的兒子大祭司約書亞啊，你也當剛強！這地的百姓，你們都當剛強做工！因為我與你們同在。這是萬軍之耶和華說的。這是照著你們出埃及我與你們立約的話，那時我的靈住在你們中間。你們不要懼怕！\\  \midrule
詩139:7&我往哪裡去躲避你的靈？我往哪裡逃躲避你的面？\\  \midrule
詩143:10&求你指教我遵行你的旨意，因你是我的神。你的靈本為善，求你引我到平坦之地。\\  \midrule

賽63:8-14&他們在一切苦難中,他也同受苦難,並且他面前的使者拯救他們。他以慈愛和憐憫救贖他們。在古時的日子,常保抱他們,懷揣他們。他們竟悖逆,使\textcolor{red}{主的聖靈}擔憂。他就轉做他們的仇敵,親自攻擊他們。那時他們想起古時的日子摩西和他百姓,說：「將百姓和牧養他全群的人從海裡領上來的在哪裡呢?將\textcolor{red}{他的聖靈}降在他們中間的在哪裡呢?使他榮耀的膀臂在摩西的右手邊行動,在他們前面將水分開，要建立自己永遠的名,帶領他們經過深處,如馬行走曠野,使他們不致絆跌的在哪裡呢?」 \textcolor{red}{耶和華的靈}使他們得安息,彷彿牲畜下到山谷,照樣,你也引導你的百姓,要建立自己榮耀的名。\\\bottomrule
\end{tabular}
\end{table}




\subsubsection{其他的灵-惡魔,沉睡的靈，恐懼的靈，謊言的靈}
\begin{table}[H]
\centering
\begin{tabular}{p{6em}p{42em}}  \toprule
经节&经文 \\  \midrule
撒上16:14-16&耶和華的靈離開掃羅，有惡魔從耶和華那裡來擾亂他。掃羅的臣僕對他說：「現在有惡魔從神那裡來擾亂你。我們的主可以吩咐面前的臣僕，找一個善於彈琴的來，等神那裡來的惡魔臨到你身上的時候，使他用手彈琴，你就好了。」 \\  \midrule

撒上16:23&從神那裡來的惡魔臨到掃羅身上的時候,大衛就拿琴,用手而彈,掃羅便舒暢爽快,惡魔離了他。 \\  \midrule

撒上18:10-12&次日，從神那裡來的惡魔大大降在掃羅身上，他就在家中胡言亂語。大衛照常彈琴，掃羅手裡拿著槍。掃羅把槍一掄，心裡說：「我要將大衛刺透，釘在牆上。」大衛躲避他兩次。掃羅懼怕大衛，因為耶和華離開自己，與大衛同在\\  \midrule

撒上19:9-10&從耶和華那裡來的惡魔又降在掃羅身上（掃羅手裡拿槍坐在屋裡），大衛就用手彈琴。掃羅用槍想要刺透大衛，釘在牆上。他卻躲開，掃羅的槍刺入牆內。當夜大衛逃走，躲避了。 \\  \midrule


賽29:9-10&你們等候驚奇吧!你們宴樂昏迷吧!他們醉了,卻非因酒;他們東倒西歪,卻非因濃酒.因為耶和華將沉睡的靈澆灌你們,封閉你們的眼,蒙蓋你們的頭.你們的眼就是先知,你們的頭就是先見.\\  \midrule
賽19:13-15&瑣安的首領都變為愚昧，挪弗的首領都受了迷惑，當埃及支派房角石的使埃及人走錯了路。耶和華使乖謬的靈摻入埃及中間，首領使埃及一切所做的都有差錯，好像醉酒之人嘔吐的時候東倒西歪一樣。埃及中，無論是頭與尾，棕枝與蘆葦，所做之工都不成就。\\  \midrule
伯4:12-21&提幔人以利法回答說：...我暗暗地得了默示，我耳朵也聽其細微的聲音。 13 在思念夜中異象之間，世人沉睡的時候， 14 恐懼、戰兢臨到我身，使我百骨打戰。 15 有靈從我面前經過，我身上的毫毛直立。 16 那靈停住，我卻不能辨其形狀，有影像在我眼前，我在靜默中聽見有聲音說： 17 『必死的人豈能比神公義嗎？人豈能比造他的主潔淨嗎？ 18 主不信靠他的臣僕，並且指他的使者為愚昧， 19 何況那住在土房，根基在塵土裡，被蠹蟲所毀壞的人呢？ 20 早晚之間，就被毀滅，永歸無有，無人理會。 21 他帳篷的繩索豈不從中抽出來呢？他死，且是無智慧而死。』\\  \midrule

王上22:21-23&隨後有一個神靈出來,站在耶和華面前說:『我去引誘他。』耶和華問他說：『你用何法呢？』他說:『我去,要在他眾先知口中做謊言的靈。』耶和華說:『這樣,你必能引誘他,你去如此行吧。』現在耶和華使謊言的靈入了你這些先知的口,並且耶和華已經命定降禍於你。」\\  \midrule

代下18:20-22&隨後有一個神靈出來,站在耶和華面前說：『我去引誘他。』耶和華問他說：『你用何法呢？』他說：『我去，要在他眾先知口中做謊言的靈。』耶和華說：『這樣，你必能引誘他，你去如此行吧。』現在耶和華使謊言的靈入了你這些先知的口，並且耶和華已經命定降禍於你。」\\  \midrule
賽30:1-2&耶和華說:「禍哉,這悖逆的兒女!他們同謀卻不由於我,結盟卻不由於我的靈,以致罪上加罪。\\  \midrule
\end{tabular}
\end{table}


\subsubsection{發怒}
\begin{table}[H]
\centering
\begin{tabular}{p{6em}p{40em}}  \toprule
经节&经文 \\  \midrule
出15:8&你大發威嚴，推翻那些起來攻擊你的；你發出烈怒如火，燒滅他們像燒碎秸一樣。 8 你發鼻中的氣，水便聚起成堆，大水直立如壘，海中的深水凝結。\\  \midrule
民16:20-24&耶和華曉諭摩西、亞倫說： 21 「你們離開這會眾，我好在轉眼之間把他們滅絕。」 22 摩西、亞倫就俯伏在地，說：「神，萬人之靈的神啊！一人犯罪，你就要向全會眾發怒嗎？」 23 耶和華曉諭摩西說： 24 「你吩咐會眾說：你們離開可拉、大坍、亞比蘭帳篷的四圍。」\\  \midrule

撒下22:16&耶和華的斥責一發，鼻孔的氣一出，海底就出現，大地的根基也顯露。 \\  \midrule
詩18:15-17&耶和華啊，你的斥責一發，你鼻孔的氣一出，海底就出現，大地的根基也顯露。…\\  \midrule

伯4:9&提幔人以利法回答說：...神一出氣，他們就滅亡；神一發怒，他們就消沒。\\  \midrule
伯15:30&惡人一生之日劬勞痛苦，強暴人一生的年數也是如此。.....他不得富足，財物不得常存，產業在地上也不加增。他不得出離黑暗，火焰要將他的枝子燒乾。因\textcolor{red}{神口中的氣}，他要滅亡。\\  \midrule
賽11:1-5&從耶西的本 a必發一條，從他根生的枝子必結果實。 2 耶和華的靈必住在他身上，就是使他有智慧和聰明的靈，謀略和能力的靈，知識和敬畏耶和華的靈。 3 他必以敬畏耶和華為樂，行審判不憑眼見，斷是非也不憑耳聞； 4 卻要以公義審判貧窮人，以正直判斷世上的謙卑人，以口中的杖擊打世界，\textcolor{red}{以嘴裡的氣}殺戮惡人。 5 公義必當他的腰帶，信實必當他脅下的帶子。\\  \midrule

\end{tabular}
\end{table}

\subsection{聖靈在新约的工作}
\subsubsection{聖靈的内在工作——更新改变}
\begin{table}[H]
\centering
\begin{tabular}{p{6em}|p{4em}|p{37em}}  \toprule
经节&工作&内容 \\  \midrule
約16:8-11&責備&他既來了,就要叫世人為罪、為義、為審判,自己責備自己。為罪,是因他們不信我;為義,是因我往父那裡去,你們就不再見我;為審判,是因這世界的王受了審判。\\  \midrule

約3:6-8&重生&從肉身生的就是肉身，從靈生的就是靈。我說你們必須重生，你不要以為稀奇。風隨著意思吹，你聽見風的響聲，卻不曉得從哪裡來，往哪裡去。凡從聖靈生的，也是如此。\\  \midrule
林前12:3&信主&被神的靈感動的，沒有說「耶穌是可咒詛」的；若不是被聖靈感動的，也沒有能說「耶穌是主」的。\\  \midrule
多3:5-7&更新&他便救了我們，並不是因我們自己所行的義，乃是照他的憐憫，藉著重生的洗和聖靈的更新。聖靈就是神藉著耶穌基督——我們救主厚厚澆灌在我們身上的，好叫我們因他的恩得稱為義，可以憑著永生的盼望成為後嗣。\\  \midrule
弗4:30&擔憂&不要叫神的聖靈擔憂，你們原是受了他的印記，等候得贖的日子來到。\\  \midrule

約14:15-17&同在&你們若愛我，就必遵守我的命令。「我要求父，父就另外賜給你們一位保惠師，叫他永遠與你們同在，就是真理的聖靈，乃世人不能接受的，因為不見他，也不認識他。你們卻認識他，因他常與你們同在，也要在你們裡面。我不撇下你們為孤兒，我必到你們這裡來。\\  \midrule
%&住在信徒心里的\\  \midrule
%罗8:9-11&住在信徒心里&\\  \midrule

林前6:19&聖靈的殿&豈不知你們的身子就是聖靈的殿嗎？這聖靈是從神而來，住在你們裡頭的。並且你們不是自己的人，\\  \midrule
弗2:19-22&聖靈居所&你們不再做外人和客旅，是與聖徒同國，是神家裡的人了。 20 並且被建造在使徒和先知的根基上，有基督耶穌自己為房角石，各房靠他聯絡得合式，漸漸成為主的聖殿，你們也靠他同被建造，成為神藉著聖靈居住的所在。\\  \midrule



羅8:1-4&治死惡行&如今，那些在基督耶穌裡的就不定罪了。因為賜生命聖靈的律在基督耶穌裡釋放了我，使我脫離罪和死的律了。律法既因肉體軟弱，有所不能行的，神就差遣自己的兒子成為罪身的形狀，做了贖罪祭，在肉體中定了罪案， 4 使律法的義成就在我們這不隨從肉體、只隨從聖靈的人身上。\\  \midrule

羅8:5-11&治死惡行&因為隨從肉體的人,體貼肉體的事；隨從聖靈的人,體貼聖靈的事。體貼肉體的就是死,體貼聖靈的乃是生命、平安。原來體貼肉體的就是與神為仇,因為不服神的律法,也是不能服；而且屬肉體的人不能得神的喜歡。如果神的靈住在你們心裡,你們就不屬肉體，乃屬聖靈了。人若沒有基督的靈,就不是屬基督的。基督若在你們心裡,身體就因罪而死,心靈卻因義而活。然而叫耶穌從死裡復活者的靈若住在你們心裡,那叫基督耶穌從死裡復活的,也必藉著住在你們心裡的聖靈,使你們必死的身體又活過來。\\  \midrule

羅8:13&治死惡行&你們若順從肉體活著,必要死；若靠著聖靈治死身體的惡行,必要活著。\\ 
 \midrule


羅8:26-27&替人禱告&況且，我們的軟弱有聖靈幫助。我們本不曉得當怎樣禱告，只是聖靈親自用說不出來的嘆息替我們禱告。鑒察人心的，曉得聖靈的意思，因為聖靈照著神的旨意替聖徒祈求。\\  \midrule


%押金/印/凭据&\\  \midrule
林后1:21-23&印記&那在基督裡堅固我們和你們，並且膏我們的就是神；他又用印印了我們，並賜聖靈在我們心裡做憑據。我呼籲神給我的心作見證，\\  \midrule
弗1:13-14&憑據印記&你們既聽見真理的道，就是那叫你們得救的福音，也信了基督，既然信他，就受了所應許的聖靈為印記。這聖靈是我們得基業的憑據，直等到神之民被贖，使他的榮耀得著稱讚。\\  \midrule
弗4:30&憑據印記&不要叫神的聖靈擔憂，你們原是受了他的印記，等候得贖的日子來到。\\  \midrule
林后5:4-5&憑據&我們在這帳篷裡嘆息勞苦，並非願意脫下這個，乃是願意穿上那個，好叫這必死的被生命吞滅了。為此培植我們的就是神，他又賜給我們聖靈做憑據。\\  \midrule

\end{tabular}%
%  \label{tab:addlabel}%
% \caption{法利賽人求神蹟}  
\end{table}



\begin{table}[H]
\centering
\begin{tabular}{p{6em}|p{4em}|p{37em}}  \toprule
经节&工作&内容 \\  \midrule

約16:12-15&啟示&我還有好些事要告訴你們，但你們現在擔當不了。只等真理的聖靈來了，他要引導你們明白一切的真理，因為他不是憑自己說的，乃是把他所聽見的都說出來，並要把將來的事告訴你們。他要榮耀我，因為他要將受於我的告訴你們。凡父所有的，都是我的，所以我說，他要將受於我的告訴你們。\\  \midrule
弗3:5&啟示&這奧祕在以前的世代沒有叫人知道，像如今藉著聖靈啟示他的聖使
徒和先知一樣。\\  \midrule

林前2:10&啟示&只有神藉着圣灵向我们显明了，因为圣灵参透万事，就是神深奥的事也参透了。\\  \midrule

路2:26-36&&在耶路撒冷有一個人名叫西面，這人又公義又虔誠，素常盼望以色列的安慰者來到，又有聖靈在他身上。 26 他得了聖靈的啟示，知道自己未死以前，必看見主所立的基督。 27 他受了聖靈的感動，進入聖殿，正遇見耶穌的父母抱著孩子進來，要照律法的規矩辦理。 28 西面就用手接過他來，稱頌神，說： 29 「主啊，如今可以照你的話，釋放僕人安然去世！ 30 因為我的眼睛已經看見你的救恩， 31 就是你在萬民面前所預備的， 32 是照亮外邦人的光，又是你民以色列的榮耀。」 33 孩子的父母因這論耶穌的話就稀奇。 34 西面給他們祝福，又對孩子的母親馬利亞說：「這孩子被立，是要叫以色列中許多人跌倒、許多人興起，又要做毀謗的話柄， 35 叫許多人心裡的意念顯露出來。你自己的心也要被刀刺透。」\\  \midrule

太12:27&趕鬼&我若是靠著別西卜趕鬼，你們的子弟又是靠著誰趕呢？這樣，他們就要審判你們了。\\  \midrule

約14:26&教訓&但保惠師，就是父因我的名所要差來的聖靈，他要將一切的事指教你們，並且要叫你們想起我對你們所說的一切話。\\  \midrule
約一2:27&教訓&你們從主所受的恩膏常存在你們心裡，並不用人教訓你們，自有主的恩膏在凡事上教訓你們。這恩膏是真的，不是假的。你們要按這恩膏的教訓住在主裡面。\\  \midrule
林前2:6-16&指教&然而，在完全的人中，我們也講智慧，但不是這世上的智慧，也不是這世上有權有位將要敗亡之人的智慧。 7 我們講的，乃是從前所隱藏、神奧祕的智慧，就是神在萬世以前預定使我們得榮耀的。 8 這智慧，世上有權有位的人沒有一個知道的，他們若知道，就不把榮耀的主釘在十字架上了。 9 如經上所記：「神為愛他的人所預備的，是眼睛未曾看見，耳朵未曾聽見，人心也未曾想到的。」 10 只有神藉著聖靈向我們顯明了。因為聖靈參透萬事，就是神深奧的事也參透了。 11 除了在人裡頭的靈，誰知道人的事？像這樣，除了神的靈，也沒有人知道神的事。 12 我們所領受的，並不是世上的靈，乃是從神來的靈，叫我們能知道神開恩賜給我們的事。 13 並且我們講說這些事，不是用人智慧所指教的言語，乃是用聖靈所指教的言語，將屬靈的話解釋屬靈的事 。然而，屬血氣的人不領會神聖靈的事，反倒以為愚拙；並且不能知道，因為這些事唯有屬靈的人才能看透。 15 屬靈的人能看透萬事，卻沒有一人能看透了他。 16 「誰曾知道主的心，去教導他呢？」但我們是有基督的心了。\\  \midrule

加5:22-24&結果子&聖靈所結的果子，就是仁愛、喜樂、和平、忍耐、恩慈、良善、信實、 23 溫柔、節制；這樣的事沒有律法禁止。 24 凡屬基督耶穌的人，是已經把肉體連肉體的邪情私慾同釘在十字架上了。\\  \midrule

徒9:31&安慰人&那時，猶太、加利利、撒馬利亞各處的教會都得平安，被建立，凡事敬畏主，蒙聖靈的安慰，人數就增多了。\\  \midrule

彼前1：2&成聖&就是照父神的先見被揀選，藉著聖靈得成聖潔，以致順服耶穌基督，又蒙他血所灑的人。願恩惠、平安多多地加給你們！\\  \midrule

%&神的灵/主/基督\\  \midrule

林后3:17&赐自由&主就是那靈，主的靈在哪裡，哪裡就得以自由。\\  \midrule

羅15:13&使人有盼望&但願使人有盼望的神，因信將諸般的喜樂、平安充滿你們的心，使你們藉著聖靈的能力大有盼望！\\ \bottomrule

%见证者&\\  \midrule
罗8:14-16&见证&因為凡被神的靈引導的，都是神的兒子。你們所受的不是奴僕的心，仍舊害怕；所受的乃是兒子的心，因此我們呼叫：「阿爸！父！」聖靈與我們的心同證我們是神的兒女。\\  \midrule
来2:1-4&&所以，我們當越發鄭重所聽見的道理，恐怕我們隨流失去。那藉著天使所傳的話既是確定的，凡干犯悖逆的都受了該受的報應，我們若忽略這麼大的救恩，怎能逃罪呢？這救恩起先是主親自講的，後來是聽見的人給我們證實了，神又按自己的旨意，用神蹟、奇事和百般的異能並聖靈的恩賜，同他們作見證。\\  \midrule
来10:15-16&&聖靈也對我們作見證，因為他既已說過：「主說：『那些日子以後，我與他們所立的約乃是這樣：我要將我的律法寫在他們心上，又要放在他們的裡面』」\\  \midrule
彼前1:10-12&&論到這救恩，那預先說你們要得恩典的眾先知早已詳細地尋求考察，就是考察在他們心裡基督的靈，預先證明基督受苦難、後來得榮耀是指著什麼時候，並怎樣的時候。他們得了啟示，知道他們所傳講的一切事，不是為自己，乃是為你們。那靠著從天上差來的聖靈傳福音給你們的人，現在將這些事報給你們；天使也願意詳細察看這些事。 \\  \midrule
\end{tabular}%
%  \label{tab:addlabel}%
% \caption{法利賽人求神蹟}  
\end{table}


\subsubsection{聖靈的外在工作——服事的能力}
\begin{table}[H]
\centering
\begin{tabular}{p{6em}|p{4em}|p{37em}}  \toprule
经节&工作&内容 \\  \midrule
約15:26&作見證&但我要從父那裡差保惠師來，就是從父出來真理的聖靈；他來了，就要為我作見證。 \\  \midrule

彼前1:10-12&&論到這救恩，那預先說你們要得恩典的眾先知早已詳細地尋求考察，就是考察在他們心裡基督的靈，預先證明基督受苦難、後來得榮耀是指著什麼時候，並怎樣的時候。他們得了啟示，知道他們所傳講的一切事，不是為自己，乃是為你們。那靠著從天上差來的聖靈傳福音給你們的人，現在將這些事報給你們；天使也願意詳細察看這些事。\\  \midrule



徒1:8&使人为主作见证&但圣灵降临在你们身上，你们就必得着能力。并要在耶路撒冷，犹大全地，和撒玛利亚，直到地极，作我的见证。\\  \midrule

徒6:10&做见证&司提反是以智慧和聖靈說話，眾人抵擋不住。\\  \midrule

徒15:7-8&做见证&你們知道神早已在你們中間揀選了我(彼得)，叫外邦人從我口中得聽福音之道，而且相信知道人心的神也為他們作了見證，賜聖靈給他們，正如給我們一樣；\\  \midrule
徒15:28&做见证&因為聖靈和我們定意不將別的重擔放在你們身上，唯有幾件事是不可少的，\\  \midrule

路12:11-12&做见证&人帶你們到會堂，並官府和有權柄的人面前，不要思慮怎麼分訴，說什麼話， 12 因為正在那時候，聖靈要指教你們當說的話。\\  \midrule

林前12:4-11&給人恩赐&恩賜原有分別，聖靈卻是一位；職事也有分別，主卻是一位；功用也有分別，神卻是一位，在眾人裡面運行一切的事。聖靈顯在各人身上，是叫人得益處。這人蒙聖靈賜他智慧的言語，那人也蒙這位聖靈賜他知識的言語，又有一人蒙這位聖靈賜他信心，還有一人蒙這位聖靈賜他醫病的恩賜，又叫一人能行異能，又叫一人能做先知，又叫一人能辨別諸靈，又叫一人能說方言，又叫一人能翻方言。這一切都是這位聖靈所運行，隨己意分給各人的。\\  \midrule
彼後1：21&說預言&因為預言從來沒有出於人意的，乃是人被聖靈感動，說出神的話來。\\  \midrule
徒11：27-28&說預言&當那些日子，有幾位先知從耶路撒冷下到安提阿。內中有一位名叫亞迦布，站起來，藉著聖靈指明天下將有大饑荒；這事到克勞迪年間果然有了。\\  \midrule

太22:41-45 &說預言&法利賽人聚集的時候，耶穌問他們說： 「論到基督，你們的意見如何？他是誰的子孫呢？」他們回答說：「是大衛的子孫。」耶穌說：「這樣，大衛被聖靈感動，怎麼還稱他為主說： 『主對我主說：你坐在我的右邊，等我把你仇敵放在你的腳下』？大衛既稱他為主，他怎麼又是大衛的子孫呢？」\\  \midrule



徒5:3&监察人心&彼得說：「亞拿尼亞，為什麼撒旦充滿了你的心，叫你欺哄聖靈，把田地的價銀私自留下幾份呢？\\  \midrule
徒5:9&监察人心&彼得說：「你們為什麼同心試探主的靈呢？埋葬你丈夫之人的腳已到門口，他們也要把你抬出去。」\\  \midrule

徒20:28&在以弗所立監督&聖靈立你們做全群的監督，你們就當為自己謹慎，也為全群謹慎，牧養神的教會，就是他用自己血所買來的\\ \bottomrule

\end{tabular}%
%  \label{tab:addlabel}%
% \caption{法利賽人求神蹟}  
\end{table}





\subsubsection{耶穌借着圣灵做的事}
\begin{table}[H]
\centering
\begin{tabular}{p{5em}p{42em}}  \toprule
经节&经文 \\  \midrule
賽48:15-16&唯有我曾說過，我又選召他，領他來，他的道路就必亨通。 16 你們要就近我來聽這話，我從起頭並未曾在隱密處說話，自從有這事，我就在那裡。」現在，主耶和華差遣我和他的靈來（耶和華和他的靈差遣我來）。\\  \midrule
賽61:1-3&主耶和華的靈在我身上，因為耶和華用膏膏我，叫我傳好信息給謙卑的人，差遣我醫好傷心的人，報告被擄的得釋放、被囚的出監牢，報告耶和華的恩年和我們神報仇的日子，安慰一切悲哀的人，賜華冠於錫安悲哀的人代替灰塵，喜樂油代替悲哀，讚美衣代替憂傷之靈，使他們稱為公義樹，是耶和華所栽的，叫他得榮耀。\\  \midrule
徒10:38&神怎樣以聖靈和能力膏拿撒勒人耶穌，這都是你們知道的。他周流四方，行善事，醫好凡被魔鬼壓制的人，因為神與他同在。\\ 
徒4:27&希律和本丟‧彼拉多，外邦人和以色列民，果然在這城裡聚集，要攻打你所膏的聖僕耶穌，\\  \midrule
路4:18&主的灵在我身上，因为祂用膏膏我，叫我传福音给贫穷的人；差遣我报告被掳的得释放，瞎眼的得看见，叫那受压制的得自由。\\  \midrule

箴1:23&你們當因我的責備回轉，我要將我的靈澆灌你們，將我的話指示你們。\\  \midrule
太3:11-12&我是用水給你們施洗，叫你們悔改；但那在我以後來的，能力比我更大，我就是給他提鞋也不配，他要用聖靈與火給你們施洗。他手裡拿著簸箕，要揚淨他的場，把麥子收在倉裡，把糠用不滅的火燒盡了\\  \midrule
可1:8&有一位在我以後來的，能力比我更大，我就是彎腰給他解鞋帶也是不配的。 8我是用水給你們施洗，他卻要用聖靈給你們施洗。\\  \midrule
路3:16-17&約翰說：「我是用水給你們施洗。但有一位能力比我更大的要來，我就是給他解鞋帶也不配，他要用聖靈與火給你們施洗。他手裡拿著簸箕，要揚淨他的場，把麥子收在倉裡，把糠用不滅的火燒盡了。」\\  \midrule
約1:33&我先前不認識他，只是那差我來用水施洗的對我說：『你看見聖靈降下來，住在誰的身上，誰就是用聖靈施洗的。』\\  \midrule

約3:31-36&從天上來的是在萬有之上，從地上來的是屬乎地，他所說的也是屬乎地。從天上來的是在萬有之上，他將所見所聞的見證出來，只是沒有人領受他的見證。 33 那領受他見證的，就\textcolor{red}{印上印}，證明神是真的。 34 神所差來的就說神的話，因為神賜聖靈給他是沒有限量的。 35 父愛子，已將萬有交在他手裡。 36 信子的人有永生，不信子的人得不著永生，神的震怒常在他身上。\\  \midrule
约20:22&说了这话，就向他们吹一口气，说，你们受圣灵。\\  \midrule

路10:21-22&正當那時，耶穌被聖靈感動就歡樂，說：「父啊，天地的主，我感謝你！因為你將這些事向聰明通達人就藏起來，向嬰孩就顯出來。父啊，是的，因為你的美意本是如此。 22 一切所有的，都是我父交付我的。除了父，沒有人知道子是誰；除了子和子所願意指示的，沒有人知道父是誰。」\\  \midrule

太12:28&我若靠着上帝的灵赶鬼，这就是上帝的国临到你们了。\\  \midrule

约6:63&叫人活着的乃是灵。\\  \midrule
约4:24&上帝是个灵，所以拜祂的，必须用心灵和诚实拜祂。\\  \midrule
林后3:17&\textcolor{red}{主就是那灵}，主的灵在那里，那里就得以自由。\\  \midrule
林前15:42-47&死人復活也是這樣：所種的是必朽壞的，復活的是不朽壞的； 43 所種的是羞辱的，復活的是榮耀的；所種的是軟弱的，復活的是強壯的； 44 所種的是血氣的身體，復活的是靈性的身體。若有血氣的身體，也必有靈性的身體。 45 經上也是這樣記著說，首先的人亞當成了有靈的活人，\textcolor{red}{末後的亞當成了叫人活的靈}。 46 但屬靈的不在先，屬血氣的在先，以後才有屬靈的。 47 頭一個人是出於地，乃屬土；第二個人是出於天。\\  \midrule

罗1:4&按圣善的灵说，因从死里复活，以大能显明是上帝的儿子。\\  \midrule
罗8:11&然而叫耶稣从死里复活者的灵，若住在你们心里，那叫基督耶稣从死里复活的，也必借着住在你们心里的圣灵，使你们必死的身体又活过来。\\  \midrule

\end{tabular}%\caption{聖靈是聖父做以下工作的途徑}  
\end{table}









\subsubsection{圣灵对圣徒个别的指引}
\begin{table}[H]
\centering
\begin{tabular}{p{5em}p{2.4em}p{40em}}  \toprule
经节&人物&经文 \\  \midrule
徒8:29&腓利&聖靈對腓利說：「你去，貼近那車走！」\\  \midrule
徒8:39&腓利&從水裡上來，主的靈把腓利提了去。太監也不再見他了，就歡歡喜喜地走路。\\  \midrule
徒10:18-20&彼得&彼得還思想那異象的時候，聖靈向他說：「有三個人來找你。起來，下去，和他們同往，不要疑惑，因為是我差他們來的。」\\  \midrule
徒11:12&彼得&聖靈吩咐我和他們同去，不要疑惑。同著我去的，還有這六位弟兄，我們都進了那人的家。\\  \midrule

徒13:2-4&掃羅&他們侍奉主、禁食的時候，聖靈說：「要為我分派巴拿巴和掃羅，去做我召他們所做的工。」 於是禁食、禱告，按手在他們頭上，就打發他們去了。他們既被聖靈差遣，就下到西流基，從那裡坐船往塞浦路斯去。\\  \midrule


徒16:6&保羅&聖靈既然禁止他們在亞細亞講道，他們就經過弗呂家、加拉太一帶地方。\\  \midrule
徒16:7-10&保羅&到了每西亞的邊界，他們想要往庇推尼去，耶穌的靈卻不許。他們就越過每西亞，下到特羅亞去。在夜間有異象現於保羅：有一個馬其頓人站著求他說：「請你過到馬其頓來幫助我們！」保羅既看見這異象，我們隨即想要往馬其頓去，以為神召我們傳福音給那裡的人聽。\\  \midrule


徒20:22-24&保羅&現在我往耶路撒冷去，心甚迫切，不知道在那裡要遇見什麼事。但知道聖靈在各城裡向我指證，說有捆鎖與患難等待我。我卻不以性命為念，也不看為寶貴，只要行完我的路程，成就我從主耶穌所領受的職事，證明神恩惠的福音。\\  \midrule
徒21:3-4&門徒&我們就在推羅上岸，因為船要在那裡卸貨。找著了門徒，就在那裡住了七天。他們被聖靈感動，對保羅說：「不要上耶路撒冷去。」\\  \midrule

徒21:8-11&亞迦布&第二天，我們離開那裡，來到愷撒利亞，就進了傳福音的腓利家裡，和他同住。他是那七個執事裡的一個。他有四個女兒，都是處女，是說預言的。我們在那裡多住了幾天，有一個先知名叫亞迦布，從猶太下來，到了我們這裡，就拿保羅的腰帶捆上自己的手腳，說：「聖靈說：猶太人在耶路撒冷要如此捆綁這腰帶的主人，把他交在外邦人手裡。」\\  \midrule

\end{tabular}%\caption{創造}  
\end{table}




\subsubsection{主，天使（使者）对圣徒个别的指引}
\begin{table}[H]
\centering
\begin{tabular}{p{5em}p{2.4em}p{40em}}  \toprule
经节&人物&经文 \\  \midrule
徒5:19-20&使徒出监&但主的使者夜間開了監門，領他們出來， 20 說：「你們去站在殿裡，把這生命的道都講給百姓聽。」 21 使徒聽了這話，天將亮的時候就進殿裡去教訓人。\\  \midrule
徒12:7-11&彼得出监&忽然，有主的一個使者站在旁邊，屋裡有光照耀。天使拍彼得的肋旁，拍醒了他，說：「快快起來！」那鐵鏈就從他手上脫落下來。 8 天使對他說：「束上帶子，穿上鞋！」他就那樣做。天使又說：「披上外衣，跟著我來！」 9 彼得就出來跟著他，不知道天使所做是真的，只當見了異象。 10 過了第一層、第二層監牢，就來到臨街的鐵門，那門自己開了。他們出來，走過一條街，天使便離開他去了。 11 彼得醒悟過來，說：「我現在真知道主差遣他的使者，救我脫離希律的手和猶太百姓一切所盼望的。」\\  \midrule

徒12:17&彼得&彼得擺手不要他們作聲，就告訴他們主怎樣領他出監，又說：「你們把這事告訴雅各和眾弟兄。」於是出去往別處去了。\\  \midrule

徒10:3-8&哥尼流&有一天，約在申初，他在異象中明明看見神的一個使者進去，到他那裡，說：「哥尼流！」 4 哥尼流定睛看他，驚怕說：「主啊，什麼事呢？」天使說：「你的禱告和你的賙濟達到神面前，已蒙記念了。 5 現在你當打發人往約帕去，請那稱呼彼得的西門來。 6 他住在海邊一個硝皮匠西門的家裡，房子在海邊上。」 7 向他說話的天使去後，哥尼流叫了兩個家人和常伺候他的一個虔誠兵來， 8 把這事都述說給他們聽，就打發他們往約帕去。\\  \midrule

徒10:30-33&哥尼流&哥尼流說：「前四天這個時候，我在家中守著申初的禱告，忽然有一個人穿著光明的衣裳，站在我面前， 31 說：『哥尼流，你的禱告已蒙垂聽，你的賙濟達到神面前已蒙記念了。 32 你當打發人往約帕去，請那稱呼彼得的西門來，他住在海邊一個硝皮匠西門的家裡。』 33 所以我立時打發人去請你。你來了很好！現今我們都在神面前，要聽主所吩咐你的一切話。」\\  \midrule

徒11:13-14&哥尼流&那人就告訴我們他如何看見一位天使站在他屋裡，說：『你打發人往約帕去，請那稱呼彼得的西門來， 14 他有話告訴你，可以叫你和你的全家得救。』\\  \midrule




徒9:3-6&掃羅&掃羅行路，將到大馬士革，忽然從天上發光，四面照著他。他就仆倒在地，聽見有聲音對他說：「掃羅，掃羅，你為什麼逼迫我？」他說：「主啊，你是誰？」主說：「我就是你所逼迫的耶穌。起來！進城去，你所當做的事，必有人告訴你。」\\  \midrule
徒22:6-10&掃羅向猶太人&我將到大馬士革，正走的時候，約在晌午，忽然從天上發大光，四面照著我。我就仆倒在地，聽見有聲音對我說：『掃羅，掃羅，你為什麼逼迫我？』我回答說：『主啊，你是誰？』他說：『我就是你所逼迫的拿撒勒人耶穌。』與我同行的人看見了那光，卻沒有聽明那位對我說話的聲音。我說：『主啊，我當做什麼？』主說：『起來！進大馬士革去，在那裡要將所派你做的一切事告訴你。』\\  \midrule
徒26:16-18&掃羅向亞基帕王做见证&王啊，我在路上，晌午的時候，看見從天發光，比日頭還亮，四面照著我並與我同行的人。我們都仆倒在地，我就聽見有聲音用希伯來話向我說：『掃羅，掃羅，為什麼逼迫我？你用腳踢刺是難的。』我說：『主啊，你是誰？』主說：『我就是你所逼迫的耶穌。「你起來站著！我特意向你顯現，要派你做執事，作見證，將你所看見的事和我將要指示你的事證明出來。我也要救你脫離百姓和外邦人的手。我差你到他們那裡去，要叫他們的眼睛得開，從黑暗中歸向光明，從撒旦權下歸向神；又因信我，得蒙赦罪，和一切成聖的人同得基業。』\\  \midrule
徒22:17-21&保羅信主后回耶京&後來我回到耶路撒冷，在殿裡禱告的時候魂遊象外，看見主向我說：『你趕緊地離開耶路撒冷，不可遲延。因你為我作的見證，這裡的人必不領受。』我就說：『主啊，他們知道我從前把信你的人收在監裡，又在各會堂裡鞭打他們。並且你的見證人司提反被害流血的時候，我也站在旁邊歡喜，又看守害死他之人的衣裳。』 主向我說：『你去吧！我要差你遠遠地往外邦人那裡去。』」\\  \midrule

徒23:11&保羅被抓&當夜，主站在保羅旁邊，說：「放心吧！你怎樣在耶路撒冷為我作見證，也必怎樣在羅馬為我作見證。」\\  \midrule

徒27:23-24&保羅海上&因我所屬、所侍奉的神，他的使者昨夜站在我旁邊說：『保羅，不要害怕！你必定站在愷撒面前；並且與你同船的人，神都賜給你了。』 」\\\bottomrule
\end{tabular}%\caption{創造}  
\end{table}
　

\subsubsection{異象对圣徒个别的指引}
\begin{table}[H]
\centering
\begin{tabular}{p{5em}p{2.4em}p{40em}}  \toprule
经节&人物&经文 \\  \midrule
徒9:10-16&亞拿尼亞&在大馬士革有一個門徒，名叫亞拿尼亞。主在異象中對他說：「亞拿尼亞！」他說：「主，我在這裡。」主對他說：「起來！往直街去，在猶大的家裡訪問一個大數人，名叫掃羅。他正禱告，又看見了一個人，名叫亞拿尼亞，進來按手在他身上，叫他能看見。」亞拿尼亞回答說：「主啊，我聽見許多人說這人怎樣在耶路撒冷多多苦害你的聖徒， 並且他在這裡有從祭司長得來的權柄，捆綁一切求告你名的人。」主對亞拿尼亞說：「你只管去！他是我所揀選的器皿，要在外邦人和君王並以色列人面前宣揚我的名。我也要指示他，為我的名必須受許多的苦難。」\\  \midrule
徒10:9-16&彼得&彼得約在午正上房頂去禱告。覺得餓了，想要吃。那家的人正預備飯的時候，彼得魂遊象外，看見天開了，有一物降下，好像一塊大布，繫著四角，縋在地上，裡面有地上各樣四足的走獸和昆蟲，並天上的飛鳥。又有聲音向他說：「彼得，起來，宰了吃！」彼得卻說：「主啊，這是不可的！凡俗物和不潔淨的物，我從來沒有吃過。」 第二次有聲音向他說：「神所潔淨的，你不可當做俗物。」這樣一連三次，那物隨即收回天上去了。\\  \midrule
徒11:5-10&彼得&我在約帕城裡禱告的時候，魂遊象外，看見異象：有一物降下，好像一塊大布，繫著四角，從天縋下，直來到我跟前。我定睛觀看，見內中有地上四足的牲畜和野獸、昆蟲並天上的飛鳥。我且聽見有聲音向我說：『彼得，起來，宰了吃！』我說：『主啊，這是不可的！凡俗而不潔淨的物從來沒有入過我的口。』第二次有聲音從天上說：『神所潔淨的，你不可當做俗物。』這樣一連三次，就都收回天上去了。\\  \midrule
徒18:9-10&保羅&夜間，主在異象中對保羅說：「不要怕，只管講，不要閉口。有我與你同在，必沒有人下手害你。因為在這城裡我有許多的百姓。」（哥林多）\\  \midrule
\end{tabular}%\caption{創造}  
\end{table}

























\subsection{圣灵恩赐}



\subsection{從林前12-14, 羅12, 弗4中看保羅對聖靈恩賜(事奉的恩賜)的教導}

耶穌在最後晚餐時，曾安慰門徒說：「我要求父，父就另外賜給你們一位保惠師 ，叫他永遠與你們同在，就是真理的聖靈」（約14:16,17a）。
耶穌復活后的第一天晚上同眾門徒在一起的時候，再次提醒門徒等候父所要差遣的聖靈（約20:21,22, 路24:49）。
耶穌升天後，五旬節門徒聚集在耶路撒冷禱告的時候被聖靈充滿，「應許的聖靈」得到應驗。「應許的聖靈」是神賜給每一位相信耶穌名的人的印記（徒2:38）。
從五旬節信徒在降下聖靈后就開始說方言的記載，可以知道「聖靈的恩賜」與聖靈同時賜給教會。
聖靈隨己意將「聖靈的恩賜」分給各人（林前12:4-11）。 

\subsubsection{恩賜的目的及意義}
\subsubsubsection{目的}
\begin{itemize}
\item 聖靈顯在各人身上，是叫人得益處(林前12：7) 
\begin{table}[H]
\centering
\begin{tabular}{p{17.5em}p{27.5em}}  \toprule
使信徒之間有區別&不同的恩賜使得信徒成爲基督身上不同的肢體彼此配搭服事（林前12:12-27）\\  \midrule
言語上：智慧的言語，知識的言語，說方言，翻譯方言和先知的預言。&保羅用身體的肢體做比喻，解釋我們每個人所得的恩賜不一樣，所以我們可
以成為基督身上不同的肢體（林前12:12），因為我們不同，所以我們可以彼
此相顧（林前12:25）。\\  \midrule
甚至人可以具有超越人類能力的作為，比如可以醫病，行異能和辨別諸靈等&不同的恩賜使得信徒彼此不可缺少，尤其是軟弱的，更是不可缺少的（林前
12:22）。\\  \midrule
工作的超人能力,比如超自然的信心，服事，治理&不同的恩賜使得信徒彼此不足，神要用榮耀的肢體去補足不俊美的肢體（林
前12:23，24）\\  \midrule
\end{tabular}%\caption{創造}  
\end{table}

\item 叫人蒙恩典\\ 
我們各人蒙恩，都是照基督所量給各人的恩賜使我們蒙恩典(弗4:7) 。 恩賜是賜給了每一個信徒的（each one of us 1 ）。
\item 造就教會(林前14：12)\\ 
做先知講道的就是造就教會的恩賜（林前14：4b）。保羅特別鼓勵信徒要切慕造就教會
的恩賜。
\item 造就信徒（林前14:4）\\ 
有的恩賜是造就個人的，比如說方言的恩賜（林前14：4a) 。
\item 為救恩做見證\\ 
這救恩起先是主親自講的，後來是聽見的人給我們證實了， 現在神又按自己的旨意，用神蹟、奇事和百般的異能並聖靈的恩賜為這救恩作見證(來2:4) 。
\end{itemize}

\subsubsubsection{意義}
\begin{itemize}
\item 恩賜是神的恩典（弗4：7-8）\\ 
恩賜不是我們天然的才幹，既然恩賜是恩典，就不是我們靠著自己的努力可以得到
的。
\item 使信徒可以做不同的神的工作（羅12：5）\\ 
比如說預言，治理神的教會。
\item 使不同的信徒成為一個身體（羅12：5）\\ 
雖然我們恩賜不同，但卻連在一個身體上，一同建造這同一個身體，就是基督的教
會。
\item 使信徒彼此不同\\ 
恩賜原有分別（林前12:4），不同的恩賜使得信徒成爲基督身上不同的肢體彼此配
搭服事（林前12:12-27）。
\begin{itemize}
\item 保羅用身體的肢體做比喻，解釋我們每個人所得的恩賜不一樣，所以我們可以成
為基督身上不同的肢體（林前12:12），因為我們不同，所以我們可以彼此相顧（
林前12:25）。
\item 不同的恩賜使得信徒彼此不可缺少，尤其是軟弱的，更是不可缺少的（林前
12:22）。
\item 不同的恩賜使得信徒彼此不足，神要用榮耀的肢體去補足不俊美的肢體（林前
12:23，24）。
\end{itemize}
\end{itemize}

\subsubsection{領受恩賜的方式}
\begin{itemize}
\item 聖靈隨己意運行分配給各人（林前12:11）
\begin{itemize}
\item 保羅做福音的執事是照神的恩賜，這恩賜是照神運行的大能賜給保羅的（弗3:7）
\item 五旬節到了，門徒都聚集在一處。忽然，從天上有響聲下來，好像一陣大風吹過，
充滿了他們所坐的屋子；又有舌頭如火焰顯現出來，分開落在他們各人頭上。他們
就都被聖靈充滿，按著聖靈所賜的口才說起別國的話來(徒2:1-3)
\item 彼得在哥尼流家裏正傳講福音的時候，聖靈降在一切聽道的人身上。那些奉割禮、
和彼得同來的信徒，見聖靈的恩賜也澆在外邦人身上，就都稀奇，因聽見他們說方
言，稱讚神為大。(徒10:44-46)
\item 恩賜是不同的，但源自同一位聖靈（林前12:4）.
\item 當人信主後，聖靈要顯在信徒的各人身上（林前12:7），就把對信徒有益處的不同
恩賜賜給他們。這種賜給是聖靈主動作的，不是信徒主動要求的，信徒是接受者。
\item 恩賜有多種，每個信徒身上的恩賜不一樣，這也不是信徒可以要求或決定的，給每
個信徒身上什麼樣的恩賜都是聖靈的工作，聖靈是隨己意分給各人的。也就是說每
個人身上的恩賜，自己不能決定要和要什麼，或者不要，全是聖靈分配的，人不能
拒絕，也不能排斥，人能作的就是把它發揮使用出來，建造教會。
\end{itemize}

\item 藉著人按手而領受聖靈的恩賜。聖經中的這類例子包括：

\begin{itemize}
\item 提摩太领受圣灵的恩賜就是從前藉著預言，在保罗（提後1：6）或者眾長老（提前
4：14）按手的時候賜給提摩太的。
\item 以弗所地区，有十二个接受施洗约翰洗礼的门徒聽見保罗说要信耶稣的話，就奉主
耶穌的名受洗。 保羅按手在他們頭上，聖靈便降在他們身上，他們就說方言，又說
預言 （徒19：6，7）
\item 撒馬利亞人領受了神的道并奉主耶穌的名受了洗，但聖靈還沒有降在他們一個人身
上，彼得、約翰从耶路撒冷到了他們那裡，就為他們禱告，使徒按手在他們頭上，
他們就受了聖靈。（徒8：18-23）
\end{itemize}
\end{itemize}

\subsubsection{恩賜的分類}
\subsubsubsection{恩賜的種類及分類簡介}

\begin{table}[htbp]
\centering
\begin{tabular}{p{6.145em}|p{17.785em}|p{21.5em}}\toprule
 \multicolumn{1}{p{6.215em}|}{\textbf{類別}} & \multicolumn{1}{p{5.215em}|}{\textbf{恩賜的名稱}} & \textbf{经节} \\\midrule
話語的恩賜& 智慧的言語&（林前12：8a）\\\midrule
&知識的言語&（林前12：8b）\\\midrule
&先知預言&（羅12：6）（林前12：10b）（林前12：28b）\\\midrule
&說方言（different kinds of tongues 2 ）&（林前12：10d）（林前12：28h）\\\midrule
&翻方言&（林前12：10e）\\\midrule
&教導&（羅12：7b）\\\midrule
&勸化&（羅12：8a）\\\midrule

異能的恩賜& 醫病（gifts of healing 3 ）&（林前12：9b）（林前12：28e）\\\midrule

&行異能&（林前12：10a）（林前12：28d）\\\midrule
&辨別諸靈&（林前12：10c）\\\midrule
工作的恩賜 &信心&（林前12：9a）\\\midrule

&執事Serve 4 (服事)&（羅12：7a）\\\midrule
&施捨&（羅12：8b）\\\midrule
&治理 Lead 5 &（羅12：8c）\\\midrule
&憐憫&（羅12：8d）\\\midrule
\end{tabular}%
\end{table}%

\subsubsubsection{話語的恩賜}
話語的恩賜包括:智慧的言語，知識的言語，先知的預言和講道，使用人間的言語或者是天使的言語說方言（different kinds of tongues 6 ），翻譯方言，教導和勸化。因爲這些恩賜的彰顯都是通過嘴巴，所以都歸類為話語的恩賜。下面用表格來簡述這几種恩賜的説明和使用原則。

\begin{table}[htbp]
\centering
\begin{tabular}{p{6.145em}|p{30.785em}|p{10.5em}}\toprule
 \multicolumn{1}{p{6.215em}|}{\textbf{恩賜的名稱}} & \multicolumn{1}{p{15.215em}|}{\textbf{原則簡單説明}} & \textbf{经节} \\\midrule
  
智慧的言語&不一定每個人都有這個恩賜&（林前12：8a）\\\midrule

知識的言語&不一定每個人都有這個恩賜&（林前12：8b）\\\midrule

&知識的言語需要有愛來搭配&（林前13:2）\\\midrule
&知識必歸於無有&(林前13:8)\\\midrule

先知&在教會中僅次於使徒&(林前12:28)\\\midrule
&可以造就教會，勸勉安慰人&（林前14:3）\\\midrule
&先知分成以下兩種。做先知的恩賜是對人說；這個恩賜最終會歸於無有&（林前13:8）\\\midrule
&先知所知道的都是有限的，直到見到主時我們所知道
的才能完全&（林前13:12）\\\midrule
&1 在教會中說預言的先知 說預言時當照著信心的程度說預言
2 在聚會中講道的先知做先知講道的人可以明白各樣的奧祕
比如保羅&（林前13:2）\\\midrule
&在聚會中，先知的講道可以造就教會。
聚會時先知若只說方言，不用啟示，或
知識，或預言，或教訓來講解就對衆人
沒有益處&（林前14:6）\\\midrule
&先知講道是為信的人做證據。要切慕做先知講道。&\\\midrule

說方言（different kinds of tongues 7 ）&

按照聖經教導，說方言之能終必停止(
林前13：8)。英文的方言是複數，所以
方言不衹是一種 8 。按照方語的種類分
有：

1，萬人的方言（林前13：1)

2，天使的言語（林前13：1)

如果按照說方言的内容分類有：



1，信息方言（會衆聚集的時候）

1)，會衆聚集的時候，說的方言也可能
是地上萬人的方言，也有可能是天使
的言語。

2)，說方言不是為信的人做證據，乃是
為不信的人(林前14：22a)

說方言的若不翻出來，使教會被造就，
那做先知講道的就比他強(林前14：5)
那說方言的，就當求著能翻出來(林前
14：13)。若沒有人翻，就當在會中閉口
(林前14:28) 。全教會聚在一處的時候，
不要都說方言 (林前14：23)，若有說方
言的，只好兩個人，至多三個人，且要
輪流著說，也要一個人翻出來 (林前
14：27)。要切慕做先知講道，也不要禁
止說方言(林前14:39)還有更多


2，禱告方言（自己獨處的時候）

禱告的方言不是對人說，乃是對神說
，因為沒有人聽出來，然而人在心靈
裡卻是講說各樣的奧祕(林前14：2)。
這種方言指用天使的言語向神禱告。
保羅說這種方言比衆人都多（林前
14：18），但他在教會中寧願說造就
教會的話（林前14：19）。

說這種方言人的是造就自己的靈(林前
14:4a)，這種禱告的言語是靈在禱告，悟
性沒有果效(林前14:14)。信徒既要用靈
禱告也要用悟性禱告(林前14:15a)。&保羅盼望信徒人人都可以說方言(林前
14：5)，由此可以推知這個恩賜是聖靈
賜給每一位信徒的恩賜。\\\midrule

%下面按照第二種分類方法來討論說方言的恩賜。
%信息方言

%禱告方言


翻方言&「翻方言」是「聖靈的恩賜」(林前
12:10)，適用於會眾聚會時使用而且是
造就教會的恩賜(林前14:5)
&
「說方言講道」時倘若沒有「翻方言」
的人，就當閉口不言。「說方言講道」
和「翻方言」是同時出現的恩賜(林前
14:13)\\\midrule

教導&當專一&\\\midrule
勸化&當專一&\\\midrule
\end{tabular}%
\end{table}%


智慧的言語
智慧的言语来自神。智慧的特点是清潔，和平，溫良柔順，滿有憐憫，多結善果，沒有偏見，沒有假冒(雅3:17)。我們中間若有缺少智慧的，應當求那厚賜與眾人、也不斥責人的神，主就必賜給信徒(雅1:5)。
\begin{table}[H]
  \centering
    \begin{tabular}{p{15.145em}|p{15.785em}|p{16.5em}}    \toprule
    \multicolumn{1}{p{15.145em}|}{\textbf{恩賜用途或特色}} & \multicolumn{1}{p{15.785em}|}{\textbf{彰顯的情形}} &  \multicolumn{1}{p{16.5em}}{\textbf{運用及原則}} \\    \midrule

根据林前6:5，保罗指
出智慧言语的恩赐可以
审断教会内部成員之間
的紛爭。& 所罗门王用智慧判断兩個婦人
誰才是活孩子的母亲，以至于
以色列眾人聽見王這樣判斷，
就都敬畏他；因見他有神的智
慧，能以斷案(王上3:28) &

1教会中间有彼此相争的事，不
要在不义的人面前求审，而是
在圣徒面前求审(林前6:1) 
2圣徒要审判世界(林前6:2) 
3圣徒要审判天使(林前6:3) \\\midrule

智慧言语的恩赐可以使
人心归向主&
施洗约翰有叫為父的心轉向兒
女，叫悖逆的人轉從義人的智
慧。(路1:17)&
使用的人必须有以利亚的心志
，行在主的前面(路1:17) \\\midrule

主賜的口才、智慧使圣
徒在君王諸侯面前做见证(路21:15)&
司提反是以智慧和聖靈說話，
眾人敵擋不住 (徒6:10)&
當立定心意，不要預先思想怎
樣分訴(路21:14)\\\midrule

这个恩赐可以使圣徒更
明白神的旨意。因此，
保罗不住的禱告祈求父
，將那賜人智慧和啟示
的靈賞給圣徒，願圣徒
在一切屬靈的智慧悟性
上，滿心知道神的旨意
（西1:9，弗1:17) &

這個恩賜可以幫助明白預言。
在這裡有智慧：凡有聰明的，
可以算計獸的數目；因為這是
人的數目，他的數目是六百六
十六(啟13:18)。
智慧的心在此可以思想。那七
頭就是女人所坐的七座山(啟
17:9)&
不可以随己意私解，或强解圣
经。\\\midrule

但以理和约瑟都提到他
们可以解梦的原因是神
赐的智慧。所以这个恩
赐可以解明奥秘的事&
但以理为王解夢(但2:30)
约瑟为法老，酒正及膳食长解
梦(创40-41:36)&
不是自己主动使用这个恩赐而
是神赐时机，约瑟和但以理解
梦的机会都是神赐的。\\\midrule

彼得在彼後3:15提到保
罗的书信是用神给的智
慧恩赐寫成的。&

新约圣经中保罗的13封书信&这个恩赐仅限于初代教会。现
在应该没有了。\\\midrule

用聖靈所指教的言語講
說神的奧祕，將屬靈的
話解釋屬靈的事(林前
2:13)&
我們傳揚他，是用諸般的智慧
，勸戒各人，教導各人，要把
各人在基督裡完完全全的引到
神面前(西1:28) &
1只限于用在完全(屬靈) 的人
中(林前2:6) 
2講從前所隱藏、神奧祕的智慧
，就是神在萬世以前預定使我
們得榮耀的(林前2:7) \\\midrule

彼此教導，互相勸戒，
歌頌神。&
&用智慧把基督的道理豐豐富富
的存在心裡(西3:16)\\\midrule
\end{tabular}%
\caption{智慧的言語}
\label{tab:addlabel}%
\end{table}%


\begin{table}[H]
  \centering
    \begin{tabular}{p{15.145em}|p{15.785em}|p{16.5em}}    \toprule
    \multicolumn{1}{p{15.145em}|}{\textbf{恩賜用途或特色}} & \multicolumn{1}{p{15.785em}|}{\textbf{彰顯的情形}} &  \multicolumn{1}{p{16.5em}}{\textbf{運用及原則}} \\    \midrule

“分別善惡對錯的知识或见识”，比如创
世纪中的分別善惡的樹,KJV英文翻译为【
善惡知识的樹(tree of knowledge of
good and evil)】(创2：9)。申命记中
，【今日不知善惡的兒女，必進入那地
（had no knowledge between good and
evil）】(申1：39) KJV英文翻译为【没
有善惡知识的兒女】&
保羅表示對於吃祭偶像之物
的態度或者看法屬於知識的
範圍。“論到祭偶像之物，
我們曉得我們都有知識（林
前8:1）”&
但知識是叫人自
高自大，惟有愛
心能造就人（林
前8:1）\\    \midrule

舊約中知識当作“舊約聖經或者律法”来
解释。比如【但以理啊，你要隱藏這話，
封閉這書，直到末時。必有多人來往奔跑
，知識就必增長】（但12:4）【我的民因
無知識而滅亡。你棄掉知識，我也必棄掉
你，使你不再給我作祭司。你既忘了你神
的律法，我也必忘記你的兒女】(何
4:6）【祭司的嘴裡當存知識，人也當由
他口中尋求律法，因為他是萬軍之耶和華
的使者 】（瑪2:7）【是蠢笨人的師傅，
是小孩子的先生，在律法上有知識和真理
的模範】（羅2:20）

主耶穌在責備律法師時，也稱律法為知識
。【因為你們把知識的鑰匙奪了去，自己
不進去，正要進去的人你們也阻擋他們
】。（路11:52 ）
&
1，猶太人做假見證控告司提
反糟蹋聖所和律法，他在大
祭司面前為自己的辯訴，應
當用的是知識的言語。他的
辯詞大量引用了摩西律法中
記載的事實，表明自己不但
熟知律法和聖殿的建造目的
，而且遵守律法（徒7）

2，保羅在彼西底的安提阿會
堂裏，對猶太人的講道用的
是知識的言語。他的整篇講
道引用了大量的聖經舊約，向猶太人論證神在舊約裏向
他們祖宗所應許的，完全應
驗在耶穌身上，使得許多人
歸主（徒13：16-41 ）。&

1，衹有在基督
裡才有全備的知
識（林前1:5）
2，知識的言語
需要有愛來搭配
（林前13:2）
3，知識必歸於
無有(林前13:8)
4, 聚會時先知
若只說方言，不
用啟示，或知識
，或預言，或教訓來講解就對衆
人沒有益處（林
前14:6）\\    \midrule

英文KJV翻譯中知識包括對於真理的認識
。他願意萬人得救，明白真道。Who will
have all men to be saved, and to
come unto the knowledge of the
truth（提前2:4）。常常學習，終久不能
明白真道 Ever learning, and never
able to come to the knowledge of the
truth（提後3:7）.&&\\    \midrule
對神的認識屬於知識的範圍。認識耶和華
榮耀的知識要充滿遍地，好像水充滿洋海
一般（哈2:14）以及經文：林前15:34 ，
林後10:5， 弗4:13，西1:10，何6:6。&&\\    \midrule
\end{tabular}%
\caption{知識的言語}
\label{tab:addlabel}%
\end{table}%


3. 先知預言
做先知的恩賜是對人說，可以造就教會，勸勉安慰人（林前14:3），這個恩賜最終
會歸於無有（林前13:8）。先知所知道的都是有限的，直到見到主時我們所知道的
才能完全（林前13:12）。先知的恩賜有四種職能：1，透過說出對將來未發生的事
物的預言，幫助教會或者教會的肢體做好準備；2，透過講解，啓示，勸勉將真理解
開辨明造就教會；3，透過唱贊美詩歌，頌詞，和靈歌勸慰鼓勵造就衆人；4，直接
聽到神的旨意。下面分別敘述。
\begin{table}[H]
  \centering
    \begin{tabular}{p{15.145em}|p{15.785em}|p{16.5em}}    \toprule
    \multicolumn{1}{p{15.145em}|}{\textbf{特色}} & \multicolumn{1}{p{15.785em}|}{\textbf{彰顯的情形}} &  \multicolumn{1}{p{16.5em}}{\textbf{運用及原則}} \\    \midrule

 
A 做先知说预言& & \\    \midrule
1，預言是人被聖
靈感動，說出神
的話來
2，先知有男(如
亞迦布，徒
21:10)也有女先
知(如執事腓利的
四個女兒.徒
21:9)&

1，亞迦布藉著聖靈指出天下
在革老丟時期將要有大饑荒(
徒11:28)
2，亞迦布告訴保羅聖靈說：
猶太人在耶路撒冷要捆綁保羅
並把他交在外邦人手裡(徒
21:10）
3，保羅在被押解去羅馬的船
上，告訴衆人喫飯，全船的人
都必要得救(徒27:21-41）&
听預言的人

1，在預言上留意，所有的預言沒
有可隨私意解說的(彼后1:19,20)
2，不要藐視先知的講論(帖前
5:20)
3，要察驗先知的講論並遵守(帖前
5:20)

說預言的人

只能傳達聖靈的意思而非自己的意
思。如亞迦布衹是忠實的講出保羅
被抓，而沒有阻止保羅違背聖靈的
意思。\\  \midrule

B 做先知講道&&\\    \midrule
1，新約教會中先
知的職能還包括
講道（林前
14：1）。
2，目地是透過傳
講把人引到基督面前（西1：28）
3, 但做先知講道
的可以造就、安
慰、勸勉人（林
前14：3）。&

1，彼得在五旬節聖靈降下后
對聚集的衆人的講論聖靈的降
臨是為耶穌做見証並勸衆人受
洗悔改(徒2:14-42）
2，彼得醫治好生來瘸腿的人
后對聚集的衆人講述醫治病人的原因並勸衆人悔改(徒2:12-
26)屬於先知的講道。
3，保羅在彼西底的安替阿會
堂的講論屬於先知講道(徒
13:16-41）&

1，說預言當照著信心的程度說預言(羅
12:6)
2，使用的途徑包括用啓示，知識，預
言，教訓，或講解的方法對教會講道（
林前14：6）
3，做先知講道的要專心祈禱和傳導(徒7:4）
4, 先知講道要把基督的奧秘揭示出來
（西4：2）\\    \midrule

C 先知唱歌&&\\    \midrule
1，根據代下
25：1的定義，唱
歌的事奉也是屬
於先知的恩賜範
圍。
2，透過唱歌可以
把基督的道理存
在心裏（西
3：16）
3，透過唱歌可以
彼此勸導歌頌神
（西3：16）
4，詩歌的信息中
也可以包括對未
來事情的預言，
如詩篇中的經節&

1，保羅和西拉在腓立比的監
獄唱詩贊美神，地大震動，監
獄門打開，囚犯的鎖鏈也鬆開
（徒13:16-41）
2，約沙法王曾經用唱歌的人
走在軍隊前面擊敗仇敵的戰例
（代下20：20-23）
3，約沙法王和其他兩個王與
摩亞打仗時，曾求問以利沙，
彈琴的時候耶和華的靈就降
在以利沙身上（王下3:13-
15）&

1, 舊約參與唱歌服事的人必須是利未
人（代上25）
2, 參與唱歌事奉的人需要接受訓練（
代上25：6-8）
3, 詩歌的歌詞中含有許多不同的信息
，所以在教會中用唱歌事奉的人需要參
照“說預言當照著信心的程度說預言(
羅12:6)”的原則，唱詩的時候需要有
信心。
4, 唱詩的事奉是站在征戰的最前綫（
代下20：20-23。
5, 要用靈歌唱也要用悟性歌唱(林前
14:15b)
6, 詩歌种类包括詩章，頌詞和靈歌（
西3：16）
6，詩歌常常和各種樂器同時使用（詩
33:3, 啟5:8-9\\    \midrule

D 直接聽到神的旨意&&\\    \midrule
聼到神的心意&安提阿教會的幾位先知和教師
禁食的時候聽到聖靈要分派巴
拿巴和掃羅去宣教(徒13:16-
41)
保羅得到馬其頓意向(徒16:6-
9)&
得到下一步行動方向或目標\\    \midrule
\end{tabular}%
\caption{先知預言}
\label{tab:addlabel}%
\end{table}%


4. 說方言 （different kinds of tongues）
按照聖經教導，說方言之能終必停止(林前13：8)。保羅盼望信徒人人都可以說方言
(林前14：5)，由此可以推知這個恩賜是聖靈賜給每一位兒女的恩賜。英文的方言是
複數，所以方言不衹是一種 2 。按照方語的種類分有：
- 萬人的方言（林前13：1)
- 天使的言語（林前13：1)
如果按照說方言的内容分類有：
- 信息方言（會衆聚集的時候）
- 禱告方言（自己獨處的時候）


\begin{table}[htbp]
  \centering
    \begin{tabular}{p{15.145em}|p{15.785em}|p{16.5em}}    \toprule
\multicolumn{1}{p{15.215em}|}{\textbf{特色}} & \multicolumn{1}{p{15.215em}|}{\textbf{彰顯的情形}} &\multicolumn{1}{p{15em}|}{\textbf{運用及原則}} \\
    \midrule
  
A信息方言（會衆聚集的時候）&& \\    \midrule
1，會衆聚集的時候，說
的方言也可能是萬人的
方言，也有可能是天使
的言語。
2，說方言不是為信的人
做證據，乃是為不信的
人(林前14：22a)&

五旬節到了，門徒都聚集在一
處，都被聖靈充滿，按著聖靈
所賜的口才說起別國的話來(
徒2:1-3)。可見五旬節門徒講
的是人間的言語。
保羅按手在他們頭上，聖靈便
降在他們身上，他們就說方言
，又說預言(徒19:6)&

說方言的若不翻出來，使教會
被造就，那做先知講道的就比
他強(林前14：5)
那說方言的，就當求著能翻出
來(林前14：13)
若沒有人翻，就當在會中閉口
。(林前14:28)
全教會聚在一處的時候，不要
都說方言 (林前14：23)
若有說方言的，只好兩個人，
至多三個人，且要輪流著說，
也要一個人翻出來。 (林前
14：27)
要切慕做先知講道，也不要禁
止說方言(林前14:39)\\    \midrule

B禱告方言（自己獨處的時候）&&\\    \midrule
禱告的方言不是對人說
，乃是對神說，因為沒
有人聽出來，然而人在
心靈裡卻是講說各樣的
奧祕(林前14：2)。這種
方言指用天使的言語向
神禱告。&

保羅說這種方言比衆人都多（
林前14：18），但他在教會中
寧願說造就教會的話（林前
14：19）。
聖靈用說不出的嘆息替我們照
著神的旨意求（羅8:26）&

說這種方言人的是造就自己的
靈(林前14:4a)
這種禱告的言語是靈在禱告，
悟性沒有果效(林前14:14)。
既要用靈禱告也要用悟性禱告(
林前14:15a)。\\    \midrule
\\    \bottomrule
    \end{tabular}%
     \caption{說方言}
  \label{tab:addlabel}%
\end{table}%




5. 翻譯方言
a. 「繙方言」是「聖靈的恩賜」(林前12:10)
b. 適用於會眾聚會時使用而且是造就教會的恩賜(林前14:5)
c. 「說方言講道」時倘若沒有「繙方言」的人，就當閉口不言。「說方言講道」和
「繙方言」是同時出現的恩賜(林前14:13)
6. 醫病
醫病的恩賜即可以醫治身體又可以醫治心靈。英文NIV用的是gifts of healing，恩
赐用的是复数。


\begin{table}[htbp]
  \centering
    \begin{tabular}{p{15.145em}|p{15.785em}|p{16.5em}}    \toprule
\multicolumn{1}{p{15.215em}|}{\textbf{特色}} & \multicolumn{1}{p{15.215em}|}{\textbf{彰顯的情形}} &\multicolumn{1}{p{15em}|}{\textbf{運用及原則}} \\
    \midrule
A，医治身体的疾病&&\\    \midrule
1 自己到主
面前求&

十個麻風病人得醫治（路17：11-19）
耶利哥的瞎子巴底買得醫治(太20:29-34)&

耶穌用叙利亚国的將軍乃缦大
麻風得潔净的例子指出不是所
有的人都得到醫治（路4:27）\\    \midrule

2 別人代求
而得到醫治&

百夫長為僕人代求（太8：5-10）
彼得的岳母的醫治（太8：14-17）&\\    \midrule

3 醫治者主動去醫治病
人&保羅在路司得城醫治生來是瘸腿的人 (徒 
14:8-10)
耶穌醫治好畢士大池子邊的癱子(約5:1-9)&
醫治者因查驗被醫治著是否有信心。\\    \midrule

B，医治灵里的疾病&&\\    \midrule
靈裏被鬼附
體而導致身
體上的疾病&

1，耶穌醫治因被鬼附而害害癲癇病的男孩
。（可9:14-2)
2，耶穌趕出的啞巴鬼使人不能説話（路
11:14）&

1，主賜人趕鬼權柄才可以把鬼
趕出去(可3：15)。猶太祭司長
士基瓦的七個兒子擅自稱主耶
穌的名趕鬼反被鬼制服了(徒
19:13-18)
2，主耶穌醫治了害癲癇病的男
孩后，告訴門徒有的鬼必須禱
告加禁食才可趕出來(太17:21)\\    \midrule

靈裏被鬼附
而導致精神
上不能自控&

保羅在腓立比從一個使女身上後來趕出她
裏面的鬼 (徒16:16-18)
格拉森被鬼附體的人（可5：1-20）&\\    \midrule

靈裏被罪捆
綁的得醫治&

1，有罪婦人膏主，因罪得赦免而得靈裏的
平安(路 7:36-50)
2，井邊打水的婦人因靈裏不滿足而換了好
幾位丈夫，因信耶穌而得靈裏的飽足（約
4：1-28）&\\    \midrule
C，同时医治身体和灵里的疾病&&\\    \midrule
罪得赦免，
身體得醫治&

1，被四個人抬的癱子罪被赦免，身體得到
了醫治（可2：1-12）
2，耶穌醫治了生來瞎眼的人，此人身體的
醫治，靈裏認識了耶穌（約9）&

彰顯神的權柄和榮耀時使用。\\    \midrule

靈裏得平安
，身體得醫
治&

1，十個大麻風中的一個被醫治后，因回來
歸榮耀給神而得拯救（路17：11-19）
2，血露婦人身體得到醫治并且得到平安（
可5：25-34）&\\    \midrule

長老和牧師
的醫治能力&

如果病人請求，長老可以為病人抹油，出
於信心的祈禱可以救這病人（雅5：14）&

\\    \bottomrule
    \end{tabular}%
     \caption{醫病}
  \label{tab:addlabel}%
\end{table}%


行異能
這種恩賜是指人行出的事情超出人的能力或者自然律。這是神爲了彰顯祂的全能而
賜給某些肢體的特殊能力 3 。
\begin{table}[H]
  \centering
    \begin{tabular}{p{20.145em}|p{23.785em}}    \toprule
\multicolumn{1}{p{20.215em}|}{\textbf{行的異能的人}} & \multicolumn{1}{p{23.215em}}{\textbf{彰顯的情形}}  \\    \midrule

使徒保羅所行
的異能或者神
跡&保羅使有法术、假充先知的犹太人巴耶稣瞎眼（徒13:5-11）
保羅和西拉在腓立比监狱里祷告地大震动(徒16:25，26)
神藉保罗的手行了些非常的奇事，甚至有人从保罗身上拿手巾或围
裙放在病人身上，病就退了，恶鬼也出去了（徒19:11-12）
保羅复活少年人犹推古（徒20:7-12）\\    \midrule

彼得所行的異
能或者神跡&彼得使多加死裏復活(徒9:36-42)
彼得使欺哄聖靈的亞拿尼亞同他的妻子撒非喇死亡（徒5:1-11)\\    \midrule

其他人行的異
能或者神跡&司提反滿得恩惠、能力，在民間行了大奇事和神蹟(徒6:8)
西門自己也信了，既受了洗，就常與腓利在一處，看見他所行的神
蹟和大異能，就甚驚奇(徒8:13)
主借使徒的手行了許多神跡奇事(徒5:12)
\\    \bottomrule
    \end{tabular}%
     \caption{行異能}
  \label{tab:addlabel}%
\end{table}%





辨別諸靈
辨别諸靈的恩赐中的靈是复数。世界上的靈有四种：动物的靈，人的靈，神的圣靈
，各种邪靈。根據約翰福音4:24記載神是個靈，我們人看不見，同樣邪靈也不能憑
肉眼可見。靈可以藉著動物和人的身體說話或者做事。所以這裡討論的辨別諸靈主
要是指辨別在人或者動物裡的諸靈。
人之所以能辨明诸靈是因為神藉著聖靈向我們顯明了(林前2:10)。屬靈的人能看透
萬事，卻沒有一人能看透了他(林前2:15)。下面讨论分辨動物裡的和人裡的諸靈。
\begin{table}[H]
  \centering
    \begin{tabular}{p{15.145em}|p{21.785em}|p{10.5em}}    \toprule
\multicolumn{1}{p{15.215em}|}{\textbf{恩賜用途或特色}} & \multicolumn{1}{p{21.785em}|}{\textbf{彰顯的情形}} &\multicolumn{1}{p{10.5em}|}{\textbf{運用及原則}} \\    \midrule
1动物裡的諸靈&&\\    \midrule

“誰知道人的靈是往上升
，獸的魂是下入地呢”(
傳道書 3:21)？中文这里
翻译的是“獸的魂”，而
NIV/KJV/ERV英文版本用
的都是獸的靈
(spirit）。神的靈和邪靈都可以附在动物身上。&

创世记第3章，记载撒旦曾经借着蛇的口
在伊甸园引诱人犯罪。

主耶稣在格拉森从一人身上赶出一群鬼，
鬼附猪身，导致两千头猪死亡（可5：1-20）。

神曾经借着驢对巴兰说话，拦阻他的愚妄
（民22：27-30）&
圣经中没有记载如何分辨动物的灵。\\    \midrule

分辨人靈里的諸靈&&\\    \midrule
分辨人心中的意念

使徒彼得和保羅在傳福音
時可以分辨出人心中抵擋的靈&

彼得在撒玛利亚时看出行邪术的西門正在
苦膽之中，被罪惡捆綁(徒8:23) 。

保羅被聖靈充滿看出行法術的以呂馬心中充滿
各樣詭詐奸惡，是魔鬼的兒子、眾善的仇
敵 (徒13:9-11)

保羅在路司得城遇到一個生來是瘸腿的人
聽保羅講道，保羅見他有信心，可得痊癒
，就醫治了他 (徒148-10)

儅主耶穌對四個人抬的癱子說：「小子，
放心吧，你的罪赦了！」他看出有幾個文
士心裡說：「這個人說僭妄的話了！」&

我們是屬神的，認
識神的就聽從我們
，不屬神的就不聽
從我們。從此我們
可以認出真理的靈
和謬妄的靈來（約
一4:6）。\\    \midrule

分辨人心中被惡者的靈充滿

靠著聖靈的恩賜人可以分
辨出人心中的惡念&

彼得分辨出撒旦充滿了亞拿尼亞的心，叫
他同妻子撒非喇欺哄聖靈，把田地的價銀
私自留下幾份。 (徒5:1-11)

耶穌分辨出猶大因自身生命的破口使得撒
旦就入了他的心（約13:27）。

儅主說要上十字架時，彼得試圖勸阻時，
主分辨出這是撒旦在利用彼得攔阻主的工
作（太16：23）。&\\    \midrule

分辨先知說預言者的靈

愚蒙人是話都信，通達人
步步謹慎。神賜給我們聖
靈可以分辨講預言的先知
是否是從神那裡差來的。&

約沙法王在撒瑪利亞與亞哈要去征戰前曾
要求問先知，亞哈招來400個先知，約沙
法發現來的都不是耶和華的先知。所以二
次要求亞哈找耶和華的先知求問（代下
18：4-6）。

伯特利的老先知通過查驗分辨猶大神人說
話的權柄，得知他的確是耶和華的先知（
王上13）。&

先知託耶和華的名
說話，所說的若不
成就，也無效驗，
這就是耶和華所未
曾吩咐的，是那先
知擅自說的（申
18：22-24）。\\    \midrule

分辨教會中先知講道者的靈

聖經告訴我們，一切的靈
，不可都信，總要試驗那
些靈是出於神的不是，因
為世上有許多假先知已經
出來了。（約一 4:1)&

保羅曾責備加拉太人沒有分辨能力，很快
的偏離了起初所領受的福音(加1:6-9)

主耶穌曾經借著使徒約翰寫信給以弗所教
會，誇獎他們曾試驗那自稱為使徒卻不是
使徒的，看出他們是假的來(啓2:2）。 

主耶穌曾經借著使徒約翰寫信給推雅推喇
教會責備他們容讓那自稱是先知的婦人耶
洗別在教會教導神的僕人(啓2:20）。 &

1，先知的靈原是順服先知的（林前14:32）

2，凡靈認耶穌基督
是成了肉身來的，
就是出於神的。凡
靈不認耶穌，就不
是出於神，這是那
敵基督者的靈（約
一 4:1-3)\\    \midrule

分辨人裏面的各种邪靈&
耶穌在一個啞巴身上分辨出附在這人裏面
使他啞巴的鬼。（路11:14)

保羅在腓利比傳福音時分辨出跟在著他的
使女靈裏有鬼（徒16:16-18)&
\\    \bottomrule
    \end{tabular}%
     \caption{辨別諸靈}
  \label{tab:addlabel}%
\end{table}%


\subsubsubsection{異能的恩賜}
異能的恩賜包括醫病（這裏所説的醫病不是用人類的醫療知識和設備，乃是神用超自然的方式醫治人類用現有的科技手段和知識無法醫治的疾病，或者使患者立刻恢復健康的醫治方式），行異能，和辨別諸靈的恩賜。這次的必讀經文中沒有提及
行異能的恩賜，所以這裏不做介紹。

\begin{table}[H]
\centering
\begin{tabular}{p{6.145em}|p{30.785em}|p{12.5em}}\toprule
 \multicolumn{1}{p{6.215em}|}{\textbf{恩賜的名稱}} & \multicolumn{1}{p{5.215em}|}{\textbf{原則簡單説明}} & \textbf{運用及原則} \\\midrule
醫病gifts ofhealing 9& 醫病的恩賜即可以醫治身體又可以醫治心靈。英文
NIV用的是gifts of healing，恩赐用的是复数。&\\\midrule
行異能&&\\\midrule
辨別諸靈&辨别諸靈的恩赐中的靈是复数& 先知的靈原是順服先知的（林前14:32）\\\midrule
\end{tabular}%
\end{table}%

\subsubsubsection{工作的恩賜}
工作的恩賜包括信心，執事（serve），施捨，治理和憐憫的恩賜。下面列表逐一簡單説明。

\begin{table}[H]
\centering
\begin{tabular}{p{6.145em}|p{30.785em}|p{12.5em}}\toprule
 \multicolumn{1}{p{6.215em}|}{\textbf{恩賜的名稱}} & \multicolumn{1}{p{5.215em}|}{\textbf{彰顯的情形}} & \textbf{運用及原則} \\\midrule

信心& 可以用於先知的講論& 照著信心的程度說預言（羅12：6）\\\midrule
&用於正確的審查評價自己&信徒當照著神分給個人信心的大小看的合乎中道（羅12：3）\\\midrule

執事Serve 10 &執事的NIV英文是serve，意思是服事&應當專一\\\midrule
施捨&& 當誠實\\\midrule
治理& Lead 11 治理的NIV英文是領導的意思& 當殷勤\\\midrule
憐憫&& 當甘心\\    \bottomrule
\end{tabular}%
\end{table}%

\subsubsection{我們對待恩賜的態度}
\begin{itemize}
\item A.要切慕屬靈的恩賜，並要求多得造就教會的恩賜（林前14:1,12）
\item B. 我們在服事中不要輕忽所得的恩賜（提前4：14）
\item C. 常提醒自己應當將所得的恩賜如火挑旺起來（提后1：6）
\item D. 所有的恩賜在使用時必須有愛，應當在愛中切慕恩賜（林前14:1）
\item E. 圣灵的恩赐不可以用钱买卖（徒8:20-23)
\end{itemize}

\subsubsection{IV. 恩賜的擁有人 （專心研讀林前 12：12-13，弗 4：1-10, 羅 12:3-8）}
恩賜雖然不同，但只要是信徒，就是恩賜的擁有者。保羅用身體和肢體比喻信徒和基督的關係（羅12:4-6），每一個信徒都是一個肢體，雖然以前的身份不同，如不同的種族或社會的地位，身份，但只要受洗歸入基督，就是從一位聖靈的感動（林前12:13），聖靈就把恩賜放在不同的肢體中。而且個人蒙恩都是照基督所量給各人的恩賜（弗4:7）說的也是一樣的道理，即每個蒙恩的人，都有神給的恩賜在身上。所以，每一位信主的信徒，都是恩賜的擁有者，不能再說我們沒有恩賜了。

\subsubsection{V. 各恩賜的重要 （林前 12:14-18, 弗 4:11-14, 羅 12:3-8, ）}


\subsubsubsection{A. 每個恩賜都不可缺}
人身上有許多的肢體，如果缺了任何一個，叫殘疾。而且每個肢體之間都是互為需要，
缺了眼睛，即使有腿走路也容易摔跤，所以保羅把每個信徒都比作一個肢體，身體需要
每一個肢體才完全，才是健康的，而每個肢體身上都有聖靈賜給的恩賜，因此，可以說
每個恩賜對教會來說也都是不可或缺的。
\subsubsubsection{B. 每個恩賜盡其職能：}
雖然恩賜各有不同，如使徒，先知，教師，傳福音的等等，但每個恩賜都不可缺少，不
僅不可缺，而且每個恩賜都有不同的職能，如傳福音的就是要按大使命去傳福音，使徒
早期的工作是到處建立教會，或像保羅把聖靈給他的啟示寫出來給各地的教會，教師把
聖經的真理教導清清楚楚，是信徒信仰的根基穩固等等，每個恩賜的職能都要被發揮出
來，即『各盡其職』（弗4：12），才能建立基督的身體，即使有異教之風，也不致被
搖動。
\subsubsubsection{C. 每個恩賜互相聯絡：}
雖然恩賜各有不同，但不是每個恩賜互不相關，也不是有哪一個恩賜大到不需要其他的
恩賜，就如身上的肢體一樣，各自有不同的用途，恩賜也是一樣，雖然每個信徒都要在
他的恩賜上要專一（作好），但要互相聯絡在一起，互補才能建造健康的身體，不能輕
看他人的恩賜或高看自己的恩賜，因為所有的事都是神賜給的，沒有一樣是我們自己靠
努力或能力得來的，所以要看得合乎中道才好（羅12：3b）。


\subsubsection{VI. 面對自己的恩賜 （林前 12:19-24, 弗 4:15-16, 羅 12:3-8）}


\subsubsubsection{A. 我不一定與別人有一樣的恩賜}
10 NIV
11 NIV

6

保羅告訴我們每個人都有恩賜，就好比身體只有一個，但肢體是多的而（林前12：19-
21）且用途不一樣（羅12：4），所以我的恩賜與別人的恩賜可能不一樣，教會的信徒
也不可能都有一樣的恩賜。
\subsubsubsection{B. 我有的恩賜與別人的一樣重要}
每個肢體都有各自的功用，不能說一個肢體不需要另個肢體（林前12：19-21），所以
神給我的恩賜，也給別的信徒恩賜，我不能因自己的恩賜就覺得自己高於他人，其他信
徒也是如此，都要看到合乎中道(羅 12:3)，因為恩賜都是神賜給，不是自己努力得來的
。
\subsubsubsection{C. 我認為不重要的恩賜神不輕看}
保羅說人以為軟弱的肢體，更不可少，人看為不體面的，神卻加上體面，讓它更俊美（
林前12：23-24）同樣，恩賜沒有哪一個是重要，哪一個是不重要的，是神給的，對教
會來說都一樣重要，我們不要輕看自己的恩賜，也不要輕看其他信徒的恩賜。
\subsubsubsection{D. 我的恩賜與別人的要互相聯絡}
每個是肢體都要連於元首基督，而且聯絡的合適，各自發揮自己的功用，才能叫身體漸
漸增長（弗4：15-16a），所以我們各自的恩賜，不是單獨的發揮作用，而是與其他肢體
的恩賜一起，各按各職，而且要彼此互助的，彼此之間的媒介就是愛，這樣才能真正發
揮作用，不但建造教會，也會在愛中建立自己（弗4：16b）。


\subsubsection{VII. 恩賜的配搭（林前 12:19-27, 弗 4:1-6, 羅 12:3-8）}

一個肢體不能成為身體，身體也不可能只有一個或都是一樣的肢體，因為肢體的不同，而且互相的配搭才能叫身體活動自如，才能健康。同樣聖靈所賜的恩賜也是不同的，多樣的，因此也需要彼此的配搭才可以建立基督的身體。恩賜的配搭原則包括：
\subsubsubsection{A. 彼此需要（林前12：21）：}
身體上有許多的肢體，每個肢體都是運用合適，身體才能正常的活動，手沒有眼睛，就
不能準確的拿到需要的東西，眼睛看見了，沒有手也是拿不到東西，可見肢體互相的需
要才可以完成一個動作，一項任務。這種互相的需要在恩賜上也是一樣。在神的教會裡
，一個恩賜不能使教會被建造，比如有傳福音的恩賜的人，把福音傳給人，接下來就需
要有人去關心新的信徒，也需要有人去教導他們聖經的真理才能在真理上成長，如果軟
弱了還需要有人去勸勉，有的信徒生活有需要了，還要有人願意奉獻所需等，這就是恩
賜上的彼此需要。
\subsubsubsection{B. 彼此看重（林前12：22-24，弗4：1-2a）：}
身體上的肢體，有的很容易受傷，比如眼睛，看起來好像很軟弱，但如果沒有眼睛行動
就困難。所以，身體上的肢體缺一不可。同樣，恩賜也是如此，聖靈賜給教會信徒的恩
賜，每一個都有不同的功用，有人恩賜多一些，有人少一些，但這是神分配的，所以信
徒之間不可以在恩賜上彼此輕看，過度看重自己或每個人的恩賜，忽略其他人的恩賜；
反而要彼此看重，個人要謙卑、溫柔、忍耐，互相寬容，這樣才與蒙召的恩相稱。
\subsubsubsection{C. 彼此相顧（林前12：25）：}
身體上的肢體互相聯絡，才可以有身體的和諧的動作。當一隻手拿東西要掉落時，眼睛
看見，大腦反應，另一隻手就會立刻幫助接住，腿也要支撐，腰也要用力，而不是眼睛
看見卻幸災樂禍不管，另一隻手也躲開。恩賜的服事也不能分開，有的恩賜服事會有明
顯的果效，如有傳福音恩賜的信徒，他一講，有人就信主，其他人的反應是與他一起歡
喜，還是心生嫉妒？還是一起去關心新的信徒，幫助他成長呢? 或者有人恩賜不明顯，
服事不得力，很灰心，就如受苦的肢體一樣，其他信徒就應當去陪伴，看如何去幫助他
。這就是恩賜服事上的彼此相顧，而不是彼此嫉妒。

7

\subsubsubsection{D. 彼此合一（弗4：2）}
不論信徒有什麼樣的恩賜，首先要讓自己的恩賜發揮出來，並且要專一（羅12：6-
8），照著神分給各自的恩賜在教會裡作好，比如教導的就專一教導，就是把教導的工
作作好。憐憫人的，要甘心的作等。雖然每個信徒要專一作自己的服事，但心裡當清楚
，神是一位，主是一位，聖靈也是一位，所以各自發揮各自的恩賜，都是要服事同一位
主，恩賜來與同一位聖靈，因此，恩賜的服事合一更重要，無論什麼樣的恩賜，什麼樣
的服事，什麼樣的職分，都要保守聖靈所賜合而為一的心。

\subsubsection{VIII. 恩賜的配合（林前 12:27-31, 弗 4:1-10, 羅 12:3-8）}
恩賜是聖靈所賜，分配在不同的肢體身上，恩賜有多種（弗4：8b），分給各人的也不同，
在教會裡不都是使徒，也不會都是教師（林前12：29-30），所以，恩賜的服事需要彼此配
搭，互相配合才能發揮所有肢體不同的用處（羅12：4）。而且在神的教會裡，神的心意是
各司其職，即每個肢體都要參與服事，把所有恩賜都充分發揮使用出來，才能建立基督的身體。比如说方言和翻方言的恩賜是彼此配合最明顯的實例，有人有說方言的恩賜，如果在聚會時，他說方言或用方言禱告，沒有會翻譯方言的恩賜的信徒在場，別人聽不懂，就不知道他在說什麼，禱告感謝後，別人也不會說阿門（羅14：16）。反之，如果沒有人會說方言，翻方言的恩賜也不可能體現出來，因此，教會裡恩賜配合好才能真正體現其意義和用途。

\subsubsection{IX. 在愛中使用恩賜 （最妙之道） （林前 13章）}

\subsubsubsection{A. 愛比恩賜更重要的原因 （13:1-3）}
保羅講了各樣的恩賜，但到最後，他告訴教會信徒要渴慕更大的恩賜，就是愛。如果沒
有愛，恩賜會成為什麼呢？
1. 沒有愛的恩賜是鳴的鑼響的鈸（林前13：1）
當我們信主後，主的愛就住在我們裡面。聖靈的恩賜主要體現在外在的表現，比如
言語方面，能力方面，或具體的工作方面，如果只有注重外在的表現，只為了完成
一項工作或一項服事，忽略了內在的因素，比如謙卑，溫柔，忍耐，和平，寬容等
，不顧被一起服事或被服事者的內心感受，不會造就人，反而會傷害人。就如鳴的
鑼響的鈸，外面聲音大而已。
2. 沒有愛的恩賜是與我們無益的（林前13：2-3）
如果一個人有信心可以挪山，或者把所有一切周濟窮人，不是出於愛的緣故，不是
為了愛神，不是為了愛服事的對象，只是顯出了自己的恩賜，就沒有任何的益處。
\subsubsubsection{B. 愛的重要性 (13:4-8)}
1. 對待別人的（林前13：4）：
我們對愛的定義，是看別人對我們有沒有愛的表現，但保羅這裡重新給我們定義了
愛-那就是我們對別人有沒有付出愛。首先說愛是恆久忍耐和恩慈，在恩賜的服事中
，如果對同工有忍耐，就不會只是看重效率。如果給人傳福音有忍耐，就會注重對
福音對象的關心，而不是他能不能立刻信主。如果傳福音的對象反對，甚至用言语
攻擊我們時，更需要忍耐，甚至還要對他跟有恩慈，如果我们也用言语回击对方，
我们就没有好的见证。愛是不嫉妒，嫉妒是看到他人比自己強，或看到他人有的比
自己多，或看到他人比自己好而產色的內心的活動，在恩賜服事上不可有嫉妒之心
，反而要與喜樂的人同樂，當其他信徒在恩賜服事上有果效，甚至比自己作的好，
要與他一同快樂。
2. 不要作的事（林前13：5-6a）：
愛是不嫉妒，不自誇，如果自己在恩賜服事上有果效，卻要有『是不自誇，不張狂

8

』才好。恩賜的服事也不要期望為自己地回報，不要求自己的益處，為他人著想，
別人作錯了，也不要記在心裡，不計算人的惡。
3. 要作的事（林前13：6-7）：
凡事包容，包容是因為別人作的不夠好，我們可以接納，或者別人作了傷害我們的
事，我們也可以忍耐，這對一起服事的同工尤其重要。凡事相信，不是相信人，而
是相信神的心意是好的；凡事盼望也是在恩賜服事上短時看不到果效時，還對神有
信心，相信神有一天會成就；凡事忍耐，要作到凡事包容，凡事相信，凡事盼望都
需要有忍耐的心。
4. 愛永存（林前13：8）：當主再來時，恩賜會終止，結束，但愛卻會存到永遠。
\subsubsubsection{C. 恩賜與愛的影響力（林前 13:8-13）}
恩賜是為要建造教會，成全聖徒而用，當主耶穌再來時，我們與主面對面，現在不
明白的，模糊的，無論是知識上的，或應該天國的事，就都清楚了，全都知道了，
所以，先知講道的恩賜不再需要，因為我們不需要有人來教導了，到那時所知道的
，就如同主知道我們一樣。因此，恩賜的服事是暫時的，如果只注重恩賜，沒有愛
，我們所作的一切就不會存到永遠。

\subsubsection{X. 我們對待恩賜的態度}
\subsubsubsection{A. 愛是恩賜服事的根基}
（所有的恩賜在使用時必須有愛，應當在愛中切慕恩賜（林前14:1）
恩賜服事中，要把愛放在首位，愛是恩賜服事的出發點，愛要貫穿於恩賜服事的整個過
程。
\subsubsubsection{B. 不要高看自己}
恩賜是聖靈按己意賜給信徒的（林前12：8，弗4：7），不是信徒自己努力得來的，所
以，無論有多少恩賜，用恩賜服事有多大的果效，都不要驕傲，是聖靈借我們來作的工
作，要把榮耀歸給神。
\subsubsubsection{C. 不要輕看別人：}
恩賜是用來建造教會，成全聖徒（弗4：12），這就要求恩賜的服事是信徒之間彼此的
互動，而不是一個人單獨的工作，沒有其他同工（恩賜）的配合，不能完成教會群體的
建造，就如保羅說的，肢體之間不能說我不需要其他的肢體。
\subsubsubsection{D. 要切慕屬靈的恩賜，並要求多得造就教會的恩賜（林前14:1,12）}
恩賜有多種，保羅教導信徒在追去愛的同時，也要切慕屬靈的恩賜，更有羨慕的是作先
知講道的恩賜，因為先知講道的恩賜是造就教會的，而有的恩賜是造就自己的，如說方
言的（林前14：1-4）。
\subsubsubsection{E. 願意與他人搭配（林前12：27，29-30）}
恩賜的服事好比身體上的肢體，需要彼此的配搭，才能起到建造教會的作用。所以在信
徒清楚自己的恩賜同時，也要在服事中主動尋找有其他恩賜的同工一起配搭，教會是個
群體，不是個體，群體的服事中要學會與他人的配合。
\subsubsubsection{F. 各別恩賜需要特別的態度}
恩賜服事需要與其他同工一起配搭，但不代表可以忽略自己恩賜服事的重要，個體服事
的重要性表現在不同的恩賜，有的需要特別的態度來使用這些恩賜。
1. 信心的態度（羅12：6）
說預言的的恩賜，在使用時，要照著信心的程度說預言，不能誇大或減小。
2. 專一的態度（羅12：7-8a）
需要有專一態度的恩賜包括執事，教導和勸化，如果一位同工有講道的恩賜，就讓他
專一教導，或者作執事，就專一作執事，作勸化的，就專一作勸化，在這些恩賜服事
時盡量不要被其他事物打擾。
3. 誠實的態度（羅12：8）
需要誠實的服事態度的恩賜是施捨的恩賜，在施捨時要誠實，或叫大方。
4. 殷勤的態度（羅12：8）
有治理恩賜的，在治理時要殷勤，不可懶惰。
5. 甘心的態度（羅12：8）
有憐憫的恩賜的人，在服事時要甘心樂意，不要表現出被迫在作。
\subsubsubsection{G. 不要排斥恩賜（林前14：1）}
保羅教導說聖靈顯在各人身上是叫人的益處，信徒要切慕屬靈的恩賜，服事中會有失敗
，挫折，軟弱，但不應以此就拒絕恩賜，排斥恩賜，因為恩賜的服事會讓基督的身體漸
漸增長，並且可以在愛中建立自己（弗4：16）。
\subsubsubsection{H. 不輕忽恩賜（林前12：23）}
保羅用身體上人以為軟如的肢體更是不可少的，他在告訴我們恩賜中我們以為小的，果
效不明顯的，或者在我們眼中看不上的，反而是更重要的，所以，我們不要輕忽任何一
個恩賜。
\subsubsubsection{I. 不強求別人與我一樣（林前12：4-6）}
聖靈所賜的恩賜是不同的，當我們作一項事工時需要同工，有時我們就期望每個人都有
跟我們一樣的恩賜就好了，甚至強迫別人作不是他恩賜所擅長的服事。
\subsubsubsection{J. 同苦同樂的心，別人得榮耀了我們不要嫉妒（林前12；26）}
保羅講到用愛來服事時，說道愛是不嫉妒，這在恩賜服事上也很重要，當其他同工服事
有效，有成就時，我們的心態是什麼？是嫉妒，為什麼神給他那麼多恩賜不給我？還是
與他一同高興快樂，看到神家的事工有果效，我們一起快樂，一起得榮耀？所以在恩賜
服事中要有同苦同樂的態度才對。
\subsubsubsection{K. 不要禁止恩赐（林前14：39）}
所以恩賜都是聖靈所賜，都是為成全聖徒，建造教會，所以不要有禁止某個恩賜的心。
保羅說要切慕先知講道，但也不要禁止說方言。
\subsubsubsection{L. 順服的態度（林前14：29-32）}
在恩賜服事中也要有彼此順服的態度，如果先知講道的，旁邊的人得了啟示，先說話的
就要閉口，不但要順服聖靈的引導，也要彼此順服。
















\subsubsection{說方言}
\begin{table}[H]
\centering
\begin{tabular}{p{6em}|p{40em}}  \toprule
经文&内容\\  \midrule

徒2:4&他們就都被聖靈充滿，按著聖靈所賜的口才說起別國的話來。\\  \midrule
%徒2:6，8，11&当日从天下各国来的人有十四个外邦城镇，十五种口音。单单犹太人就有三种，即本土的犹太人，从罗马来的犹太人与各国来的犹太人。他们的口音并不相同（参拙作使徒行传讲义第二章五至十三节详解）。\\  \midrule

徒10:44-46&彼得還說這話的時候，聖靈降在一切聽道的人身上。 45 那些奉割禮、和彼得同來的信徒，見聖靈的恩賜也澆在外邦人身上，就都稀奇， 46 因聽見他們說方言，稱讚神為大。\\  \midrule
徒19:6-7&保羅按手在他們頭上，聖靈便降在他們身上，他們就說方言，又說預言。一共約有十二個人\\  \midrule
林前14:21-22&律法上記著：「主說：『我要用外邦人的舌頭和外邦人的嘴唇向這百姓說話，雖然如此，他們還是不聽從我。』」這樣看來，說方言不是為信的人做證據，乃是為不信的人；做先知講道不是為不信的人做證據，乃是為信的人。\\  \midrule
徒8:18-24&西門看見使徒按手便有聖靈賜下，就拿錢給使徒，說：「把這權柄也給我，叫我手按著誰，誰就可以受聖靈。」彼得說：「你的銀子和你一同滅亡吧！因你想神的恩賜是可以用錢買的。你在這道上無份無關，因為在神面前你的心不正。你當懊悔你這罪惡，祈求主，或者你心裡的意念可得赦免。我看出你正在苦膽之中，被罪惡捆綁。」西門說：「願你們為我求主，叫你們所說的，沒有一樣臨到我身上！」\\  \midrule

林前12:8-11&这人从圣灵领受智言，那人从圣灵领受知识；圣灵赐这人信心，赐那人医病的恩赐；圣灵使这人能行神迹，使那人能做先知讲道；使这人能辨别诸灵，使那人能说方言，又使另一人能翻译方言。这一切都是同一位圣灵所赐的，是祂按自己的旨意分别赐给各人的。\\  \midrule
\end{tabular}
\end{table}





\subsection{没有分类}
腓立比書 2:1
所以，在基督裡若有什麼勸勉，愛心有什麼安慰，聖靈有什麼交通，心中有什麼慈悲憐憫，
%約 4:14人若喝我所賜的水，就永遠不渴。我所賜的水要在他裡頭成為泉源，直湧到永生。」

%以弗所書 1:13,14你們既聽見真理的道，就是那叫你們得救的福音，也信了基督，既然信他，就受了所應許的聖靈為印記。 14這聖靈是我們得基業的憑據，直等到神之民被贖，使他的榮耀得著稱讚。

%以弗所書 4:30不要叫神的聖靈擔憂，你們原是受了他的印記，等候得贖的日子來到。
%哥林多後書 1:22他又用印印了我們，並賜聖靈在我們心裡做憑據。

%路11:13你們雖然不好，尚且知道拿好東西給兒女，何況天父，豈不更將聖靈給求他的人嗎？」

%路24:49我要將我父所應許的降在你們身上，你們要在城裡等候，直到你們領受從上頭來的能力。」


加拉太書 3:14
這便叫亞伯拉罕的福，因基督耶穌可以臨到外邦人，使我們因信得著所應許的聖靈。

 51 你們這硬著頸項、心與耳未受割禮的人，常時抗拒聖靈！你們的祖宗怎樣，你們也怎樣。 





%圣灵应验了主同在的应许
%出埃及记 33：14 耶和华说：“我必亲自和你同去，使你得安息。”
%请参见 
%赛 63：11; 
%赛 63：14; 
%该 2：5

圣灵实施所应许的新约
耶利米书 31：33–34 耶和华说：“那些日子以后，我与以色列家所立的约乃是这样：我要将我的律法放在他们里面，写在他们心上。我要作他们的神，他们要作我的子民。 他们各人不再教导自己的邻舍和自己的弟兄说：‘你该认识耶和华。’因为他们从最小的到至大的，都必认识我。我要赦免他们的罪孽，不再记念他们的罪恶。”这是耶和华说的。
请参见 
罗 8：2; 
罗 8：9–10; 
林后 3：3; 
林后 3：6; 
来 8：10–11





所应许的圣灵将赋予弥赛亚的生命和工作以能力
以赛亚书 61：1–2 主耶和华的灵在我身上，因为耶和华用膏膏我，叫我传好信息给谦卑的人，差遣我医好伤心的人，报告被掳的得释放，被囚的出监牢； 报告耶和华的恩年，和我们神报仇的日子，安慰一切悲哀的人；
请参见 
赛 11：2; 
路 4：18–19




徒 11：16

耶稣基督应许了圣灵

约 15：26; 
约 16：7–15; 
徒 1：5–8
在耶稣基督复活和升天之后的时期应许赐予圣灵
约翰福音 7：39 耶稣这话是指着信他之人要受圣灵说的。那时还没有赐下圣灵来，因为耶稣尚未得着荣耀。
请参见 

徒 2：33

所应许的圣灵应验了神对亚伯拉罕的话语
加拉太书 3：14 这便叫亚伯拉罕的福，因基督耶稣可以临到外邦人，使我们因信得着所应许的圣灵。

所应许的圣灵向信徒们保证将来的产业
以弗所书 1：13–14 你们既听见真理的道，就是那叫你们得救的福音，也信了基督，既然信他，就受了所应许的圣灵为印记。 这圣灵是我们得基业的凭据，直等到神之民被赎，使他的荣耀得着称赞。



\section{講道}
\subsection{彼得的講道}
\subsection{司提反的講道}
\subsection{保羅的講道}


\section{保羅的同工團隊}
\begin{table}[htbp]
  \centering
    \begin{tabular}{|l|l|l|l|}
    \toprule
   名字 & 同工時間  & 種族，職業 &   經文\\
    \midrule
       阿尼色弗&阿尼色弗一家的人屢次使保羅暢快 & &提後1:16\\ 
&在羅馬的時候，殷勤的找保羅，並且找著了 & &提後1:17\\ 
&在以弗所多多的服事保羅 & &提後1:18\\
&問阿尼色弗一家人安 & &提後4:19\\ \midrule      
    阿尼西母& 保羅寫歌羅西書時受差送信&歌羅西教會的人 &歌羅西書4:9  \\
  & 保羅寫腓利門書為此人求情&  &腓利門書1:11  \\   \midrule
    巴拿巴  &把賣田地的價銀全部奉獻給使徒 & 生在居比路利未人 & 4:36-37\\
    &在耶路撒冷第一個接納保羅 & 名叫約瑟 & 9:26-30\\
    &在安提阿服事，使多人信主 &好人被聖靈充滿  & 11:22-24\\   
     &把保羅從大數找回安提阿  &  & 11:25-26\\
      &與保羅把捐項捐項送往耶路撒冷 &   & 11:27-30\\    
      &與保羅第一次宣教 &   & 13，14\\
        &與耶路撒冷參加會議 &   & 15:1-35 \\
          &第二次宣教時，因馬可與保羅分開 &   & 15:36-41 \\   \midrule 
  畢羅斯的兒子所巴特 &第三次宣教的馬其頓開始加入&庇哩亞人  & 20:1-4  \\
    \midrule 
     布田&寫提後時，與保羅同工  &  & 提後4:21 \\ \midrule
     大馬哩&第二次宣教的雅典 &雅典  &17:34 \\   \midrule
     德丟 &為保羅代筆寫羅馬書 & &羅16:22\\ \midrule
     底馬& 保羅寫歌羅西書時與保羅在一起&  &歌羅西書4:11  \\     
   & 保羅寫提後時因貪愛世界去了帖撒羅尼迦&  &提後4:10  \\
  &保羅寫腓利門書時，與保羅同工 & &腓利門書1:24\\ 
 \midrule
     丟尼修&第二次宣教的雅典 &雅典,亞略巴古的官  &17:34 \\    \midrule
     循都基&在福音上曾與保羅一同勞苦 &腓立比教會同工 &腓立比書4:2\\ \midrule
  非比 &第三次宣教,給羅馬教會帶信&堅革哩教會中的女執事&羅16:1 \\
     \midrule
     腓利門&  & 腓立比教會同工&腓利門書1:1\\ \midrule
     福徒拿都&保羅寫林前時，與保羅同在 & &林前16:16-17\\ \midrule 
     該猶   &第三次宣教的以弗所 & 馬其頓,特庇人 &19:29 \\
      &與保羅一起去亞西亞 &   &20:4  \\
      & 接待全教會，與保羅一起問羅馬教會安 &   & 羅16:23 \\
      &  保羅為他施洗  &  & 林前1:14 \\  
        &  與保羅一起做監  &  & 歌羅西4:10 \\  
      \midrule    
 革老底亞&寫提後時，與保羅同工  &  & 提後4:21 \\ \midrule 
 革利免&和保羅一同做工 & 腓立比教會同工&腓立比書4:3\\ 
     \bottomrule
    \end{tabular}%
      \caption{保羅的同工團隊1}
  \label{tab:addlabel}%
\end{table}%       
\begin{table}[htbp]
  \centering
    \begin{tabular}{|l|l|l|l|}
    \toprule
   名字 & 同工時間  & 種族，職業 &   經文\\     
    \midrule
       基利司布 &第二次宣教的哥林多 &哥林多,管會堂 &18:8\\  
     &  保羅為他施洗  &  & 林前1:14 \\   
    \midrule   
    括土 &保羅寫羅馬書時，與保羅同工 & &羅16:24\\ \midrule
 利奴&寫提後時，與保羅同工  &  & 提後4:21 \\ \midrule    
     路加 &第二次宣教（特羅亞）-最後 &特羅亞人 &16:10 \\   
    &寫提後時,與保羅在一起 &  &提後4:11 \\
    &寫歌羅西書時，與保羅在一起問歌羅西教會安 &  &歌羅西4:14 \\ 
      &保羅寫腓利門書時，與保羅同工 & &腓利門書1:24\\  \midrule    
      路求 & 在安提阿教會與保羅一起禁食禱告&安提阿教會同工 &13:1  \\
 & 第三次宣教寫羅馬書時與保羅同在& 保羅親屬,古利奈人 & 羅16:21 \\ \midrule
     呂底亞&第二次宣教的腓立比 &推雅推喇人 &16:14-15 \\     \midrule
    馬可（約翰）   &保羅送捐項後，他從耶路撒冷來安提阿 &猶太人 &12:25 \\
    &第一次宣教時做幫手 & &13:5 \\
       &第一次宣教到旁非利亞的別加離開保羅 & &13:13 \\
       &第二次宣教時與巴拿巴坐船去居比路 & &15:36-39 \\      
     &寫提後時吩咐提摩太把馬可帶來 & &提後4:11 \\
       &寫歌羅西時與保羅一起，可能受差去歌羅西 & &歌羅西4:10 \\
       &保羅寫腓利門書時，與保羅同工 & &腓利門書1:24\\
       \midrule 
  司提反&服事聖徒& 亞該亞初結的果子&林前16:15\\
&保羅寫林前時，與保羅同在 & &林前16:16-17\\ \midrule  
     所提尼 &第二次宣教的哥林多，寫林前 &哥林多,管會堂 &18:17  \\ 
          &與保羅一起寫林前 &  & 林前1:1 \\  
  所西巴德 &保羅寫羅馬書，問安時提及此人 & &羅16:21 \\ 
     \bottomrule
    \end{tabular}%
      \caption{保羅的同工團隊2}
  \label{tab:addlabel}%
\end{table}%

\begin{table}[htbp]
  \centering
    \begin{tabular}{|l|l|l|l|}
    \toprule
   名字 & 同工時間  & 種族，職業 &   經文\\     
    \midrule
 
   特羅非摩  &第三次宣教的馬其頓開始加入&亞西亞以弗所人&20:1-4 \\
    &保羅被捕前與保羅一起在耶路撒冷& &21:29 \\
   &寫提後時病了，被保羅留在米利都 & &提後4:20 \\ \midrule   
       提多猶士都 &第二次宣教的哥林多 & &18:7 \\   \midrule   
    提多  &耶路撒冷會議 &希利尼人，未受割禮 &加2:1,4 \\
    &寫提後時,提多去撻馬太 &  &提後4:10 \\
     &寫提多書時,提多在革哩底 &  &提多1:5 \\
       &寫提多書時,打發提多去尼哥波立見保羅 &  &提多3:12 \\
         &寫林後時,提多從哥林多回來 &  &林後7:6，8章 \\  \midrule    
    提摩太   & 在路司得加入保羅的第二次宣教 &路司得人 & 16:1-3 \\
        &在庇哩亞傳道時，保羅離開後與西拉留守&  &17:10-15 \\         
        &在亞西亞時和以拉都一起去被打發去馬其頓&  &19:22 \\          
    & 寫羅馬書時，與保羅同工  & 父:希利尼，母:猶太 &  羅16:21  \\  
    &寫林前時 ，打發提摩太去哥林多教會  && 林前4:17，16:10 \\   
   & 和保羅一起寫林後  && 林後1:1 \\ 
    & 二次宣教時與保羅去哥林多  &  &  林後1:19  \\ 
    & 和保羅一起寫腓利門書&& 腓利門1:1 \\ 
      & 和保羅一起寫歌羅西書  &&歌羅西1:1 \\  
        & 和保羅一起寫腓立比書  && 腓立比1:1 \\ 
           & 寫腓立比書後打發提摩太去腓立比教會送信  && 腓立比2:19-23 \\  
            &被保羅打發去帖撒羅尼迦教會 && 帖前3:2 \\     
             & 從帖回來後與保羅一起寫帖前 && 帖前1:1 \\                                            
               & 與保羅一起寫帖後 && 帖後1:1 \\        
    & 寫提摩太前書時被保羅安置在以弗所教會  && 提前1:3 \\   
       & 寫提摩太後書時提到其母親祖母  &外祖母羅以，母親友尼基& 提後1:5 \\        
       &被保羅打發去帖撒羅尼迦教會 && 帖前3:2 \\        
      &曾經入獄後被釋放 && 希伯來書13:23 \\        
    \midrule    
   推基古   &第三次宣教的馬其頓開始加入  &亞西亞人& 20:1-4\\
    & 寫提多書時，提到可能打發推基古去提多那裡  & &提多3:12\\
    & 寫以弗所書時，保羅打發推基古去以弗所教會送信  & &弗6:21\\ 
  & 寫歌羅西書時，保羅打發推基古去歌羅西教會送信  &  忠心的執事&歌4:7，8\\ 
   & 寫提後時，保羅打發推基古去以弗所  &  &提後4:12\\   
 
    \bottomrule
    \end{tabular}%
      \caption{保羅的同工團隊3}
  \label{tab:addlabel}%
\end{table}%


\begin{table}[htbp]
  \centering
    \begin{tabular}{|l|l|l|l|}
    \toprule
   名字 & 同工時間  & 種族，職業 &   經文\\ 
          \midrule
 亞里達古&第三次宣教的以弗所開始加入& 帖撒羅尼迦人  & 19:29 \\
      &與保羅一起去亞西亞 &   &20:4  \\
      & 與保羅一起在被捆鎖中去羅馬&    &  27:2  \\
     &保羅寫腓利門書時，與保羅同工 & &腓利門書1:24\\ 
    \midrule 
 亞力山大&第三次宣教的以弗所&猶太人  & 19:33-34  \\
     &提後書中保羅讓提摩太提防此人，他曾苦害保羅&銅匠  & 提後4:14  \\
    \midrule

亞提馬&寫提多書時，提到可能打發亞提馬去提多那裡  & &提多3:12\\ \midrule
亞波羅&保羅在林前時勸他去哥林多，遭拒&& 林前16:12\\ \midrule 
亞該古&保羅寫林前時，與保羅同在 & &林前16:16-17\\ \midrule
亞居拉妻百基拉&保羅二次宣教在哥林多遇見他們 & 猶太人生在本都&18:1-4 \\
    &與保羅一起去以弗所 &製造帳篷為業 &18:18-19 \\ 
     &在以弗所接待亞波羅給他講福音 & &18:24-26\\   
    &保羅寫羅馬書問他們羅馬家中的教會安 &  &羅16:3-5 \\  
     &寫提後時，保羅問他們安  &  & 提後4:19\\ 
    &保羅寫林前時，與保羅同在 & &林前16:19\\ \midrule
亞腓亞& 保羅在腓利門書裡問她安 & 腓利門妹妹      &腓利門書1:2\\ \midrule
亞基布& 保羅在腓利門書裡提及此人 & 與保羅同當兵      &腓利門書1:2\\ \midrule
亞迦布& 從耶路撒冷到安提阿預言天下有打飢荒 & 先知     &11:27-28\\  \midrule  
亞拿尼亞 &在大馬色奉主命令為保羅施洗  & 大馬色門徒      &9:10-18\\  \midrule  
 
耶數(猶士都)& 保羅寫歌羅西書時與保羅一同做工&猶太人 &歌羅西書4:11  \\
 \midrule
耶孫&第二次宣教的帖撒羅尼迦 &帖撒羅尼迦人?  &17:1-9  \\  
    &寫羅馬書時，與保羅同工 &  & 羅16:21 \\   \midrule  

以巴弗提 & 被腓立比教會所差服事保羅&腓立比人 &腓立比2:25-30  \\ \midrule      
以巴弗&將歌羅西教會情況匯報給保羅 &基督忠心的執事 &歌羅西書1:7，8\\
&保羅寫歌羅西書時與保羅一起禱告 & &歌羅西書4:12\\
&為歌羅西,老底嘉並希拉波立的弟兄多多的勞苦 & &歌羅西書4:13\\
&保羅寫腓利門書時，以巴弗與保羅一同坐監 & &腓利門書1:23\\
 \midrule
    以拉都&第三次宣教的亞西亞-以弗所 &幫助保羅的人 &19:21 \\
    &第三次宣教寫羅馬書 & 城內管銀庫&羅16:24 \\
   &寫提後時，以拉都在哥林多住下 &  &提後4:20 \\      
    \midrule  
    
 友布羅 &寫提後時，與保羅同工  &  & 提後4:21 \\ \midrule
友阿爹&在福音上曾與保羅一同勞苦 &腓立比教會同工 &腓立比書4:2\\  

    \bottomrule
    \end{tabular}%
      \caption{保羅的同工團隊4}
  \label{tab:addlabel}%
\end{table}%

\begin{table}[htbp]
  \centering
    \begin{tabular}{|l|l|l|l|}
    \toprule
   名字 & 同工時間  & 種族，職業 &   經文\\ 
      \midrule
      西拉 &耶路撒冷會議後與保羅一起去安提阿 &猶太人，先知 &15:22-35 \\
     &在腓立比獄中與保羅一起唱詩禱告&  &16:25-40 \\
      &在帖撒羅尼迦傳道被迫與保羅提前離開去庇哩亞&  &17:1-9 \\  
       &在庇哩亞傳道時，保羅離開後與提摩太留守&  &17:10-15 \\       
       &二次宣教時被留在馬其頓，後到哥林多會保羅 &  &18:5 \\
     &第二次宣教 &  &15:36-41，林後1:19 \\
      & 與保羅一起寫帖前 && 帖前1:1 \\  
          & 與保羅一起寫帖後 && 帖後1:1 \\                
     \midrule   
     西公都    & 第三次宣教的馬其頓開始加入&帖撒羅尼迦人&20:1-4 \\
     \bottomrule
    \end{tabular}%
      \caption{保羅的同工團隊5}
  \label{tab:addlabel}%
\end{table}%

